\chapter{Mathematical modelling: real problem to math}
\label{vol02:ch18}

\epigraph{All models are wrong, but some are useful.}{Aphorism in
working circulation among statisticians and engineers, popularised by
George Box (1976). The chapter takes the second half of the sentence as
seriously as the first.}

\chapmeta{Half-life: foundational. Archetypes: ethics (responsible
modelling), uncertainty, scaling, balance, optimisation, feedback. Exercise target: 25. Project
track: analysis only. Hazard class: none. Reader path default: standard,
with sections explicitly tagged. Named cases: none in chapter prose; the
recognition-level discussion in section 18.8 names financial and
operational episodes without invoking the registry. Companion assets:
\texttt{volumes/02-form/ch18-mathematical-modelling/figures/} (eighteen
TikZ figures, including the Buckingham-Pi data collapse, the SIR
compartment diagram, the Kleiber log-log scaling line, the
thermal-mass model-revision pair, the M/M/c queue state chain with its
congestion nonlinearity, the queue prediction-versus-simulation and
staffing-cost optimisation pair, the model-selection penalised-fit
curve, the contact-network contagion pair, the feedback-loop block
diagram with its closed-loop step response, the analytic-versus-Monte-Carlo
propagation histogram, and the bias-variance decomposition); \texttt{code/} (nine Python scripts:
cooling-cup fit, logistic
fit, identifiability demonstration, bootstrap, cross-validation, overfit
demonstration, Buckingham-Pi tool, sensitivity propagation, model
triangulation); \texttt{data/} (cooling-cup notebook, synthetic logistic
data, pendulum periods, beam-deflection perturbation table).}

\noindent
We close Volume II by codifying a habit the reader has been practising
for seventeen chapters. Modelling is the act of choosing what to keep,
what to throw away, and what to test against the world. The fifteen
hundred pages of this volume have built the apparatus by which a
quantity becomes a function, a function becomes an equation, an equation
acquires parameters, parameters acquire intervals, and intervals acquire
predictions that an outside engineer can check. This chapter puts the
moves in one place, names the act of modelling as the act it is, and
asks the reader to take one quantity from their Volume I notebook and
turn it into a model that survives a hostile review.

The chapter occupies the same structural slot in Volume II that
Chapter 9 occupies in Volume I. Volume I closed with estimation as a
moral discipline; Volume II closes with modelling as the same
discipline at greater apparatus. The reader who has carried the
estimation habit through eighteen chapters of mathematics has been
prepared, by repeated practice, for the move this chapter formalises.
Every estimate the reader has made was a small model. Every worked
example in chapters 5 through 17 was a model held to a standard. The
chapter's job is to name the standard, expose the choices the standard
hides, and turn the volume's apparatus into a working method the
reader can apply to a question outside this book.

% Figure: the modelling cycle. Six labelled stages around a loop, with a
% feedback arrow from validation back to abstraction (the iterative
% revision the chapter installs as discipline).
\begin{figure}[htbp]
\centering
\begin{tikzpicture}[
  stage/.style={draw=engblack, fill=engblue!12, rounded corners=3pt,
    minimum width=2.6cm, minimum height=0.85cm, font=\small\bfseries,
    align=center},
  arrow/.style={-{Stealth}, thick, draw=engblack!70},
  feedback/.style={-{Stealth}, thick, dashed, draw=engred!80},
  node distance=0.4cm and 0.6cm,
]
% Six stages around a hexagon
\node[stage] (q)  at ( 0.0,  3.0) {Question};
\node[stage] (a)  at ( 4.0,  1.8) {Abstraction};
\node[stage] (i)  at ( 4.0, -1.8) {Identification};
\node[stage] (v)  at ( 0.0, -3.0) {Validation};
\node[stage] (u)  at (-4.0, -1.8) {Uncertainty};
\node[stage] (l)  at (-4.0,  1.8) {Look-back};

% Forward cycle
\draw[arrow] (q) -- (a);
\draw[arrow] (a) -- (i);
\draw[arrow] (i) -- (v);
\draw[arrow] (v) -- (u);
\draw[arrow] (u) -- (l);
\draw[arrow] (l) -- (q);

% Feedback: validation forces revised abstraction
\draw[feedback] (v) to[bend left=25] node[midway, fill=white, inner sep=2pt,
  font=\scriptsize\itshape, align=center] {revise on\\ residual\\ structure} (a);

% Centre note
\node[font=\small\itshape, align=center, color=engblack!70] at (0, 0)
  {the modelling cycle\\ (one iteration)};
\end{tikzpicture}
\caption[The modelling cycle]{The modelling cycle as the chapter
installs it. The solid arrows mark one full pass: question, abstraction,
identification, validation, uncertainty report, look-back. The dashed
arrow marks the iterative revision triggered when validation reveals
residual structure that the current abstraction does not capture; the
revised model re-enters the cycle at the abstraction stage. A model
that has not closed the loop at least once has not yet been validated;
a model that has closed the loop without revision has either been
correct on the first abstraction (rare) or has been built without
inspecting the residuals (common, and the failure mode of
section~18.8).}
\label{fig:vol02:ch18:modelling-cycle}
\end{figure}


The cycle in figure~\ref{fig:vol02:ch18:modelling-cycle} is the
visual table of contents for the chapter. The first three sections
(\S 18.1--18.3) install the forward arc: what a model is, how
abstraction selects what to keep, how identification and validation
turn data into a checked model. Sections~18.4 and~18.5 install the
return arc: parameter and structural uncertainty, the look-back that
revises the abstraction when residuals carry information the model
does not yet capture. Sections~18.6 and~18.7 walk the full loop on
three worked systems; section~18.8 names the modes of failure that
follow from short-circuiting any arc. The project closes the loop on
a quantity of the reader's choice.

\section{What modelling is}\pathtag{core}

A \engterm{model} is a finite description of a system, kept honest by
the act of being checked against the system. The description may be a
single equation, a system of equations, an algorithm, a graph, a table,
or any combination of these; what makes it a model rather than a
notation is the discipline of comparing its predictions against
measurement.

The definition is narrower than the everyday use of the word. A drawing
of a building is a model of the building's geometry, but only when the
draughtsman has measured the building and recorded the discrepancies
does the drawing acquire the engineering status of a model. A
spreadsheet that estimates next year's revenue is a model of the firm's
finances, but only when last year's spreadsheet has been compared
against last year's actual revenue does the spreadsheet acquire the
engineering status. The check separates a model from a story.

\subsection{The model as a hypothesis}

A model is a \engterm{hypothesis} about how a quantity in the world is
produced. The hypothesis takes the form: if the world behaves according
to this description, then a measurement of the world should produce a
result within a stated interval. The interval is part of the
hypothesis. A model that predicts a single number with no interval is
not yet a model in this sense; it has not yet exposed itself to
falsification.

Consider the simplest model the reader has met. Volume I Chapter 4
asked the reader to measure the period of a pendulum. Volume II
Chapter 7 derived, from the small-amplitude linear pendulum equation,
the prediction
\[
T = 2\pi \sqrt{\frac{L}{g}}.
\]
That equation is a model. It carries an underlying hypothesis: the
restoring force is proportional to the displacement, the damping is
negligible, the support is rigid, and the amplitude is small. Each
hypothesis can be falsified. The reader who measures the period of a
swing at a children's playground, with $L = 2.5\,\si{\meter}$, expects
$T \approx 3.17\,\si{\second}$. A measured period of $3.20\,\si{\second}$
falls within the model's expected interval (set by amplitude
nonlinearity, support compliance, and the reader's stopwatch). A
measured period of $4.50\,\si{\second}$ does not. The discrepancy is the
information; the model has done its job.

\subsection{The model as a working object}

A useful model is also a \engterm{working object}. The engineer can
manipulate it, ask it questions about regimes that have not yet been
tested, and produce predictions that guide intervention. The pendulum
equation, manipulated, produces the design rule that doubling the
length increases the period by a factor of $\sqrt{2}$. The reader who
needs a clock pendulum with a period of $2\,\si{\second}$ inverts the
equation and obtains $L = g (T / 2\pi)^{2} \approx 0.994\,\si{\meter}$.
The model has supplied a design value the engineer can build to.

The working-object character of a model is what separates engineering
modelling from descriptive science. A descriptive scientist may be
satisfied with a model that explains the data. An engineer needs a
model that, in addition, predicts the behaviour of a system not yet
built or not yet observed, with stated uncertainty, in time for a
decision. The engineering standard is more demanding because the
consequence of a wrong prediction is a built artefact that does not
work or a process that does not run.

\subsection{The model and the system}

A model and the system it represents are different things, and the
discipline that keeps them separate is one of engineering modelling's
quiet defences. The pendulum equation describes a swing; the swing is
the physical object; the equation lives on the page. The distinction
is well-worn, and we restate it because the well-wornness is what
protects the engineer from a real failure mode. A model that is
treated as the system loses the discipline of being checked against
the system. A reader who reasons within the pendulum equation about
swings that violate small-amplitude or rigid-support assumptions has
left the model's domain of validity and has not yet noticed.

The remedy is mechanical. Every model carries a domain of validity:
the range of inputs and conditions under which it has been checked,
or under which the underlying assumptions are defensible. A model
applied outside its domain is no longer a model; it is a guess. The
engineer who states the domain explicitly, before applying the model,
keeps the model honest.

\begin{archetype}[Ethics, first introduction]
The engineer's responsibility for the predictions the model makes,
the decisions the model supports, and the world the model leaves
unmodelled. Volume X gives the archetype its formal home; Volume II
introduces it here because a model that is built without a sense of
responsibility for its consequences is the most consequential
failure mode this volume describes. Responsibility takes three
working forms. First: stating the domain of validity, so that the
model is not applied to questions it cannot answer. Second:
reporting the uncertainty alongside the prediction, so that the
decision the model supports is informed by what the model does not
know. Third: documenting what the model leaves out, so that a
subsequent engineer can recover the modelling choices and revise
them when new data or new questions arrive. The archetype recurs
in every later volume that asks the engineer to act on a model's
prediction.
\end{archetype}

\section{Choosing what to keep and what to throw away}\pathtag{core}

The first hard move in modelling is the move of \engterm{abstraction}:
deciding which features of the system are essential to the question
being asked and which features are not. The decision is the modeller's
craft, and it is the move that distinguishes a working model from a
caricature or from an unmanageable simulacrum.

\subsection{The pendulum, three ways}

The same physical system admits many models. Consider a pendulum
swung from a string in still air at small amplitude.

The simplest model treats the bob as a point mass, the string as a
massless rigid rod, the pivot as frictionless, and the air as absent.
The equation of motion is $\ddot\theta + (g/L) \sin\theta = 0$, which
under the small-angle approximation $\sin\theta \approx \theta$
linearises to $\ddot\theta + (g/L)\theta = 0$, with period
$T = 2\pi\sqrt{L/g}$. Three lines, four assumptions, one prediction.

A second model relaxes the small-angle assumption. The exact period of
the simple pendulum involves an elliptic integral,
\[
T(\theta_{0}) = 4\sqrt{\frac{L}{g}}\, K(\sin(\theta_{0}/2)),
\]
where $K$ is the complete elliptic integral of the first kind and
$\theta_{0}$ is the maximum angular amplitude. The model is more
accurate at large amplitude. Its closed-form simplicity is gone; the
period now depends on amplitude, and the engineer who needs a clock
pendulum has acquired a new design constraint (keep the amplitude
small, so the period does not drift with the spring drive's decay).

A third model adds linear damping and a forcing term, giving the
driven damped oscillator
\[
\ddot\theta + 2\zeta\omega_{0}\dot\theta + \omega_{0}^{2}\theta
  = F(t)/(mL),
\]
with $\omega_{0}^{2} = g/L$. The equation predicts amplitude decay,
resonance, and the steady-state phase lag. It is the model the reader
needs when designing a pendulum-driven seismometer or analysing a
clock's escapement, because those applications turn on damping and
forcing rather than on free oscillation.

Three models, one system, three different questions. The first model
is the right model for a children's swing or a first-pass clock
estimate. The second is the right model for a high-amplitude
mechanical metronome. The third is the right model for a damped
oscillator under a periodic drive. None of the three is universally
better; each is the correct abstraction for its question.

% Figure: three models of the same pendulum, arranged by abstraction
% level. The axes are not numerical; they are conceptual (parameter
% count vs domain of validity).
\begin{figure}[htbp]
\centering
\begin{tikzpicture}[
  modelbox/.style={draw=engblack, fill=englime!18, rounded corners=2pt,
    minimum width=4.0cm, minimum height=1.1cm, font=\small, align=center},
  axislab/.style={font=\small\itshape, color=engblack!70},
  arrow/.style={-{Stealth}, thick, draw=engblack!70},
]
% Axes (conceptual)
\draw[arrow] (0, 0) -- (10.0, 0) node[right, axislab] {parameter count};
\draw[arrow] (0, 0) -- (0, 5.0) node[above, axislab, align=center] {domain\\ of validity};

% Three models
\node[modelbox] (m1) at (1.8, 1.0) {Linear small-angle\\
  $\ddot\theta + (g/L)\theta = 0$\\ (2 parameters)};
\node[modelbox] (m2) at (5.0, 2.6) {Exact small-amplitude\\
  $T = 4\sqrt{L/g}\,K(\sin(\theta_0/2))$\\ (3 parameters)};
\node[modelbox] (m3) at (8.2, 4.1) {Damped driven\\
  $\ddot\theta + 2\zeta\omega_0\dot\theta + \omega_0^2\theta = F/(mL)$\\
  (5 parameters)};

% Linking arrows (one model added a feature)
\draw[arrow, dashed, draw=engred!70] (m1) -- node[midway, fill=white,
  inner sep=2pt, font=\scriptsize\itshape, align=center] {drop small-\\angle
  assumption} (m2);
\draw[arrow, dashed, draw=engred!70] (m2) -- node[midway, fill=white,
  inner sep=2pt, font=\scriptsize\itshape, align=center] {add damping\\
  and forcing} (m3);

% Annotation
\node[align=left, font=\footnotesize, color=engblack!70] at (5.0, -0.9)
  {Each step adds parameters and extends the model's domain.
   The right model is the simplest model that answers the question.};
\end{tikzpicture}
\caption[Three pendulum models, arranged by abstraction]{Three models
of the same physical pendulum, arranged by parameter count and domain
of validity. The linear small-angle model has two parameters and a
narrow domain (small amplitude, undriven, undamped); the exact
small-amplitude model relaxes the linearisation at the cost of a third
parameter (amplitude $\theta_0$); the damped driven model extends to
forced and damped regimes with two further parameters. No model in
the figure dominates the others; each is the correct abstraction for
a question whose precision and regime it serves. The discipline of
abstraction is to choose the smallest model on the chain that the
question requires.}
\label{fig:vol02:ch18:abstraction-hierarchy}
\end{figure}


Figure~\ref{fig:vol02:ch18:abstraction-hierarchy} arranges the three
models on conceptual axes (parameter count, domain of validity) so
the trade-off is visible. Each step rightward adds a parameter and
extends the regime in which the model is defensible. The discipline
of abstraction is the discipline of choosing the leftmost model
whose domain covers the question, because every parameter beyond
that point carries a cost the reader pays in identification effort,
in identifiability ambiguity, and in the chance of misuse.

\subsection{The discipline of abstraction}

Abstraction has rules. The reader who abstracts well follows them;
the reader who abstracts badly produces models that are either too
simple to predict or too complex to manipulate.

\engterm{Rule one: name the question first}. A model is built for a
question. A model built without a question is a notational exercise.
The question fixes the output the model must predict, the units the
output is in, and the precision the prediction must meet. Without
that fix, the modeller has no way to decide which features are
essential.

\engterm{Rule two: name the dominant scales}. Volume I taught the
reader to estimate before calculating; the same habit applies to
modelling. Before any equation is written, the reader estimates the
order of magnitude of every quantity in the candidate model. The
estimate identifies which terms in a candidate equation are large
and which are small. Small terms can often be dropped; large terms
must be kept. The estimate is performed in dimensionless ratios,
which is what the next rule formalises.

\engterm{Rule three: nondimensionalise}. The candidate model is
rewritten in dimensionless form, with all variables and parameters
expressed as ratios to characteristic scales. The dimensionless
groups that emerge identify the regimes the model occupies. Reynolds
number, P{\'e}clet number, Mach number, and the like are not
artefacts of fluid mechanics; they are the structural information
that nondimensionalisation extracts from any candidate model. The
groups whose magnitude is large or small relative to unity identify
the terms the model should retain or drop.

\engterm{Rule four: check the limits}. The candidate model is
exercised in its limits: send a parameter to zero, send another to
infinity, set two parameters equal, impose a symmetry. The model
must produce defensible behaviour in each limit, or the
abstraction has discarded a feature the limit needed. Volume I
Chapter 9 introduced the limiting-case check as a sanity check on
calculations; here the check applies to the model itself.

\engterm{Rule five: choose the simplest model that answers the
question}. Among the candidate models that pass the previous four
rules, the simplest is preferred. The principle is sometimes stated
as Ockham's razor; in engineering practice it has a sharper form.
A simpler model has fewer parameters to identify, fewer
opportunities for unmeasured covariates to contaminate the fit,
fewer degrees of freedom to overfit the data, and fewer
opportunities for a subsequent engineer to misuse the model
outside its domain. Complexity is a cost the engineer pays only
where the question requires it.

\subsection{The Meadows discipline: dynamic systems}

Donella Meadows's \emph{Thinking in Systems} \cite{gen:meadows2008}
codifies the abstraction discipline for dynamic systems. The book
is written in plain prose, without engineering notation, and yet
the rules it lays down are the rules an engineer follows. We name
three of Meadows's working principles because they translate
directly into modelling decisions the engineer has to make.

\engterm{Stocks and flows are the structural primitives}. A stock
is an accumulation; a flow is a rate of accumulation or
depletion. The mathematical apparatus is the integral and the
derivative the reader met in Volume II Chapters 5 and 6: the
stock is the integral of the net flow. A model that confuses
stocks with flows, that treats a rate as if it were an
accumulation or vice versa, has the wrong structural primitives
and will produce wrong predictions even when its parameters are
correct.

\engterm{Feedback loops set the dynamics}. Whether a stock grows,
shrinks, oscillates, or settles depends on the feedback loops in
which it sits. A balancing loop returns the stock toward a
setpoint; a reinforcing loop drives the stock further from
equilibrium. The model that omits a loop will miss the dynamic
the loop produces. Volume II Chapter 7 introduced the
mathematical apparatus (linear ODE stability, eigenvalue
analysis); Meadows's contribution is the discipline of finding
the loops in the system before writing the equations.

\engterm{Delays change everything}. A balancing loop with no delay
is stable; a balancing loop with a long delay relative to its
correction time is oscillatory or unstable. The engineer who
omits a delay because the delay is short on the eye's timescale
has discarded a feature the model needs. The modelling decision
is to identify, before writing equations, which delays are short
relative to the question's timescale and can be dropped, and
which are not.

\subsection{What gets thrown away}

A useful model throws something away. The list of what gets
thrown away is part of the model's documentation, not an
embarrassment to be hidden.

For the small-angle pendulum: the rod's mass, the air drag, the
support compliance, the higher harmonics, the tide-driven
variation of $g$, the relativistic correction. Each is small in
the regime the model addresses; each can become large in a
regime where the model would be wrong; each is named, so a
subsequent engineer can recover the assumption.

For the RC circuit of Volume II Chapter 7: the resistor's
inductance, the capacitor's leakage current, the wire's resistance,
the voltage source's internal impedance, the temperature
dependence of the components, the tolerance band of the
components. Each is small for the operating regime; each can
become large outside it; each is named in the model's
documentation.

For a population dynamics model of Chapter 13: the spatial
distribution, the age structure, the genetic variation, the
weather, the pathogens not in the model, the human action that
the model treats as exogenous. Each can be small or large
depending on the question; each is named.

The principle is structural. The discipline of naming what gets
thrown away is what allows the model to be revised. A model whose
omissions are documented can be improved by a subsequent
engineer; a model whose omissions are silent cannot be.

\section{Dimensional analysis and the Buckingham Pi theorem}\pathtag{core}

Rule three of the abstraction discipline said to nondimensionalise. The
rule has a formal engine, and the engine deserves a section of its own
because it is the cheapest source of structural information a modeller
ever gets. Before a single parameter is identified, before a single
datum is taken, dimensional analysis tells the modeller how many
independent dimensionless groups govern the problem, what those groups
are, and what functional form the answer must take. The cost is a few
lines of bookkeeping; the return is a model whose structure is fixed by
the dimensions of its variables alone.

\subsection{Dimensions constrain equations}

Every physical equation must be dimensionally homogeneous: the
dimensions of the left side equal the dimensions of the right side.
The principle is older than the calculus and was put on a formal
footing by Fourier and later by Buckingham
\cite{paper:buckingham1914} and Bridgman \cite{text:bridgman1922}. Its
force is that a candidate equation that is not dimensionally homogeneous
is wrong before it is tested, and a candidate equation that is
homogeneous is constrained to a narrow family of forms.

Consider the period of a pendulum, approached as if the small-angle
solution were unknown. The period $T$ (dimension $\mathsf{T}$) can
depend only on the length $L$ ($\mathsf{L}$), the mass $m$
($\mathsf{M}$), and the gravitational acceleration $g$
($\mathsf{L}\mathsf{T}^{-2}$). The only dimensionless combination of
$L$, $m$, $g$, $T$ is $T\sqrt{g/L}$, because $m$ cannot appear (no other
variable carries mass, so mass cannot be cancelled). Dimensional
homogeneity therefore forces
\[
T \sqrt{g/L} = \text{const},
\qquad\text{hence}\qquad
T = \text{const} \cdot \sqrt{L/g}.
\]
The constant is dimensionless and cannot be found by dimensional
analysis alone; the small-angle solution supplies $2\pi$. The structure
of the answer, the $\sqrt{L/g}$ dependence and the irrelevance of mass,
came from dimensions before any dynamics were written. A modeller who
proposed a period that depended on mass would be caught by this check
immediately.

\subsection{The Buckingham Pi theorem}

The general statement organises the bookkeeping. A relation among $n$
physical variables that involves $k$ independent primary dimensions
(mass, length, time, temperature, charge, and so on) can be rewritten
as a relation among $n - k$ independent dimensionless groups, the
\engterm{Pi groups}. The theorem is named for Buckingham's 1914 paper
\cite{paper:buckingham1914}, though the result was in circulation
earlier; Barenblatt's monograph \cite{text:barenblatt1996} is the
modern reference that connects the theorem to scaling and
self-similarity.

The recipe is mechanical. List the $n$ variables and their dimensions.
Count the number $k$ of independent primary dimensions that appear.
Choose $k$ of the variables as a dimensionally independent basis (the
\engterm{repeating variables}). Form $n - k$ dimensionless products,
each combining one of the remaining variables with powers of the
repeating variables chosen to cancel all dimensions. The relation among
the original variables is equivalent to a relation among the $n - k$ Pi
groups, and that relation has fewer arguments, which is the reduction
the theorem buys.

\subsection{Worked case: drag on a sphere}

The drag force $F$ on a sphere moving through a fluid is the textbook
case, and it is worth working in full because the reduction it produces
is the foundation of the whole experimental tradition of fluid
mechanics. The drag can depend on the sphere diameter $d$
($\mathsf{L}$), the relative speed $U$ ($\mathsf{L}\mathsf{T}^{-1}$),
the fluid density $\rho$ ($\mathsf{M}\mathsf{L}^{-3}$), and the fluid
dynamic viscosity $\mu$ ($\mathsf{M}\mathsf{L}^{-1}\mathsf{T}^{-1}$).
With the force $F$ ($\mathsf{M}\mathsf{L}\mathsf{T}^{-2}$), there are
$n = 5$ variables and $k = 3$ primary dimensions
($\mathsf{M}, \mathsf{L}, \mathsf{T}$), so the theorem predicts
$n - k = 2$ dimensionless groups.

Choose $\rho$, $U$, $d$ as the repeating variables (they are
dimensionally independent: $\rho$ carries the only mass, $U$ the only
time, and the three together span length). Form the first group from
the force:
\[
\Pi_1 = \frac{F}{\rho U^2 d^2},
\]
which the reader can verify is dimensionless by substituting the
dimensions, and which is, up to the conventional factor of
$\tfrac{1}{2}$ and the use of frontal area $A = \pi d^2/4$, the
\engterm{drag coefficient} $C_D = F / (\tfrac{1}{2}\rho U^2 A)$. Form
the second group from the viscosity:
\[
\Pi_2 = \frac{\rho U d}{\mu} = \mathrm{Re},
\]
the \engterm{Reynolds number}. The theorem now asserts that the entire
five-variable relation collapses to
\[
C_D = \phi(\mathrm{Re})
\]
for some single-argument function $\phi$ that dimensional analysis
cannot supply but experiment can measure once and for all.

The reduction is the modelling payoff.
Figure~\ref{fig:vol02:ch18:buckingham-collapse} shows what it means in
the laboratory. Drag measured on spheres of three different diameters,
in fluids of different density and viscosity, forms three separate
families when plotted as force against speed. Plotted as $C_D$ against
$\mathrm{Re}$, the three families collapse onto one curve. An
experimenter who did not know the theorem would measure drag for every
combination of diameter, speed, density, and viscosity that the
application might encounter, a four-dimensional grid of experiments. The
theorem reduces the grid to a single curve in one variable. The same
$\phi(\mathrm{Re})$ governs a sand grain settling in water and a weather
balloon rising through air, because both occupy the same range of the
one group that matters.

% Figure: data collapse under nondimensionalisation. Sphere drag data,
% scattered when plotted dimensionally, collapsing onto one curve when
% plotted as drag coefficient versus Reynolds number.
\begin{figure}[htbp]
\centering
\begin{tikzpicture}[font=\small]
% ===== Left panel: dimensional scatter =====
\begin{scope}[xshift=0cm]
\draw[->, thick, draw=engblack!70] (0,0) -- (5.2, 0)
  node[right, font=\scriptsize] {speed $U$ (m/s)};
\draw[->, thick, draw=engblack!70] (0,0) -- (0, 4.3)
  node[above, font=\scriptsize] {drag force $F$ (N)};
\node[font=\small, anchor=south west] at (0.2, 4.1) {\textbf{(a) dimensional}};
% Three families of points (three different sphere diameters / fluids),
% each a different quadratic-ish curve: F = c U^2 with different c.
\foreach \c/\col in {0.18/engred!85, 0.34/engblue!85, 0.09/engblack!60} {
  \foreach \U in {0.6,1.2,1.8,2.4,3.0,3.6} {
    \pgfmathsetmacro\Fy{\c*\U*\U}
    \filldraw[\col] (\U*1.3, {\Fy}) circle (0.05);
  }
}
\node[anchor=west, font=\scriptsize] at (1.4, 3.7) {\textcolor{engred!85}{$\bullet$} $d_1$};
\node[anchor=west, font=\scriptsize] at (2.4, 3.7) {\textcolor{engblue!85}{$\bullet$} $d_2$};
\node[anchor=west, font=\scriptsize] at (3.4, 3.7) {\textcolor{engblack!60}{$\bullet$} $d_3$};
\end{scope}
% ===== Right panel: nondimensional collapse =====
\begin{scope}[xshift=7.2cm]
\draw[->, thick, draw=engblack!70] (0,0) -- (5.2, 0)
  node[right, font=\scriptsize] {$\log_{10}\mathrm{Re}$};
\draw[->, thick, draw=engblack!70] (0,0) -- (0, 4.3)
  node[above, font=\scriptsize] {$\log_{10} C_D$};
\node[font=\small, anchor=south west] at (0.2, 4.1) {\textbf{(b) collapsed}};
% Ticks Re
\foreach \xt/\xv in {0.5/0, 1.6/1, 2.7/2, 3.8/3, 4.9/4} {
  \draw[draw=engblack!70] (\xt,0) -- (\xt,-0.07);
  \node[font=\scriptsize, anchor=north] at (\xt,-0.09) {\xv};
}
% Standard C_D(Re) curve (Stokes 24/Re at low Re, plateau ~0.4 in
% intermediate, drag crisis dip near Re~3e5). Plot schematically.
\draw[very thick, draw=engblue!80, smooth, samples=80, domain=0.3:4.6]
  plot (\x*1.1+0.1, { 2.5 - 1.0*\x + 0.16*\x*\x });
% All three families' points now on one curve
\foreach \dx in {0.6,1.2,1.8,2.4,3.0,3.6,4.2} {
  \pgfmathsetmacro\yy{2.5 - 1.0*\dx + 0.16*\dx*\dx}
  \filldraw[engred!85] (\dx*1.1+0.1, \yy) circle (0.045);
}
\node[anchor=west, font=\scriptsize, color=engblack!70] at (1.4, 3.5)
  {one curve: $C_D = C_D(\mathrm{Re})$};
\end{scope}
\end{tikzpicture}
\caption[Data collapse under the Buckingham Pi reduction]{Sphere-drag
measurements at three diameters and fluid combinations. (a) Plotted in
dimensional variables, the force-versus-speed data form three separate
families, one per sphere. (b) Plotted as the dimensionless drag
coefficient $C_D = F / (\tfrac{1}{2}\rho U^2 A)$ against the Reynolds
number $\mathrm{Re} = \rho U d / \mu$, the three families collapse onto a
single curve. The Buckingham Pi reduction (four dimensional variables,
three primary dimensions, hence one dimensionless group beyond $C_D$)
predicted the collapse before any datum was taken: the dimensionless drag
coefficient can depend only on the dimensionless Reynolds number. The
collapse both compresses the data and exposes the regimes (Stokes
$1/\mathrm{Re}$ at low $\mathrm{Re}$, the intermediate plateau, the
drag-crisis dip near $\mathrm{Re}\approx 3\times 10^5$).}
\label{fig:vol02:ch18:buckingham-collapse}
\end{figure}


\begin{archetype}[Scaling, working home]
The discipline of carrying dimensionless groups through a model so that
a result measured at one scale transfers to another. The archetype has
its canonical home in Volume I Chapter 9 (orders of magnitude) and
recurs wherever a model must cross scales: the wind-tunnel test of a
scale aircraft (match Reynolds and Mach number, not size), the
pilot-plant scale-up of a chemical reactor (match Damk\"ohler and
Reynolds number, not volume), the animal-model dosing of a drug (match
the dimensionless exposure, not the milligram). A model whose
dimensionless groups are matched between the test article and the
deployed article transfers; a model matched on dimensional quantities
alone does not. The recurring engineering error is to scale a prototype
by geometry and assume the physics scales with it. The square-cube law
that the next section works is the simplest instance of why it does not.
\end{archetype}

\begin{mastery}
The choice of repeating variables is free, within the constraint that
they be dimensionally independent, and different choices produce
different but equivalent sets of Pi groups. The set
$\{C_D, \mathrm{Re}\}$ above came from choosing $\rho, U, d$; choosing
$\mu, U, d$ instead produces $\{F/(\mu U d), \mathrm{Re}\}$, and since
$F/(\mu U d) = C_D \cdot \mathrm{Re} \cdot (\pi/8)$ the two
descriptions carry the same information. The freedom is occasionally
useful: in the Stokes regime ($\mathrm{Re} \ll 1$), where drag is
linear in viscosity, the group $F/(\mu U d)$ tends to a constant
(Stokes found $3\pi$) while $C_D$ diverges as $24/\mathrm{Re}$, so the
viscosity-based group is the better coordinate at low Reynolds number.
The lesson generalises: nondimensionalisation has a gauge freedom, and
the skilled modeller chooses the gauge in which the function of interest
is simplest in the regime of interest. The companion script
\texttt{dimensional\_analysis.py} forms the dimension matrix for a list
of variables, computes its rank (the number $k$), and returns a basis
for the null space, which is a complete set of independent Pi-group
exponents; the reader can feed it any variable list and read off the
groups.
\end{mastery}

\section{Scaling and orders of magnitude}\pathtag{standard}

Dimensional analysis tells the modeller which groups matter; scaling
tells the modeller how a quantity changes when a characteristic size
changes. The two are the same discipline viewed from two angles, and
scaling is the angle the modeller uses when the question is ``what
happens if this gets bigger?''

\subsection{The square-cube law}

Geometric similarity is the assumption that an object is scaled
uniformly: every length multiplied by the same factor $\lambda$. Under
geometric similarity, areas scale as $\lambda^2$ and volumes as
$\lambda^3$. The mismatch between the two exponents is the
\engterm{square-cube law}, and it is the source of more failed
scale-ups than any other single fact.

A structure's weight scales with its volume, as $\lambda^3$. The
cross-sectional area of the columns that must carry that weight scales
as $\lambda^2$. The stress in the columns, weight divided by area,
therefore scales as $\lambda^3/\lambda^2 = \lambda$: it grows in direct
proportion to the scale factor. A design that is safe at one size is
overstressed by the factor $\lambda$ at scale $\lambda$. Galileo stated
the principle in 1638, observing that a large animal cannot have bones
geometrically similar to a small animal's, because the bones would
crush under the animal's own weight. The engineering form is identical:
a bridge cannot be scaled up by photographic enlargement, because the
stresses do not scale with the photograph.

The same exponent mismatch governs heat. Heat generated by a body
scales with its volume or mass, as $\lambda^3$; heat dissipated through
its surface scales with area, as $\lambda^2$. The ratio of generation to
dissipation scales as $\lambda$, so a large body runs hotter than a
geometrically similar small one. A small electric motor cools itself by
its surface; a motor scaled up by a factor of ten generates a thousand
times the heat and dissipates only a hundred times as much, and runs
ten times closer to its thermal limit. The modeller who scales a
prototype motor by geometry and assumes the thermal margin scales with
it has made the square-cube error in its thermal form.

The square-cube law is the simplest member of a general family, and
naming the family shows the reader when the naive exponents hold and
when nature routes around them. Any quantity $y$ that scales with a
characteristic size $\ell$ as a power, $y \propto \ell^{\,a}$, is an
\engterm{allometric} relation, and the exponent $a$ is fixed by the
geometry only when the object stays geometrically similar as it grows.
The interesting engineering cases are the ones where it does not.
Galileo's bones crush under the naive $\lambda^1$ stress rise precisely
because bone strength scales as cross-section, $\lambda^2$, while weight
scales as $\lambda^3$; nature's response is that large animals are not
geometrically similar to small ones, but grow disproportionately thick
bones, so that bone cross-section scales faster than $\lambda^2$ and the
stress stays bounded. The bone exponent departs from the geometric value
by exactly the amount needed to hold the stress constant, which is the
signature of a system that has been selected, or designed, to defeat the
square-cube law rather than to obey it. The engineering lesson is the
converse of the naive one: the square-cube exponents are the prediction a
model makes \emph{under geometric similarity}, and a measured exponent
that departs from them is not an error but a report that the system
abandoned geometric similarity to survive its own scaling. A designer
scaling a structure reads the departure as a specification: to hold a
stress or a temperature constant across scale, some dimension must grow
disproportionately, and the required disproportion is the difference
between the geometric exponent and the constant-quantity exponent. This
is the same move Kleiber's $3/4$ exponent forced in the previous
subsection, generalised: a scaling exponent that refuses its geometric
value is the structural information the next model must explain.

\subsection{Power laws and the log-log straight line}

A quantity $y$ that scales as a power of another quantity $x$,
$y = a x^{b}$, plots as a straight line on log-log axes, because
$\log y = \log a + b \log x$. The slope is the exponent $b$ and the
intercept is $\log a$. The log-log straight line is the modeller's
first diagnostic when a power law is suspected: a straight line confirms
the power law and measures its exponent in a single fit, and a curve on
log-log axes is the signal that the relation is not a pure power.

% Figure: a power law as a straight line on log-log axes. Metabolic rate
% versus body mass (Kleiber's 3/4 law) as the worked scaling example.
\begin{figure}[htbp]
\centering
\begin{tikzpicture}[font=\small]
\draw[->, thick, draw=engblack!70] (0,0) -- (8.0,0)
  node[right, font=\scriptsize] {$\log_{10}$ body mass (kg)};
\draw[->, thick, draw=engblack!70] (0,0) -- (0,5.2)
  node[above, font=\scriptsize] {$\log_{10}$ metabolic rate (W)};
% x ticks: -3 .. 3
\foreach \xt/\xv in {0.4/-3, 1.6/-2, 2.8/-1, 4.0/0, 5.2/1, 6.4/2, 7.6/3} {
  \draw[draw=engblack!70] (\xt,0) -- (\xt,-0.07);
  \node[font=\scriptsize, anchor=north] at (\xt,-0.09) {\xv};
}
% y ticks: -1 .. 4
\foreach \yt/\yv in {0.4/-1, 1.4/0, 2.4/1, 3.4/2, 4.4/3} {
  \draw[draw=engblack!70] (0,\yt) -- (-0.07,\yt);
  \node[font=\scriptsize, anchor=east] at (-0.09,\yt) {\yv};
}
% Kleiber line: slope 3/4. y = 0.75 x + const. Draw across the panel.
% Pick: at x=-3 (0.4), y ~ 0.5; at x=3 (7.6), y ~ 4.8.
\draw[very thick, draw=engblue!80] (0.4, 0.55) -- (7.6, 4.95);
\node[anchor=west, font=\scriptsize, color=engblue!80, rotate=27] at (3.2, 2.5)
  {slope $= 3/4$};
% Data points scattered along the line (mouse -> elephant)
\foreach \px/\py in {0.7/0.8, 1.5/1.35, 2.6/2.05, 3.6/2.75, 4.4/3.35, 5.4/4.0, 6.6/4.55, 7.2/4.85} {
  \pgfmathsetmacro\jit{0.10*sin(\px*220)}
  \filldraw[engred!85] (\px, {\py+\jit}) circle (0.06);
}
% Annotation labels for endpoints
\node[font=\scriptsize, anchor=south east, color=engblack!70] at (1.4, 1.4) {mouse};
\node[font=\scriptsize, anchor=north west, color=engblack!70] at (6.8, 4.5) {elephant};
\end{tikzpicture}
\caption[Kleiber's law as a straight line on log-log
axes]{Metabolic rate plotted against body mass for organisms spanning
roughly six decades of mass, on log-log axes. A power law
$P = a M^{b}$ becomes the straight line $\log P = \log a + b \log M$; the
slope $b$ reads directly off the plot. Kleiber's measurements give
$b \approx 3/4$ rather than the $2/3$ a naive surface-area argument
predicts, a discrepancy that the discipline of scaling treats as
information about the underlying transport network rather than as noise.
The log-log straight line is the diagnostic the modeller reaches for
first when a quantity is suspected of scaling as a power of another:
a straight line confirms the power law and measures its exponent in one
step.}
\label{fig:vol02:ch18:scaling-loglog}
\end{figure}


Figure~\ref{fig:vol02:ch18:scaling-loglog} shows the diagnostic on
Kleiber's law, the empirical observation that an organism's metabolic
rate scales with its body mass to a power near $3/4$ across roughly six
decades of mass. A naive surface-area argument (heat is lost through the
surface, surface scales as mass to the $2/3$) predicts an exponent of
$2/3$; the measured exponent is closer to $3/4$, and the discrepancy is
not noise to be smoothed away but information about the underlying
distribution network, which later modelling traced to the geometry of
space-filling transport systems. The discipline is the one the chapter
keeps returning to: a residual that survives a careful fit, here the gap
between $2/3$ and $3/4$, is the next model's specification.

\subsection{Worked case: the natural frequency of a cantilever}

Dimensional analysis fixes the form of a scaling law before any dynamics
are written, and a second worked case shows the move on a vibration
problem the reader will meet again in Volume III. The fundamental natural
frequency $f$ (dimension $\mathsf{T}^{-1}$) of a cantilever beam can
depend on its length $L$ ($\mathsf{L}$), its cross-sectional area moment
of inertia $I$ ($\mathsf{L}^4$), its Young's modulus $E$
($\mathsf{M}\mathsf{L}^{-1}\mathsf{T}^{-2}$), and its mass per unit length
$m$ ($\mathsf{M}\mathsf{L}^{-1}$). Five variables, three primary
dimensions, so the Buckingham count gives $5 - 3 = 2$ dimensionless
groups. Choosing $L$, $E$, $m$ as the repeating set and combining with $f$
and with $I$ gives the groups
\[
\Pi_1 = f L^2 \sqrt{\frac{m}{E I_0}}, \qquad \Pi_2 = \frac{I}{L^4},
\]
where $I_0$ is any fixed reference inertia used to make $\Pi_1$
dimensionless. Folding $\Pi_2$ back in, dimensional homogeneity forces
\[
f = \frac{\text{const}}{L^2}\sqrt{\frac{E I}{m}},
\]
and the dynamics (Euler-Bernoulli beam theory) supply only the
dimensionless constant, $1.875^2/(2\pi) \approx 0.56$ for the fundamental
mode. The structure, $f \propto L^{-2}$ and $f \propto \sqrt{EI/m}$, came
from dimensions alone. The scaling consequence is immediate and
counterintuitive: a beam scaled up geometrically by $\lambda$ in every
dimension has $I \propto \lambda^4$, $m \propto \lambda^2$, and
$L \propto \lambda$, so $f \propto \lambda^{-2}\sqrt{\lambda^4/\lambda^2} =
\lambda^{-1}$. The larger structure rings at a lower frequency in inverse
proportion to its size, which is why a large bell tolls deep and a small
bell rings high, and why a scale model's vibration test must correct the
measured frequency by $\lambda$ before it speaks to the full-size
structure. A modeller who measured a model bridge's resonance and applied
it to the prototype without the $\lambda^{-1}$ correction would predict
the wrong resonant wind speed, the failure the scaling archetype warns
against.

\begin{estimation}
\textbf{Question.} A model railway is built at $1{:}87$ scale (HO
gauge). A real locomotive masses $120\,\text{t}$ and develops a
drawbar pull of $200\,\text{kN}$. If the model were a perfect
geometric copy in the same materials, estimate its mass and the
gravitational stress on its couplers relative to the prototype's. State
why the model cannot be a dynamically faithful copy.

\smallskip
\textit{Estimate, before any calculation.}

\smallskip
Mass scales as $\lambda^3$ with $\lambda = 1/87$, so the model mass is
$120\,\text{t} \times (1/87)^3 \approx 120 \times 10^3\,\text{kg} \times
1.5\times 10^{-6} \approx 0.18\,\text{kg}$, a fifth of a kilogram, which
matches the heft of a real HO locomotive. Coupler stress is the gravity
load (scaling as $\lambda^3$) divided by the coupler cross-section
(scaling as $\lambda^2$), so the gravitational stress scales as
$\lambda$: the model's couplers see $1/87$ of the prototype's
gravitational stress. The model is therefore vastly overbuilt against
gravity, which is why it survives being dropped while the prototype
would not survive a proportional fall.

\smallskip
\textbf{Postmortem.} The model cannot be dynamically faithful because
the dimensionless groups do not match at constant gravity. Froude number
$U^2/(gL)$ governs the dynamics of a body under gravity; matching it at
$1/87$ scale requires the model speed to scale as $\sqrt{\lambda}$, so a
prototype at $90\,\text{km/h}$ corresponds to a model at
$90/\sqrt{87} \approx 9.6\,\text{km/h}$, not $90/87 \approx
1\,\text{km/h}$. A modeller who scales speed by length gets the dynamics
wrong by a factor of $\sqrt{\lambda}$. This is the same lesson the
scaling archetype names: geometry scales by $\lambda$, physics scales by
the dimensionless groups, and the two are reconciled only by deciding
which groups to hold fixed.
\end{estimation}

\section{Compartment and rate models}\pathtag{standard}

A large class of engineering models shares one structure: quantities
accumulate in compartments, and flows move them between compartments at
rates that depend on the current contents. The Meadows stocks-and-flows
discipline of section 18.2 is the qualitative version; the compartment
model is its quantitative form, and it is worth a section because it
recurs across domains the reader will meet in later volumes, from
pharmacokinetics to reactor design to epidemiology to ecology.

\subsection{The structure}

A compartment model is a system of first-order ordinary differential
equations, one per compartment, in which the rate of change of each
compartment's stock equals the sum of the inflows minus the sum of the
outflows. The apparatus is the conservation law of Volume I applied to
whatever is conserved (mass, charge, head-count, energy), and the
mathematics is the coupled linear or nonlinear ODE system of Volume II
Chapter 7. The modelling work is in identifying the compartments, the
flows, and the rate laws; once those are fixed, the equations write
themselves.

The simplest case is a single well-mixed tank of volume $V$ holding a
solute at concentration $c$, fed by an inflow at volumetric rate $Q$
and concentration $c_{\text{in}}$, and drained at the same rate $Q$.
Mass conservation gives
\[
V \dv{c}{t} = Q c_{\text{in}} - Q c,
\]
which is the same first-order linear equation as Newton's cooling, with
time constant $\tau = V/Q$ and steady state $c \to c_{\text{in}}$. The
modeller who recognises the structure inherits the entire solution
apparatus at once: the exponential approach, the time constant as the
dominant scale, the step and impulse responses. Recognising that two
physically different systems share a compartment structure is itself a
modelling move, and it is the move that lets the engineer reuse a
solution rather than rederive it.

\subsection{Worked case: the SIR epidemic model}

The susceptible-infectious-removed model is the canonical nonlinear
compartment model, and Kermack and McKendrick's 1927 derivation
\cite{paper:v6c08-kermack-mckendrick-1927} is one of the founding
results of mathematical epidemiology. Three compartments partition a
population of size $N$: the susceptible $S$, the infectious $I$, and the
removed $R$ (recovered or deceased). The flows are infection, $S \to I$,
and removal, $I \to R$. The rate laws are the modelling content. New
infections occur at a rate proportional to the number of
susceptible-infectious encounters, $\beta S I / N$; removals occur at a
rate proportional to the current infectious count, $\gamma I$. The
equations are
\[
\dv{S}{t} = -\frac{\beta S I}{N},
\qquad
\dv{I}{t} = \frac{\beta S I}{N} - \gamma I,
\qquad
\dv{R}{t} = \gamma I.
\]
The three right sides sum to zero, which is the conservation law:
$S + I + R = N$ for all time, so one equation is redundant and the
system is effectively two-dimensional.

The structural insight, available before any data, is dimensional. The
parameters $\beta$ and $\gamma$ both have dimension
$\mathsf{T}^{-1}$, so their ratio
\[
R_0 = \frac{\beta}{\gamma}
\]
is dimensionless and is the single group that governs the qualitative
behaviour. At the outset, with $S \approx N$, the infectious equation
reads $\dot I \approx (\beta - \gamma) I = \gamma(R_0 - 1) I$. When
$R_0 > 1$ the infectious stock grows exponentially and an epidemic
occurs; when $R_0 < 1$ it decays and no epidemic occurs. The threshold
$R_0 = 1$ is the model's central prediction, and it falls out of the
dimensionless group rather than from numerical integration. The
engineering quantity a hospital planner needs, the peak infectious
fraction $I_{\max}$ and its timing, is set by $R_0$ as well, and can be
read from the trajectories in
figure~\ref{fig:vol02:ch18:compartment-sir}.

% Figure: compartment diagram (stocks and flows) for the SIR model, with
% the resulting S, I, R trajectories alongside.
\begin{figure}[htbp]
\centering
\begin{tikzpicture}[font=\small,
  stock/.style={draw=engblack!75, fill=engblue!12, minimum width=1.2cm,
    minimum height=0.9cm, font=\small\bfseries},
  flow/.style={->, thick, draw=engblack!70}]
% ===== Top: compartment diagram =====
\begin{scope}[yshift=0cm]
\node[stock] (S) at (0,0) {$S$};
\node[stock] (I) at (3,0) {$I$};
\node[stock] (R) at (6,0) {$R$};
\draw[flow] (S) -- node[above, font=\scriptsize] {$\beta SI/N$} (I);
\draw[flow] (I) -- node[above, font=\scriptsize] {$\gamma I$} (R);
\node[font=\scriptsize, anchor=north] at (S.south) {susceptible};
\node[font=\scriptsize, anchor=north] at (I.south) {infectious};
\node[font=\scriptsize, anchor=north] at (R.south) {removed};
\node[font=\small, anchor=south west] at (-1.0, 0.9) {\textbf{(a) stocks and flows}};
\end{scope}
% ===== Bottom: trajectories =====
\begin{scope}[yshift=-4.2cm]
\draw[->, thick, draw=engblack!70] (0,0) -- (8.5,0)
  node[right, font=\scriptsize] {time $t$ (days)};
\draw[->, thick, draw=engblack!70] (0,0) -- (0,3.6)
  node[above, font=\scriptsize] {fraction of $N$};
\foreach \yt/\yv in {0/0, 1.5/0.5, 3.0/1.0} {
  \draw[draw=engblack!70] (0,\yt) -- (-0.07,\yt);
  \node[font=\scriptsize, anchor=east] at (-0.09,\yt) {\yv};
}
% S(t): high then falls
\draw[very thick, draw=engblue!80, smooth, samples=80, domain=0:8]
  plot (\x, { 3.0/(1 + exp(1.4*(\x-3.4))) * 0.95 + 0.05*3.0 });
% I(t): bump
\draw[very thick, draw=engred!85, smooth, samples=80, domain=0:8]
  plot (\x, { 3.0*0.42*exp(-((\x-3.6)*(\x-3.6))/1.6) });
% R(t): rises to plateau
\draw[very thick, draw=engblack!65, smooth, samples=80, domain=0:8]
  plot (\x, { 3.0*0.9/(1 + exp(-1.4*(\x-3.8))) });
\draw[very thick, draw=engblue!80] (5.6,3.2) -- (5.95,3.2);
\node[anchor=west, font=\scriptsize] at (6.0, 3.2) {$S$};
\draw[very thick, draw=engred!85] (5.6,2.8) -- (5.95,2.8);
\node[anchor=west, font=\scriptsize] at (6.0, 2.8) {$I$};
\draw[very thick, draw=engblack!65] (5.6,2.4) -- (5.95,2.4);
\node[anchor=west, font=\scriptsize] at (6.0, 2.4) {$R$};
\node[font=\small, anchor=south west] at (0.2, 3.4) {\textbf{(b) trajectories}};
\draw[dashed, draw=engblack!45] (3.6,0) -- (3.6, 1.5)
  node[above, font=\scriptsize, anchor=south] {peak $I$};
\end{scope}
\end{tikzpicture}
\caption[The SIR compartment model: stocks, flows, and
trajectories]{(a) The susceptible-infectious-removed (SIR) model drawn as
stocks and flows. Each compartment is a stock, an accumulation measured
in head-count; each arrow is a flow, a rate measured in head-count per
unit time. The infection flow $\beta S I / N$ is a reinforcing coupling
(more infectious cases drive more new infections); the removal flow
$\gamma I$ is a draining loop. (b) The trajectories that the coupled rate
equations produce: $S$ falls monotonically, $I$ rises to a peak then
falls, $R$ accumulates toward the final epidemic size. The peak in $I$ is
the engineering quantity a hospital planner needs; its height and timing
are set by the basic reproduction number $R_0 = \beta/\gamma$, the single
dimensionless group that governs the model. Kermack and McKendrick
derived the structure in 1927~\cite{paper:v6c08-kermack-mckendrick-1927}.}
\label{fig:vol02:ch18:compartment-sir}
\end{figure}


The full modelling cycle runs as before. Identification fits $\beta$ and
$\gamma$ from case-count time series, typically by nonlinear least
squares on the numerically integrated $I(t)$ or by fitting the
exponential early-growth rate $\gamma(R_0 - 1)$ and the recovery rate
$\gamma$ separately. Validation holds out the late-epidemic data and
checks whether the model fit to the early growth predicts the peak.
Structural error is large and well-catalogued: the model assumes
homogeneous mixing (every individual equally likely to meet every
other), a closed population (no births, deaths from other causes, or
migration), permanent immunity after removal, and a constant $\beta$
(no behaviour change, no seasonality). Each assumption fails in a known
regime, and each failure has a structured extension (age-structured
contact matrices, $SEIR$ with a latent compartment, waning immunity,
time-varying $\beta(t)$). The discipline of section 18.2 governs which
extension to add: find the residual structure, and add the term it
demands.

\begin{mastery}
The final epidemic size, the fraction of the population ever infected,
satisfies a transcendental equation that the dimensionless group again
controls. Dividing the $S$ equation by the $R$ equation and integrating
gives $S(t) = S_0 \exp(-R_0 [R(t) - R_0^{\text{init}}]/N)$, and taking
$t \to \infty$ with $I(\infty) = 0$ yields the final-size relation
\[
1 - \frac{S_\infty}{N} = 1 - \exp\!\left(-R_0\, \frac{N - S_\infty}{N}\right),
\]
an implicit equation for the ever-infected fraction $1 - S_\infty/N$ in
terms of $R_0$ alone. The equation has the trivial root (no epidemic)
for $R_0 < 1$ and a nontrivial root that rises steeply just above
$R_0 = 1$. The modelling lesson is that a quantity of direct
operational interest, the total caseload, depends on the data only
through the single dimensionless combination $R_0$; an enormous
reduction in what must be estimated, and a direct application of the
nondimensionalisation discipline to a nonlinear model. The
root is found by the fixed-point or Newton iteration of Volume II
Chapter 16.
\end{mastery}

\section{Identification and validation}\pathtag{core}

A model has parameters. A pendulum equation has $L$ and $g$. An RC
circuit has $R$ and $C$. A logistic growth equation has $r$ and $K$.
The act of \engterm{identification} is the act of fixing the
parameters from data. The act of \engterm{validation} is the act of
checking the identified model against data the identification did not
use.

The two acts are different. A reader who confuses them will report a
model that fits the data it was fit on and will be surprised when the
model fails on data it was not fit on. The confusion has a name in the
machine-learning literature, where it is called \engterm{overfitting};
the engineering name for the same failure mode is \engterm{validating
against the training data}, and the discipline that protects against it
is the same discipline a careful experimenter has always followed.

\subsection{Identification: from data to parameters}

Parameter identification turns measurements into numbers the model
can use. The methods are the methods of Volume II Chapter 14:
maximum likelihood, least squares, method of moments, Bayesian
inference. Each method assumes a relation between the data and the
parameters. Least squares assumes additive Gaussian noise. Maximum
likelihood assumes a chosen probability model for the noise. The
method of moments assumes the parameters can be recovered from
sample moments of the data.

For the pendulum, a reader who has measured $T$ at known $L$ can
identify $g$ by fitting
\[
T^{2} = (4\pi^{2}/g) L
\]
as a linear regression of $T^{2}$ on $L$. The slope estimate
$\hat\beta$ gives $\hat g = 4\pi^{2}/\hat\beta$. The interval on
$\hat\beta$ from Chapter 14 gives an interval on $\hat g$ via the
delta method or via direct propagation. The identification is
quantitative; the reader has not just a value but a stated
uncertainty.

For the logistic equation $\dot N = r N (1 - N/K)$, a reader who has
measured $N(t)$ at several times can identify $(r, K)$ by nonlinear
least squares on the analytical solution
\[
N(t) = \frac{K}{1 + (K/N_{0} - 1) e^{-rt}}.
\]
The fit returns point estimates and a covariance matrix, from which
intervals on $r$ and $K$ follow. If the data set is small, the
intervals are wide; if the data set spans only the early-exponential
regime, $K$ is barely identified, and the interval on $K$ swallows
its point estimate. The intervals are part of the identification's
honest report.

The identification can fail. The model can be \engterm{unidentifiable}:
the data, no matter how plentiful, do not constrain the parameters
to a unique value. The pendulum equation written as
$T = 2\pi\sqrt{(\alpha L) / (\alpha g)}$ for any $\alpha$ has the
same predictions for any rescaling of $L$ and $g$ together; the
combination $L/g$ is identifiable, but $L$ and $g$ separately are
not. The engineer who tries to identify $L$ and $g$ from period
measurements alone will produce a fit that has converged to one of
the infinitely many equivalent solutions; the engineer who does not
notice this will report a misleading interval. Identifiability is a
structural property of the model, checkable before any data are
collected, and a model that is not identifiable in the relevant
parameters is not a usable model regardless of how rich the data.

\subsection{Structural and practical identifiability}

The pendulum's $L$-and-$g$ ambiguity is an example of
\engterm{structural unidentifiability}: no quantity of data, however
rich, would separate $L$ from $g$, because the model's equations are
invariant under the transformation $(L, g) \to (\alpha L, \alpha g)$.
Cobelli and DiStefano's review of identifiability classes
\cite{paper:v2ch18-cobelli-distefano-1980} settled the language the
engineering community now uses. A parameter (or a combination of
parameters) is structurally identifiable if and only if the model's
input-output map is injective in that parameter. Structural
identifiability is a property of the model, checkable by symbolic
algebra on the equations alone.

\engterm{Practical identifiability} is a weaker, data-dependent
property. A model can be structurally identifiable in $(r, K)$ and
yet, on a given data set, return such a wide interval on $K$ that
the parameter is for working purposes unknown.
Figure~\ref{fig:vol02:ch18:identifiability} shows the regime: when
logistic data span only the exponential phase, $K$ is structurally
identifiable but practically not, because the data have not seen
the saturation that pins $K$ down. The same model, fit to data that
span the inflection, returns $K$ tightly. The practical-identifiability
diagnostic is therefore part of experimental design, not just of model
analysis; the engineer who needs $K$ designs the data acquisition to
include the regime in which $K$ is informative.

% Figure: parameter-identifiability map. Three regions in (r, K) space
% for the logistic equation, showing how the data window selects which
% parameters are tightly identified.
\begin{figure}[htbp]
\centering
\begin{tikzpicture}[font=\small]
% Axes
\draw[->, thick, draw=engblack!70] (0, 0) -- (8.5, 0) node[right, font=\small] {$r$ (intrinsic growth rate, $\si{\per\hour}$)};
\draw[->, thick, draw=engblack!70] (0, 0) -- (0, 5.5) node[above, font=\small, align=center] {$K$\\ (carrying capacity)};
% Tick labels
\foreach \xt/\xv in {1.5/0.5, 3.0/1.0, 4.5/1.5, 6.0/2.0, 7.5/2.5} {
  \draw[draw=engblack!70] (\xt, 0) -- (\xt, -0.08);
  \node[font=\scriptsize, anchor=north] at (\xt, -0.1) {\xv};
}
\foreach \yt/\yv in {1.0/1, 2.0/5, 3.0/10, 4.0/50} {
  \draw[draw=engblack!70] (0, \yt) -- (-0.08, \yt);
  \node[font=\scriptsize, anchor=east] at (-0.1, \yt) {\yv};
}

% Region A: data span only exponential regime; K barely identified
\fill[engred!15, opacity=0.7] (0.4, 0.4) rectangle (8.0, 4.8);
\draw[draw=engred!60, dashed, thick] (0.4, 0.4) rectangle (8.0, 4.8);

% Region B: data span the inflection; both identified
\fill[engblue!20, opacity=0.7] (2.0, 1.5) rectangle (5.5, 3.5);
\draw[draw=engblue!70, thick] (2.0, 1.5) rectangle (5.5, 3.5);

% Region C: data span saturation only; r barely identified
\fill[englime!25, opacity=0.7] (3.0, 2.0) ellipse (1.5 and 0.8);
\draw[draw=engblack!70, thick] (3.0, 2.0) ellipse (1.5 and 0.8);

% Labels
\node[font=\footnotesize, align=center, color=engblack] at (6.5, 4.2)
  {(A) data span\\ exponential only:\\ $r$ identified, $K$ wide};
\node[font=\footnotesize, align=center, color=engblack] at (3.75, 2.5)
  {(B) data span\\ inflection: both identified};
\node[font=\footnotesize, align=center, color=engblack, fill=white, inner sep=1pt] at (3.0, 0.9)
  {(C) data span saturation: $K$ identified, $r$ wide};
\end{tikzpicture}
\caption[Parameter identifiability and the data window]{The
identifiability of $(r, K)$ in the logistic model depends on the
window of time over which the data are collected. (A) Data confined to
the early exponential phase identify $r$ tightly (it sets the early
slope on a log scale) but leave $K$ almost unconstrained (the
saturation has not been observed). (B) Data spanning the inflection
point (population near $K/2$) identify both: the inflection's location
fixes $K$ and the slope through it fixes $r$. (C) Data confined to
the saturation tail identify $K$ tightly but leave $r$ poorly
constrained because the transient has already passed. The identifiability
diagnostic is therefore part of the experimental design: choose the
data window such that the parameters the engineer needs are in
region~(B). The discipline appears in Cobelli and Distefano's
identifiability literature \cite{paper:v2ch18-cobelli-distefano-1980}
and underwrites Exercise 18.18 (design of an identification experiment).}
\label{fig:vol02:ch18:identifiability}
\end{figure}


The working diagnostic is the condition number of the model's
Jacobian, evaluated at the fitted parameter values. A condition number
near unity says the parameters are well separated by the data; a
condition number of $10^{3}$ or more says the data leave the
parameters nearly collinear, and small noise in the data will produce
large excursions in the fit. The companion script
\texttt{logistic\_fit.py} prints the Jacobian condition number
alongside the fit; the identifiability demonstration in
\texttt{identifiability\_demo.py} replays the figure's three regimes
on synthetic data so the reader can see the parameter interval widen
as the data window narrows.

\begin{mastery}
The condition number is a fast local diagnostic; the
\engterm{profile likelihood} is the more complete one and is worth the
reader's acquaintance because it diagnoses practical non-identifiability
in a form that does not assume the quadratic approximation the
condition number relies on. To profile a parameter $\theta_j$, fix it at
a grid of values, and at each value re-optimise the remaining
parameters to maximise the likelihood. The resulting curve, the maximum
achievable likelihood as a function of $\theta_j$, is the profile. A
parameter that is well identified produces a sharp profile with a clear
minimum in the negative log-likelihood and finite confidence bounds
where the profile rises by the chi-square threshold. A structurally
non-identifiable parameter produces a flat profile: the likelihood is
unchanged as $\theta_j$ moves, because the other parameters compensate
exactly, and the pendulum's $L$ with $g$ free is the textbook flat
profile. A practically non-identifiable parameter produces a profile
that is bounded on one side and flat on the other, the signature of data
that constrain the parameter from below but not above, which is exactly
the logistic $K$ fit to exponential-phase data. The profile likelihood
thus distinguishes the three cases the condition number lumps together,
and it does so without linearising. Raue and coauthors brought the
method into systems-biology practice; the engineering reader takes from
it the working rule that a flat or one-sided profile is a refusal to
report that parameter as identified, however small the residuals.
\end{mastery}

\subsection{Validation: checking the model against held-out data}

Validation checks the identified model against data the
identification did not use. The minimum discipline is the
\engterm{train-validate split}: a fraction of the data is set aside
before identification, and the identified model is checked against
that fraction. A model that fits the training data well and the
validation data poorly has overfit; a model that fits both well has
generalised to data the identification did not see, which is the
property the engineer wanted.

For small data sets, the train-validate split discards information.
\engterm{Cross-validation} recovers the information by performing
the split many times, identifying on each training portion,
validating on each held-out portion, and aggregating the validation
scores. $k$-fold cross-validation is the standard scheme: the data
is partitioned into $k$ subsets, the model is identified on $k - 1$
of them, validated on the remaining one, and the procedure repeats
$k$ times. The aggregate validation score estimates how the model
will perform on data the identification did not see.

For models with parameters that lie in a continuous space, a
related discipline is the \engterm{prediction interval} on a
new measurement. The model, fit to past data, predicts the next
measurement with a stated interval; when the next measurement is
taken, the engineer checks whether it falls in the interval at
the stated rate. A 95\% prediction interval that captures the
next measurement 95\% of the time, across many trials, has earned
the percentage. A 95\% prediction interval that captures the
next measurement 70\% of the time has not.

The distinction between a prediction interval and a confidence
interval is the one readers most often collapse, and the collapse
under-reports uncertainty by a definite amount. A confidence
interval bounds where the model's \emph{expected} value lies: it
narrows toward zero as the data accumulate, because more data pin
the fitted curve more tightly. A prediction interval bounds where
the \emph{next single measurement} will fall, and it does not
narrow toward zero, because a new measurement carries its own noise
that no amount of past data removes. For a linear fit the two are
related by
\[
\sigma_{\text{pred}}^2 = \sigma_{\text{noise}}^2
  + \sigma_{\text{fit}}^2,
\]
the sum of the irreducible measurement noise and the shrinking
uncertainty in the fitted value. As the sample grows,
$\sigma_{\text{fit}} \to 0$ and the prediction interval floors at
$\sigma_{\text{noise}}$, the noise the world will always add. An
engineer who reports the confidence interval where a downstream user
needs the prediction interval has promised a precision the next
measurement cannot honour, and the promise fails exactly when a
single new datum, not an average, is what the decision rests on. The
coverage check above tests the right object: it asks whether single
new measurements land inside the interval at the stated rate, which
only the prediction interval claims.

\subsection{Validation is not fitting}

The most consequential discipline error in identification is the
substitution of fitting for validation. A model is \engterm{fit}
when its parameters have been chosen to minimise residuals on a
data set; the model is \engterm{validated} only when its
predictions on data outside the fit have been compared to those
data. The two are different acts, performed at different times,
producing different evidence.

The substitution is easy to make. A reader fits a polynomial to a
data set, observes that the residuals are small, and reports the
polynomial as a validated model. The polynomial may have ten
parameters and the data set ten points; the residuals are small
because the polynomial has interpolated the data. The model has
no validated predictive power. A test on data the polynomial did
not see typically produces residuals orders of magnitude larger
than the training residuals; the reader has reported a model
that does not work.

The discipline is procedural. The validation set is partitioned
off the data \emph{before} the model is chosen. The model is
identified on the training set without reference to the
validation set. The validation score is computed once, after the
model is final. A reader who, after seeing the validation score,
adjusts the model and re-validates has contaminated the
validation; the adjusted model has effectively been fit on the
validation set as well. The procedural defence is the only
defence; the human capacity to rationalise a re-tuned model into
a validated one is unlimited.

\section{Parameter estimation}\pathtag{standard}

Volume II Chapter 14 developed point and interval estimation for
parameters of statistical models. This section recapitulates the
relevant pieces in the modelling context, and adds the
modelling-specific habits that turn an estimation into a usable
input for an engineering decision.

\subsection{Point and interval}

A \engterm{point estimate} is a single value for the parameter. An
\engterm{interval estimate} is a range, with stated confidence or
credibility, that the parameter is supposed to lie in. The
engineer reports both. A point estimate without an interval is a
working number that is not yet ready to leave the desk; an
interval without a point estimate, while occasionally appropriate
in deep-uncertainty problems, is the exception rather than the
rule.

For a parameter $\theta$ estimated by maximum likelihood from a
sample, the point estimate is $\hat\theta$, the value that
maximises the likelihood. Under regularity conditions, the
sampling distribution of $\hat\theta$ is approximately Gaussian
with mean $\theta_{0}$ (the true value) and variance the inverse
Fisher information $I(\theta_{0})^{-1}/n$. A 95\% confidence
interval is approximately $\hat\theta \pm 1.96 \hat{\mathrm{SE}}$
where $\hat{\mathrm{SE}}$ is the estimated standard error.

For a small sample, the Gaussian approximation to the sampling
distribution can be poor. A bootstrap interval, in which the
sampling distribution is approximated by resampling the data with
replacement, often has better coverage. For a parameter
constrained to a range (a probability between 0 and 1, a rate
above 0), the symmetric Gaussian interval can produce values
outside the range; a logit or log transformation followed by an
inverse transformation produces a constrained interval that
respects the parameter's domain.

\subsection{Least squares, maximum likelihood, and what the choice assumes}

The reader has two workhorse identification methods from Chapter 14, and
the modelling context makes the difference between them concrete.
Least squares chooses the parameters that minimise the sum of squared
residuals; maximum likelihood chooses the parameters that maximise the
probability of the observed data under an explicit noise model. The two
coincide exactly when the noise is additive, independent, and Gaussian
with constant variance, because then the log-likelihood is, up to a
constant, the negative sum of squared residuals. The reader who fits by
least squares is fitting by maximum likelihood under that Gaussian
assumption, whether or not the assumption was stated.

The assumption is load-bearing, and naming it is the modelling
discipline. Three common situations break it. When the noise variance
grows with the signal (heteroscedasticity), as it does for count data
and for many sensors near full scale, ordinary least squares gives too
much weight to the large-signal points; weighted least squares, with
each residual divided by its standard deviation, restores the maximum-
likelihood estimate. When the noise is multiplicative rather than
additive, as it often is for concentrations and rates that span decades,
fitting the logarithm of the data converts the multiplicative noise to
additive and makes least squares appropriate again, which is exactly why
the cooling-cup and pendulum fits in this chapter regress a logarithm.
When the noise has heavy tails or outliers, the squared-residual
penalty over-weights the outliers; a robust loss (absolute residuals, or
Huber's blend of quadratic and linear) corresponds to maximum likelihood
under a heavier-tailed noise model and resists the contamination.

The worked contrast is the logistic fit. Cell-count data taken by
plating dilutions carry noise that is roughly proportional to the count,
multiplicative on the linear scale. A naive least-squares fit of $N(t)$
weights the high-count saturation points heavily and the low-count early
points barely, biasing $r$ toward the saturation behaviour. Fitting
$\ln N(t)$, or equivalently weighting each residual by $1/N$, recovers
the early-growth information and returns an $\hat r$ whose interval is
both narrower and honestly placed. The companion script
\texttt{logistic\_fit.py} runs both fits on the same data so the reader
can see the estimate and its interval shift when the noise model is
corrected. The lesson is general: the identification method encodes an
assumption about the noise, and a modeller who fits by least squares
without checking the residual variance has made that assumption silently.

\subsection{Sensitivity and identifiability}

Beyond the point and the interval, the engineer needs a third
quantity: the \engterm{sensitivity} of the prediction to the
parameter. A model whose prediction is very sensitive to a
parameter and whose interval on the parameter is wide is a model
whose predictions are correspondingly uncertain, even if its
point prediction is sharp.

The sensitivity is computed by partial derivative. For a
prediction $y = f(\theta_{1}, \theta_{2}, \ldots, \theta_{p})$,
the sensitivity to $\theta_{i}$ is
\[
S_{i} = \pdv{y}{\theta_{i}}.
\]
A first-order propagation of uncertainty gives
\[
\Var(y) \approx \sum_{i, j} S_{i} S_{j} \Cov(\theta_{i}, \theta_{j}),
\]
which collapses, for independent parameters, to
$\Var(y) = \sum_{i} S_{i}^{2} \Var(\theta_{i})$. Volume I Chapter 7
introduced linear error propagation in the measurement context;
the same algebra applies here, with the parameter intervals from
the identification taking the place of the measurement intervals.

A small sensitivity in absolute terms can still be large in
relative terms if the parameter's interval is wide. The
working diagnostic is the dimensionless \engterm{relative
sensitivity}
\[
\tilde S_{i} = \frac{\theta_{i}}{y} \pdv{y}{\theta_{i}},
\]
which gives the percentage change in $y$ for a one-percent change
in $\theta_{i}$. Parameters with large $\tilde S_{i}$ are the
parameters whose identification matters most; parameters with
small $\tilde S_{i}$ may be poorly identified without harming
predictions.

Identifiability and sensitivity are related but distinct. A
parameter can be identifiable in principle but, in the data the
engineer has, only weakly constrained. The pendulum's $L/g$
ratio is identifiable from period measurements; the period is
much more sensitive to changes in $L/g$ than to changes in any
omitted higher-order term. A model where the most consequential
parameter has a wide interval is a model whose validation
priority is to acquire data that narrow that interval.

% Figure: tornado diagram showing relative sensitivity of beam-deflection
% prediction to each input parameter. Standard visualisation for
% sensitivity analysis.
\begin{figure}[htbp]
\centering
\begin{tikzpicture}[font=\small]
% Centre axis
\draw[thick, draw=engblack!70] (5, 0) -- (5, 5.5);
\node[font=\small\itshape, anchor=south] at (5, 5.6) {nominal $y_{\max}$};
% Horizontal axis below
\draw[->, thick, draw=engblack!70] (0, -0.3) -- (10, -0.3)
  node[right, font=\small] {$\Delta y_{\max} / y_{\max}$};
% Tick labels
\foreach \xt/\xv in {1/-40, 2/-30, 3/-20, 4/-10, 5/0, 6/+10, 7/+20, 8/+30, 9/+40} {
  \draw[draw=engblack!70] (\xt, -0.3) -- (\xt, -0.4);
  \node[font=\scriptsize, anchor=north] at (\xt, -0.42) {\xv\%};
}

% Parameter bars (ordered by sensitivity, largest at top)
% L (sens 4 * 0.5% = 2%): bar from 4.8 to 5.2
% E (sens -1 * 8% = 8%): from 4.2 to 5.8
% I (sens -1 * 3% = 3%): from 4.7 to 5.3
% w (sens 1 * 5% = 5%): from 4.5 to 5.5

% L^4 dominates: 4 * 0.5% = 2%
\fill[engblue!50] (4.8, 4.5) rectangle (5.2, 5.1);
\draw[draw=engblack!70] (4.8, 4.5) rectangle (5.2, 5.1);
\node[anchor=east, font=\footnotesize] at (4.7, 4.8) {$L^{\pm 0.5\%}$};
\node[anchor=west, font=\scriptsize, color=engblack!70] at (5.3, 4.8) {$\pm 2\%$};

% E
\fill[engblue!50] (4.2, 3.5) rectangle (5.8, 4.1);
\draw[draw=engblack!70] (4.2, 3.5) rectangle (5.8, 4.1);
\node[anchor=east, font=\footnotesize] at (4.1, 3.8) {$E^{\pm 8\%}$};
\node[anchor=west, font=\scriptsize, color=engblack!70] at (5.9, 3.8) {$\pm 8\%$};

% w
\fill[engblue!50] (4.5, 2.5) rectangle (5.5, 3.1);
\draw[draw=engblack!70] (4.5, 2.5) rectangle (5.5, 3.1);
\node[anchor=east, font=\footnotesize] at (4.4, 2.8) {$w^{\pm 5\%}$};
\node[anchor=west, font=\scriptsize, color=engblack!70] at (5.6, 2.8) {$\pm 5\%$};

% I
\fill[engblue!50] (4.7, 1.5) rectangle (5.3, 2.1);
\draw[draw=engblack!70] (4.7, 1.5) rectangle (5.3, 2.1);
\node[anchor=east, font=\footnotesize] at (4.6, 1.8) {$I^{\pm 3\%}$};
\node[anchor=west, font=\scriptsize, color=engblack!70] at (5.4, 1.8) {$\pm 3\%$};

% Total bar (root-sum-square): sqrt(2^2 + 8^2 + 5^2 + 3^2) = sqrt(102) ~ 10.1%
\fill[engred!40] (3.99, 0.5) rectangle (6.01, 1.1);
\draw[draw=engred!80, thick] (3.99, 0.5) rectangle (6.01, 1.1);
\node[anchor=east, font=\footnotesize\bfseries] at (3.89, 0.8) {$y_{\max}$ total};
\node[anchor=west, font=\scriptsize, color=engred!80] at (6.11, 0.8) {$\pm 10\%$ (RSS)};

\end{tikzpicture}
\caption[Tornado diagram: parameter sensitivity for the beam-deflection
model]{Relative contribution of each input parameter's interval to the
prediction interval on $y_{\max} = 5 w L^4 / (384 E I)$ for a simply
supported beam under uniform load. Each blue bar is the half-width of
the prediction's relative interval that a $\pm$one-standard-deviation
move in the parameter would produce, computed from
$\tilde S_i = (\theta_i / y) \partial y / \partial \theta_i$ times the
parameter's relative uncertainty. The length scaling ($L^4$) means
$L$'s small interval of $\pm 0.5\%$ contributes $\pm 2\%$, while $E$'s
relatively loose $\pm 8\%$ contributes $\pm 8\%$. The root-sum-square
total (red) is the propagated interval on the prediction, here
$\pm 10\%$; the visualisation makes the ranking explicit and tells the
engineer that improving $E$'s estimate is the highest-yield move.
Exercise 18.5 asks the reader to construct this diagram.}
\label{fig:vol02:ch18:sensitivity-tornado}
\end{figure}


Figure~\ref{fig:vol02:ch18:sensitivity-tornado} is the standard
visual report for a sensitivity analysis: a tornado diagram for the
beam-deflection prediction $y_{\max} = 5 w L^{4} / (384 E I)$, with
bars ranked by each parameter's contribution to the propagated
relative interval. The reader notices immediately that $L$, whose
relative interval is only $\pm 0.5\%$, contributes $\pm 2\%$ to the
prediction (because of the fourth power); and that $E$, with its
loose $\pm 8\%$, contributes the same $\pm 8\%$ undiluted (because
of the first power). The diagram tells the engineer where to spend
the next measurement budget. The Saltelli primer
\cite{text:v2ch18-saltelli-sensitivity-2008} extends the approach to
global, model-free indices; the local-derivative version shown here
is the right tool when the model is differentiable, the parameter
intervals are small, and the engineer needs the rank order rather
than full nonlinear coverage. The companion script
\texttt{sensitivity\_propagation.py} computes both the analytic and
the finite-difference sensitivities for any of the chapter's
worked-example models.

\begin{mastery}
The local-derivative sensitivity of this section is the right tool when
the model is smooth, the parameter intervals are small, and the engineer
wants a rank order. Three situations break those conditions, and each has
a standard heavier instrument. When the model is nonlinear and the
parameter intervals are wide, the first-order propagation
$\Var(y) = \sum S_i^2 \Var(\theta_i)$ understates the true output spread,
and the engineer turns to \engterm{Monte-Carlo uncertainty
quantification}: draw thousands of parameter vectors from their joint
distribution, run the model on each, and report the empirical distribution
of $y$ directly. The method makes no linearity assumption and captures the
full output distribution, including its skew and its tails, at the cost of
many model evaluations. When the engineer needs not just the output spread
but the share of it each parameter causes, including the share carried by
interactions the tornado diagram cannot show, the \engterm{Sobol indices}
decompose the output variance into contributions from each parameter and
each combination, the global, model-free generalisation of the relative
sensitivity, which Saltelli's primer
\cite{text:v2ch18-saltelli-sensitivity-2008} treats in full. When each
model evaluation is expensive, a finite-element run of hours or a field
trial of weeks, the Monte-Carlo and Sobol calculations become
unaffordable, and the engineer fits a \engterm{surrogate model}: a cheap
approximation (a polynomial response surface, a Gaussian-process emulator)
trained on a modest design of expensive runs and then queried thousands of
times in their place. The surrogate is itself a model, with its own
identification, validation, and domain of validity, and the discipline of
this chapter applies to it unchanged: a surrogate trusted outside the
design region from which it was trained is the regime-extension fallacy
wearing a faster coat. The progression from local derivative to
Monte-Carlo to surrogate-accelerated global analysis is a progression in
cost and in fidelity, and the modeller's judgement is to climb it only as
far as the decision requires.
\end{mastery}

\subsection{The estimate as part of a working report}

A modelling report includes, for each parameter, the point
estimate, the interval, the method by which the interval was
constructed, the data on which the identification was performed,
the sensitivity of the prediction to the parameter, and the
identifiability check. Volume I Chapter 8 introduced the
working-report discipline for measurements; the same discipline
extends to modelling. A model whose parameter estimates are
reported without intervals is a model whose users will treat
the point estimates as if they were known to many decimal
places, and the model will make predictions to that false
precision. The engineer who states the interval has done the
work that a downstream user would otherwise do badly.

The report also carries the correlations among the parameters, not
only their marginal intervals, and the omission of the correlations
is a quiet source of wrong downstream predictions. Two parameters
reported as $\hat\theta_1 = 3.0 \pm 0.4$ and
$\hat\theta_2 = 5.0 \pm 0.6$ tell a downstream user nothing about
whether the two move together or in opposition, yet a prediction that
depends on their difference $\theta_1 - \theta_2$ has variance
$\sigma_1^2 + \sigma_2^2 - 2\,\mathrm{Cov}(\theta_1, \theta_2)$, which
a strong positive correlation can shrink to a fraction of the
independent estimate, or a negative one can inflate. The identification
returns the full covariance matrix; the working report either passes
it forward or states the correlation coefficients alongside the
marginal intervals, so the downstream user propagates the right
quantity. The habit matters most for exactly the models the chapter
has flagged as trade-off prone, where a correlation near $\pm 1$
means the marginal intervals badly misrepresent what the data actually
constrain. A report that lists only the marginals, for such a model,
has hidden the ridge the covariance would have exposed, and the
downstream user who treats the parameters as independent will compute
a prediction interval that is wrong in a direction the report gave no
way to detect.

\begin{estimation}
\textbf{Question.} A reader fits the simple pendulum equation
$T^{2} = (4\pi^{2}/g) L$ to five measurements of period at known
length, taken with a stopwatch (period uncertainty
$0.05\,\si{\second}$) and a tape measure (length uncertainty
$2\,\si{\milli\meter}$). What confidence interval on $g$ does the
reader expect, before doing the fit?

\smallskip
\textit{Estimate, before the calculation.}

\smallskip
The period uncertainty is $0.05\,\si{\second}$ on a period of order
$2\,\si{\second}$, so $\sigma_{T}/T \approx 2.5\,\%$. The length
uncertainty is $2\,\si{\milli\meter}$ on a length of order
$1\,\si{\meter}$, so $\sigma_{L}/L \approx 0.2\,\%$. The relation
$g = 4\pi^{2} L / T^{2}$ propagates these to
\[
\frac{\sigma_{g}}{g} \approx \sqrt{\left(\frac{\sigma_{L}}{L}\right)^{2}
  + \left(2 \frac{\sigma_{T}}{T}\right)^{2}}
\approx \sqrt{0.002^{2} + 0.05^{2}} \approx 5\,\%,
\]
dominated by the period measurement (the factor of two from the
square). With five measurements, the standard error scales as
$1/\sqrt{5}$, so the per-fit interval shrinks to about $2\,\%$,
giving a 95\% confidence half-width of roughly $4\,\%$ or
$\pm 0.4\,\si{\meter\per\square\second}$.

\smallskip
\textbf{Postmortem.} The dominant factor is the period
measurement, despite its small absolute uncertainty, because the
period enters the relation squared. A reader who wants a tighter
interval should reduce the period uncertainty (longer pendulum,
multiple-period timing, automated timing) before reducing the
length uncertainty. The estimate identifies the lever before any
data are taken; the discipline that this volume installs is the
discipline of identifying that lever first.
\end{estimation}

\section{Model uncertainty and structural error}\pathtag{core}

A parameter interval describes uncertainty about the value of a
parameter inside a fixed model. A larger source of uncertainty is
the choice of model itself: the parametric form, the variables
included, the assumptions about noise, the boundary conditions.
This source is called \engterm{model uncertainty} or
\engterm{structural error}, and a working engineer who treats it
honestly produces predictions whose intervals are wider, whose
domain of validity is narrower, and whose predictive performance
on new data is closer to what the model claims.

\subsection{The model is one of many}

For most engineering questions, several candidate models could be
fit to the same data. A linear regression and a quadratic
regression and a piecewise-linear regression and an exponential
regression can all be fit to a data set; each will produce a
point estimate and an interval. The intervals from any one of
these models do not include the disagreement among models. A
reader who reports the linear-regression interval as the
prediction's uncertainty has confused the within-model interval
with the across-model interval.

The standard remedy is \engterm{model averaging}: fit several
candidate models, weight them by their evidence (information
criterion, Bayesian posterior probability, cross-validation
score), and report the combined prediction with an interval that
includes the across-model variance. The remedy is honest and
expensive; it requires fitting many models and assigning weights
that the data may not strongly support.

A simpler remedy is \engterm{model triangulation}: fit two or
three structurally distinct models, compare their predictions,
and take the disagreement as a lower bound on the structural
uncertainty. Where the models agree, the prediction is insensitive to
structural assumptions; where they disagree, the prediction is
sensitive to which model is used, and the engineer reports the
range. Triangulation does not require the formal weights of model
averaging and produces a defensible report from a small number of
fits. David MacKay's \emph{Sustainable Energy: Without the Hot Air}
\cite{gen:mackay2009} demonstrates the discipline at book length on
energy questions: the same quantity, computed by independent
back-of-envelope routes, becomes a triangulated estimate whose
disagreement is itself the report.

\subsection{Model selection: choosing among candidates by penalised fit}

Triangulation reports the disagreement among models; selection chooses
one. When the candidates are nested in complexity, a polynomial of
increasing degree, a compartment model of increasing compartment count, a
regression with increasing covariates, the choice is governed by a
trade-off the chapter has named in pieces and now states whole. A richer
model always fits the training data at least as well, because the simpler
model is a special case the richer one can reproduce. The training
residual therefore falls monotonically with complexity and cannot, by
itself, select a model; a rule that minimised training residual would
always choose the most complex candidate and overfit.

The selection rule penalises complexity. The \engterm{Akaike information
criterion} adds to the negative log-likelihood a penalty of $2p$ for a
model with $p$ free parameters; the \engterm{Bayesian information
criterion} uses a penalty of $p \ln n$ that grows with the sample size
$n$ and so favours simpler models more strongly in large samples. Both
turn the monotone training curve into a U-shape with an interior minimum,
the simplest model whose added parameters still buy enough fit to pay for
their penalty. Cross-validation reaches the same end by a different route,
estimating held-out performance directly rather than penalising parameter
count, and on most problems the two agree on where the optimum sits even
when they disagree on the score.
Figure~\ref{fig:vol02:ch18:model-selection} shows the three curves
together: the training residual falling toward zero, the penalised and
held-out scores turning over at the optimum $p^\star$, the underfit regime
to its left and the overfit regime to its right.

The worked instance is the cooling cup. The Newton's-law model has one
parameter ($k$); a two-term model adding a radiative correction has two; a
three-term model adding an evaporative loss has three. Fitting all three
to the thirty-point notebook, the training residual falls from one model
to the next, as it must. The AIC, computed as $n \ln(\text{RSS}/n) + 2p$,
is lowest for the two-parameter model on a hot-start notebook (the
radiative term earns its parameter) but lowest for the one-parameter model
on a notebook that never exceeded $50\,\degC$ excess (the radiative term
is below the noise, and its parameter is not paid for). The criterion
encodes the chapter's Rule Five mechanically: the simplest model that
answers the question is the simplest model the data can distinguish from
its richer rivals, and AIC names that model by number rather than by
judgement. Box's response-surface monograph
\cite{text:v2ch18-box-empirical-1987} and the statistical-learning text
\cite{text:v2ch18-hastie-statistical-learning-2009} treat the
information-criterion and cross-validation machinery in the depth the
working engineer occasionally needs; the chapter's claim is the narrow
one that selection must read the optimum off a complexity-penalised curve,
never off the training residual.

% Figure: model-selection scores against model complexity, showing the
% in-sample fit improving monotonically while the penalised/held-out
% scores turn over at an interior optimum.
\begin{figure}[htbp]
\centering
\begin{tikzpicture}[font=\small]
\draw[->, thick, draw=engblack!70] (0,0) -- (9.0,0)
  node[right, font=\scriptsize] {model complexity (parameters $p$)};
\draw[->, thick, draw=engblack!70] (0,0) -- (0,4.6)
  node[above, font=\scriptsize] {score (lower is better)};
\foreach \xt/\xv in {1/1, 2/2, 3/3, 4/4, 5/5, 6/6, 7/7, 8/8} {
  \draw[draw=engblack!70] (\xt,0) -- (\xt,-0.07);
  \node[font=\scriptsize, anchor=north] at (\xt,-0.09) {\xv};
}
% In-sample residual (training error): monotone decreasing toward zero
\draw[very thick, draw=enggray, dashed, smooth, samples=80, domain=1:8]
  plot (\x, { 0.35 + 3.6*exp(-0.55*(\x-1)) });
\node[anchor=west, font=\scriptsize, text=enggray] at (6.1,0.55)
  {training residual};
% AIC: training fit + 2p penalty -> U-shape, min near p=3
\draw[very thick, draw=engblue!85, smooth, samples=120, domain=1:8]
  plot (\x, { 0.35 + 3.6*exp(-0.55*(\x-1)) + 0.30*(\x-1) });
\node[anchor=west, font=\scriptsize, text=engblue!85] at (1.0,4.2) {AIC};
% Cross-validation error: U-shape, min near p=3, slightly above AIC at right
\draw[very thick, draw=engred!85, smooth, samples=120, domain=1:8]
  plot (\x, { 0.55 + 3.4*exp(-0.55*(\x-1)) + 0.42*(\x-1)*(\x-1)*0.12 });
\node[anchor=west, font=\scriptsize, text=engred!85] at (5.9,2.9)
  {CV error};
% mark the optimum near p=3
\draw[dashed, draw=engblack!45] (3,0) -- (3,1.85);
\fill[engblack!75] (3,1.85) circle (1.8pt);
\node[font=\scriptsize, anchor=south] at (3,1.9) {$p^\star$};
% underfit / overfit zones
\node[font=\scriptsize, anchor=south, text=engamber] at (1.5,0.05)
  {underfit};
\node[font=\scriptsize, anchor=south, text=engamber] at (7.0,0.05)
  {overfit};
\end{tikzpicture}
\caption[Model selection: in-sample fit versus penalised and held-out
score]{Three scores as the model's parameter count grows. The training
residual (dashed) falls monotonically: a richer model always fits the
data it was fit on at least as well. That curve alone would select the
most complex model and overfit. The Akaike information criterion (blue)
adds a penalty proportional to the parameter count, and the
cross-validation error (red) estimates performance on held-out data
directly; both turn over at an interior optimum $p^\star$, the simplest
model whose added complexity still buys predictive accuracy. To the left
of $p^\star$ the model underfits (it omits structure the data contain);
to the right it overfits (it fits noise as signal). The two penalised
curves usually agree on the location of $p^\star$ even when they disagree
on the score; the engineering rule is to read the optimum off a curve that
penalises complexity, never off the training residual.}
\label{fig:vol02:ch18:model-selection}
\end{figure}


% Full-page figure: model triangulation. Four independent routes to the
% same engineering quantity (UK personal-transport energy demand, after
% MacKay), with their values and intervals plotted on a common axis,
% showing both agreement and disagreement.
\begin{figure}[p]
\centering
\begin{tikzpicture}[font=\small]

% ===== Title block =====
\node[font=\large\bfseries, align=center] at (5.5, 13.0)
  {Four routes to one quantity: model triangulation};
\node[font=\small\itshape, align=center, color=engblack!70] at (5.5, 12.4)
  {Annual personal-transport energy demand per UK resident, current as of 2008-09.};
\node[font=\small\itshape, align=center, color=engblack!70] at (5.5, 12.0)
  {Four independent computation routes; the spread of their estimates is the structural-uncertainty floor.};

% ===== Route boxes =====

% Route 1
\node[draw=engblack, fill=engblue!10, rounded corners, minimum width=10cm,
  minimum height=1.9cm, align=left, anchor=north west, font=\footnotesize] at (0.5, 11.0)
  {\textbf{Route 1: vehicle-kilometre $\times$ energy intensity.} \\
   $E_1 = D \cdot \eta = (13{,}500\,\text{km}/\text{yr})(0.80\,\text{kWh}/\text{km})
   = 10{,}800\,\text{kWh}/\text{yr}.$\\
   $D$ from UK National Travel Survey 2009; $\eta$ from typical petrol-car
   energy-per-km (cars only). \\
   \textit{Uncertainty: $\pm 20\%$ from $\eta$ spread across vehicle fleet.}};

% Route 2
\node[draw=engblack, fill=engblue!10, rounded corners, minimum width=10cm,
  minimum height=1.9cm, align=left, anchor=north west, font=\footnotesize] at (0.5, 8.9)
  {\textbf{Route 2: fuel sales $\div$ population.}\\
   $E_2 = F / N = (44 \cdot 10^9\,\text{L petrol-equiv}/\text{yr})
   (9.7\,\text{kWh}/\text{L}) / (61 \cdot 10^6) = 7{,}000\,\text{kWh}/\text{yr}.$\\
   $F$ from UK Department for Transport fuel-sales tables; $N$ from ONS.\\
   \textit{Uncertainty: $\pm 5\%$ from sales aggregation; lower than Route 1.}};

% Route 3
\node[draw=engblack, fill=engblue!10, rounded corners, minimum width=10cm,
  minimum height=1.9cm, align=left, anchor=north west, font=\footnotesize] at (0.5, 6.8)
  {\textbf{Route 3: travel-time budget $\times$ average power.}\\
   $E_3 = t \cdot P = (1.1\,\text{h}/\text{day})(365)(15\,\text{kW}\,\text{avg per car})/4
     = 4{,}500\,\text{kWh}/\text{yr per person}.$\\
   $t$ from time-use surveys; $P$ from typical engine load; factor 4 occupancy correction.\\
   \textit{Uncertainty: $\pm 30\%$ from occupancy and power assumptions.}};

% Route 4
\node[draw=engblack, fill=engblue!10, rounded corners, minimum width=10cm,
  minimum height=1.9cm, align=left, anchor=north west, font=\footnotesize] at (0.5, 4.7)
  {\textbf{Route 4: CO\textsubscript{2} emissions $\div$ carbon intensity.}\\
   $E_4 = (M_\text{CO\textsubscript{2},transport} / I_\text{CO\textsubscript{2}}) / N
     \approx 8{,}000\,\text{kWh}/\text{yr}.$\\
   $M$ from UK national GHG inventory (transport sector minus freight);
   $I$ from fuel carbon content.\\
   \textit{Uncertainty: $\pm 15\%$ from sectoral attribution.}};

% ===== Common axis with intervals =====
\draw[->, thick, draw=engblack!70] (1, 2.5) -- (10, 2.5)
  node[right, font=\small] {$E$ (kWh/yr per resident)};
\foreach \xt/\xv in {2/2000, 3.5/4000, 5/6000, 6.5/8000, 8/10000, 9.5/12000} {
  \draw[draw=engblack!70] (\xt, 2.5) -- (\xt, 2.4);
  \node[font=\scriptsize, anchor=north] at (\xt, 2.35) {\xv};
}

% Interval bars on the axis
% Route 1: 10800 +/- 20% => bar from 8640 to 12960, x-coord = 1 + 9 * E/12000
% map: x = 1 + (E/12000)*9
\foreach \E/\sig/\col/\label in {
  10800/2160/engblue!85/{R1},
  7000/350/engblue!85/{R2},
  4500/1350/engblue!85/{R3},
  8000/1200/engblue!85/{R4}} {
  \pgfmathsetmacro\xc{1 + (\E/12000)*9}
  \pgfmathsetmacro\xl{1 + ((\E-\sig)/12000)*9}
  \pgfmathsetmacro\xr{1 + ((\E+\sig)/12000)*9}
  \draw[thick, draw=\col] (\xl, 1.7) -- (\xr, 1.7);
  \draw[thick, draw=\col] (\xl, 1.6) -- (\xl, 1.8);
  \draw[thick, draw=\col] (\xr, 1.6) -- (\xr, 1.8);
  \filldraw[\col] (\xc, 1.7) circle (0.08);
  \node[font=\scriptsize, anchor=south, color=\col] at (\xc, 1.85) {\label};
}

% Annotation: triangulation envelope
\draw[dashed, draw=engred!80, thick] (4.375, 1.3) -- (4.375, 2.0);
\draw[dashed, draw=engred!80, thick] (8.6, 1.3) -- (8.6, 2.0);
\node[font=\footnotesize, color=engred!80, align=center, anchor=north] at (6.5, 1.25)
  {triangulation envelope: $\sim$4{,}500 to $\sim$11{,}000 kWh/yr,\\
   the structural-uncertainty range the four-route comparison reports};

% Conclusion box
\node[draw=engblack!60, fill=engblack!5, rounded corners,
  minimum width=10cm, minimum height=1.1cm, align=left, anchor=north west,
  font=\footnotesize] at (0.5, 0.9)
  {\textbf{Reading.} The four routes agree on the order of magnitude
   ($\sim$$10^4$ kWh/yr per resident) and disagree on the factor-of-two
   level. The disagreement is the structural floor on any single-route
   estimate. A reader who reports only Route 2's tight $\pm 5\%$
   interval has under-reported the modelling uncertainty by a factor of
   five, because Route 2's interval is the within-model interval; the
   across-route interval is what the planner using the number needs.};

\end{tikzpicture}
\caption[Model triangulation: four routes to UK personal-transport
energy demand]{Four independent computation routes to a single
engineering quantity: annual personal-transport energy demand per
UK resident, current as of 2008-09. Each route's value and within-route
uncertainty are plotted on the common axis; the dashed red envelope
marks the spread across routes, which sets the lower bound on the
structural uncertainty of any single-route estimate. The across-route
spread is six times Route~2's reported within-route interval, which is
the chapter's point: a model whose interval is reported without
triangulation under-reports the uncertainty the decision-maker
actually faces. The example is drawn from David MacKay's
\emph{Sustainable Energy: Without the Hot Air}
\cite{gen:mackay2009}, which builds the triangulation discipline at
book length on energy questions.}
\label{fig:vol02:ch18:triangulation}
\end{figure}


Figure~\ref{fig:vol02:ch18:triangulation} reproduces one of MacKay's
triangulation tableaus: four routes to the same quantity (UK
personal-transport energy per resident, current as of 2008-09), each
with its within-route interval, the across-route spread overlaid as
the structural-uncertainty envelope. The within-route intervals
disagree by less than $\pm 30\%$ individually; the across-route
spread covers a factor of more than two. The triangulated number a
planner can rely on is the across-route envelope; the single-route
interval is the within-model report, which understates what the
decision-maker faces by roughly a factor of five. The discipline
generalises: a single model whose interval is reported without
triangulation is reporting one component of a multi-component
uncertainty.

\subsection{Structural error is what the model cannot say}

A model has a residual: the difference between the model's
prediction and the measurement. The residual contains two pieces.
The first is measurement noise, which the model treats by
construction. The second is structural error, which the model
does not treat. A model whose residuals show structure (a trend,
a periodicity, a heteroscedasticity, a correlation with a
variable not in the model) has not captured the structure; the
residual structure is what the model cannot say.

The diagnostic is mechanical and was introduced in Volume II
Chapter 14. The residuals are plotted against each input
variable, against the prediction itself, and against any
candidate omitted variable. A residual plot with no structure is
a residual plot consistent with the model's noise assumptions.
A residual plot with a clear trend is a residual plot demanding
that the model be revised.

The discipline is to revise the model, not to enlarge the
interval. A reader who notices a trend in the residuals and
inflates the interval to cover the trend has reported an honest
interval but has discarded the structural information the
residuals contained. The structural information is the next
model's specification: the trend that the current model misses
is the term the next model needs.

\subsection{The honest interval}

A modelling report's prediction interval should reflect both
parameter uncertainty and structural uncertainty. The reader
who reports only parameter uncertainty produces an interval that
will be overconfident on new data: the next measurement, drawn
from a regime where the structural error is non-negligible,
will fall outside the interval more often than the interval's
percentage suggests.

The honest interval inflates the parameter-uncertainty interval by
a factor that reflects the structural uncertainty's contribution.
Where the structural uncertainty is small, the inflation is
small; where the structural uncertainty is large, the inflation
is large. The inflation factor can be estimated from the
discrepancy between training and validation residuals: a model
whose validation residuals are twice the training residuals has a
structural-uncertainty contribution roughly equal to the
parameter contribution, and the honest interval is roughly
$\sqrt{2}$ times the within-model interval.

The discipline appears in the engineering-statistics literature
under the name \engterm{predictive variance}: the variance of a
new measurement under the model's posterior, including parameter
uncertainty and the model's noise term. A predictive variance
that has been augmented for structural error gives the engineer
the interval that the engineer's downstream user actually needs.
Saltelli and coauthors' sensitivity-analysis primer
\cite{text:v2ch18-saltelli-sensitivity-2008} and Hastie, Tibshirani,
and Friedman's statistical-learning text
\cite{text:v2ch18-hastie-statistical-learning-2009} together cover the
modern apparatus the engineer can reach for; both books are
recommended companions for the reader who wants to push past the
chapter's depth.

\begin{estimation}
\textbf{Question.} A reader has fit a cooling-cup model to seven
points and obtained $\hat k = 0.04\,\si{\per\minute}$ with a 95\%
parameter interval of $\pm 12\%$. The reader then runs five-fold
cross-validation and finds the validation-fold residuals have an RMS
of $0.5\,\degC$ against an in-sample residual RMS of
$0.3\,\degC$. By what factor should the prediction interval on
$t_{60}$ be inflated to account for structural error?

\smallskip
\textit{Estimate, before any formal calculation.}

\smallskip
The ratio of validation to in-sample RMS is approximately $1.67$. The
structural-error variance contributes $\sigma_{s}^{2} = \sigma_{v}^{2}
- \sigma_{p}^{2}$ where $\sigma_{v}$ is the validation RMS and
$\sigma_{p}$ the in-sample RMS, on the assumption that the two are
approximately independent. Plugging in: $\sigma_{s}^{2} \approx
(0.5)^{2} - (0.3)^{2} = 0.16$, so $\sigma_{s} \approx 0.4\,\degC$.
The predictive standard deviation is $\sigma_{\text{pred}} =
\sqrt{\sigma_{p}^{2} + \sigma_{s}^{2}} \approx
\sqrt{0.09 + 0.16} = 0.5\,\degC$, which is $\sqrt{0.25/0.09}
\approx 1.67$ times $\sigma_{p}$. The parameter-only interval should
be widened by a factor of about $1.7$ to honestly cover new
measurements. The $\pm 12\%$ parameter-only interval on $\hat k$
propagates into a $\pm 1.8\,\text{min}$ interval on $t_{60}$ (from
$\tau = 1/\hat k$); the honest interval is $\pm 3.1\,\text{min}$.

\smallskip
\textbf{Postmortem.} The inflation factor is the modeller's audit on
the parameter-only report. A factor of $1$ means the model's noise
assumption is calibrated; a factor near $2$ means the model is
missing structure on the order of the noise; a factor of $5$ or
more means the model is wrong and should be revised before any
prediction interval is reported. The companion script
\texttt{cross\_validation.py} computes this factor for the
cooling-cup data automatically; the reader's own measurement
notebooks should be exercised the same way before the project
write-up declares an interval honest.
\end{estimation}

\subsection{The polya discipline: solve a related problem}

George P{\'o}lya's \emph{How to Solve It} \cite{gen:polya1945} is
the canonical short reference on problem-solving heuristics, and
the heuristics translate directly into modelling decisions. We
name two.

\engterm{Solve a related problem first}. When a model is hard,
identify a simpler model whose solution the modeller can write
down, and solve that. The solution gives a benchmark against which
the harder model can be checked, and a starting point from which
the harder model's behaviour can be understood as a perturbation.
The pendulum's small-angle linearisation is exactly this move: the
nonlinear pendulum is hard; the linearised pendulum is the related
problem the modeller can solve; the nonlinear corrections are the
perturbation. The engineer who attempts the nonlinear problem
without the linear benchmark has discarded the diagnostic the
benchmark would have provided.

\engterm{Look back}. After a model is built, the modeller asks:
does the result use all the data? Does it explain the
discrepancies the simpler model could not explain? Does it admit
a check by an independent route? The look-back move is what
catches the model that fit the training data without learning
anything about the system. The look-back is procedural, not
analytical; the discipline is to perform it before reporting the
model, not after a downstream user has reported a problem.

\section{Worked example: from a measurement notebook to a model and back}\pathtag{core}

We work an example end-to-end. The example uses the cooling-cup
measurement from Volume I Chapter 4, which produced a notebook of
temperatures recorded every minute as a cup of hot water cooled
toward room temperature. The reader who has done the Volume I
project has the data; the reader who has not can take new data in
half an hour with a thermometer and a kitchen timer.

\subsection{The data}

The notebook records $T_{i}$ at times $t_{i} = 0, 1, 2, \ldots, 30$
minutes, with $T_{0} = 90\,\degC$, $T_{30}$ approaching the room
temperature $T_{\infty} \approx 22\,\degC$. The thermometer's
resolution is $0.1\,\degC$ and its accuracy, from Volume I
Chapter 3's calibration, is $\pm 0.5\,\degC$. The timer is
accurate to a second on a minute scale.

\subsection{The candidate model}

Newton's law of cooling models the rate of cooling as proportional
to the temperature difference between the body and its environment:
\[
\dv{T}{t} = -k (T - T_{\infty}),
\]
which integrates to
\[
T(t) = T_{\infty} + (T_{0} - T_{\infty}) e^{-kt}.
\]
The model has one identifiable parameter, $k$, with units of
inverse time, and assumes that $T_{\infty}$ is known and that the
heat transfer coefficient between cup and air is constant over the
cooling range.

\subsection{Identification}

The reader linearises the model by taking the logarithm of the
temperature excess:
\[
\ln(T(t) - T_{\infty}) = \ln(T_{0} - T_{\infty}) - kt.
\]
A linear regression of $\ln(T_{i} - T_{\infty})$ on $t_{i}$
returns a slope estimate $-\hat k$ with standard error
$\hat{\mathrm{SE}}(\hat k)$. For our notebook,
$\hat k \approx 0.084\,\si{\per\minute}$ with a 95\% confidence
interval of roughly $\pm 0.005\,\si{\per\minute}$ from the
regression standard error. The reader records the value, the
interval, the method, and the data set.

\subsection{Validation}

Validation requires data the identification did not use. The
reader takes a new cup, of the same shape and volume, and records
$T(t)$ for a fresh cooling. Newton's-law prediction with
$\hat k = 0.084\,\si{\per\minute}$ gives
$T(15) \approx 22 + 68 e^{-1.26} \approx 41.3\,\degC$. The
measured value, on a representative trial, is $42.1\,\degC$. The
discrepancy of $0.8\,\degC$ is within the thermometer's accuracy;
the model's prediction has been validated against held-out data.

A more demanding validation looks at the residuals from the fit
itself. The residuals plotted against $t$ show no obvious trend;
plotted against $T - T_{\infty}$ they show a faint pattern, with
slightly larger residuals at large temperature excess. The pattern
suggests that the heat transfer coefficient is not perfectly
constant: at large $T - T_{\infty}$, the convective transfer
is slightly enhanced by buoyancy. The pattern is small in our
data; in a hotter cup, or one cooling in still air, it would
matter more.

% Figure: model-vs-data, with prediction band and residual panel.
% Cooling cup: Newton's law fit, validation residuals below.
\begin{figure}[htbp]
\centering
\begin{tikzpicture}[font=\small]
% ===== Top panel: data and fit =====
\begin{scope}[yshift=0cm]
% Axes
\draw[->, thick, draw=engblack!70] (0,0) -- (8.5, 0) node[right, font=\small] {$t$ (min)};
\draw[->, thick, draw=engblack!70] (0,0) -- (0, 4.5) node[above, font=\small] {$T$ ($\si{\degreeCelsius}$)};
% Ticks: t = 0, 5, 10, 15, 20, 25, 30 mapped to x = 0..7
\foreach \xt/\xv in {0/0, 1.17/5, 2.33/10, 3.5/15, 4.67/20, 5.83/25, 7.0/30} {
  \draw[draw=engblack!70] (\xt, 0) -- (\xt, -0.08);
  \node[font=\scriptsize, anchor=north] at (\xt, -0.1) {\xv};
}
% Ticks: T = 20, 40, 60, 80, 90; with T_inf = 22 at y = 0
\foreach \yt/\yv in {0/22, 1.0/40, 2.13/60, 3.27/80, 3.83/90} {
  \draw[draw=engblack!70] (0, \yt) -- (-0.08, \yt);
  \node[font=\scriptsize, anchor=east] at (-0.1, \yt) {\yv};
}
% Prediction band (Newton's law, k=0.084 +/- 0.005); shaded region
\fill[engblue!10, opacity=0.7]
  plot[smooth, samples=40, domain=0:7]
    (\x, {(90-22)*exp(-0.079*\x*30/7) * 0.057}) --
  plot[smooth, samples=40, domain=7:0]
    (\x, {(90-22)*exp(-0.089*\x*30/7) * 0.057}) -- cycle;
% Best-fit Newton's law curve (k=0.084 per min, mapped)
\draw[thick, draw=engblue!80]
  plot[smooth, samples=60, domain=0:7]
    (\x, {68 * exp(-0.084 * \x * 30/7) * 0.0567});
% Data points
\foreach \xt/\Tval in {0/90.0, 1.17/70.4, 2.33/55.8, 3.5/45.0, 4.67/37.0, 5.83/31.2, 7.0/27.0} {
  \pgfmathsetmacro\yval{(\Tval-22) * 0.0567}
  \filldraw[engred!90] (\xt, \yval) circle (0.07);
}
% T_inf line
\draw[dashed, draw=engblack!50] (0, 0) -- (7.5, 0) node[right, font=\scriptsize] {$T_\infty$};
% Legend
\node[anchor=west, font=\scriptsize] at (4.6, 3.6) {\textcolor{engred!90}{$\bullet$} data};
\node[anchor=west, font=\scriptsize] at (4.6, 3.3) {\textcolor{engblue!80}{(line)} fit ($\hat k = 0.084$)};
\node[anchor=west, font=\scriptsize] at (4.6, 3.0) {\textcolor{engblue!50}{$\blacksquare$} 95\% interval};
\node[font=\small, anchor=south west] at (0.2, 4.3) {\textbf{(a) Fit and data}};
\end{scope}

% ===== Bottom panel: residuals =====
\begin{scope}[yshift=-3.0cm]
\draw[->, thick, draw=engblack!70] (0,0) -- (8.5, 0) node[right, font=\small] {$t$ (min)};
\draw[->, thick, draw=engblack!70] (0,-0.8) -- (0, 0.8) node[above, font=\small] {residual ($\si{\degreeCelsius}$)};
\draw[dashed, draw=engblack!50] (0, 0) -- (7.5, 0);
% Tick labels
\foreach \yt/\yv in {-0.6/-1.0, 0.0/0.0, 0.6/+1.0} {
  \draw[draw=engblack!70] (0, \yt) -- (-0.08, \yt);
  \node[font=\scriptsize, anchor=east] at (-0.1, \yt) {\yv};
}
% Synthetic residuals (small, slight curvature: model misses radiative term)
\foreach \xt/\rval in {0/0.5, 1.17/-0.1, 2.33/-0.3, 3.5/-0.2, 4.67/0.05, 5.83/0.15, 7.0/-0.1} {
  \filldraw[engred!90] (\xt, {\rval * 0.6}) circle (0.06);
}
\node[font=\small, anchor=south west] at (0.2, 0.85) {\textbf{(b) Residuals}};
\node[anchor=west, font=\scriptsize, color=engblack!70] at (4.6, -0.55)
  {faint trend at large $t$: radiative term};
\end{scope}
\end{tikzpicture}
\caption[Cooling-cup fit and validation residuals]{(a) Cooling-cup
notebook (seven points, fifteen-minute increments) with the Newton's-law
fit $T(t) = T_\infty + (T_0 - T_\infty) e^{-\hat k t}$ overlaid in
blue, the 95\% prediction band shaded. The fitted decay constant is
$\hat k = 0.084\,\si{\per\minute}$ with interval $\pm
0.005\,\si{\per\minute}$. (b) Residuals on the same time axis: their
amplitude is within the thermometer's accuracy ($\pm 0.5\,\degC$), but
the faint trend at small $t$ suggests the model under-predicts cooling
when the temperature excess is large (the radiative correction the model
omits). The trend is the diagnostic the model leaves the modeller; the
discipline of section~18.5 reads it as the next model's specification.}
\label{fig:vol02:ch18:validation-residuals}
\end{figure}


Figure~\ref{fig:vol02:ch18:validation-residuals} shows the fit and
the residuals as the modelling report would carry them. The
residuals subpanel is the part the engineer learns from: the
amplitude is within thermometer accuracy, but a faint trend at small
$t$ marks the regime in which the radiative term begins to bite.
The companion script \texttt{cooling\_fit.py} reproduces the fit and
its propagated interval on $t_{60}$; \texttt{cross\_validation.py}
runs the $k$-fold split and reports the inflation factor by which
the parameter-only interval understates the structural uncertainty.

\subsection{The model's structural error}

What does the model leave out? Several things. The cup's geometry
is not in the model: a tall thin cup cools differently from a
shallow wide cup. The radiative cooling at high temperature is not
in the model: the Stefan-Boltzmann term contributes about
$\sigma (T^{4} - T_{\infty}^{4})$ rather than $T - T_{\infty}$,
which becomes significant above about $50\,\degC$ excess. The
evaporative cooling, where the water surface is exposed and the
air is dry, is not in the model: the latent heat of vaporisation
is large. The thermometer's thermal mass is not in the model: the
thermometer reads its own equilibrated temperature, not the
water's instantaneous temperature, which produces a small phase
lag at fast cooling.

Each omission is small for the regime our notebook covers. Each
becomes large in a different regime. The reader who states the
omissions has stated the model's domain of validity; the reader
who would apply the model to a kettle cooling from boiling, or a
thin film of water evaporating, or a cup in a refrigerator, knows
which omissions become large in those regimes.

\subsection{Sensitivity and the engineering question}

The engineering question motivating the model is: how long until
the cup is cool enough to drink, $T \leq 60\,\degC$? Inverting the
model:
\[
t_{60} = \frac{1}{k} \ln\!\left(\frac{T_{0} - T_{\infty}}{60 - T_{\infty}}\right)
       = \frac{1}{k} \ln\!\left(\frac{68}{38}\right)
       \approx \frac{0.58}{k}.
\]
With $\hat k = 0.084 \pm 0.005\,\si{\per\minute}$, the prediction
is $t_{60} \approx 6.9 \pm 0.4$ minutes (the $\pm$ from propagating
the interval on $k$ through the inverse relation). The prediction
interval is the engineer's report; a reader who said ``$7$
minutes'' without the interval has hidden the model's honest
uncertainty.

\subsection{The look-back}

Does the result use all the data? Yes; the regression used all
thirty points. Does it explain the small residual trend at high
temperature excess? Partially; the radiative correction is the
candidate next term. Does it admit an independent check? Yes; a
second cup, same conditions, can be measured and compared, and
the validation step did exactly that. Does the model survive
hostile review? Within its stated domain (cooling from below
$\sim 80\,\degC$ to room temperature, in still air, with the
thermometer in the water), yes; outside its domain, the model is
explicitly not applied.

The example is small. The discipline it demonstrates is general.
A measurement, a candidate model, an identification, a validation,
a structural-error account, a sensitivity propagation, a look-back.
Seven moves; the project at the close of the chapter asks the
reader to perform them on a quantity of their choice from the
Volume I notebook.

\section{Worked example, in full: a room's thermal response}\pathtag{core}

The cooling cup ran the cycle on a model that did not need revising. The
more instructive case is one where the first model fails its own
residual check and the cycle closes back on the abstraction. We work a
building thermal-response problem at full length, because the revision
step is where the discipline earns its place, and because the system is
one the reader can instrument with a thermometer and a clock in their
own home. The model reaches for the whole toolkit of the volume: the
ODEs of Chapter 7, the linear algebra of Chapters 9 and 10 to diagonalise
the coupled system, the least-squares fitting of Chapter 14, and the
numerical integration of Chapter 16.

\subsection{The question and the data}

A room is heated by an electric radiator on a thermostat. The
engineering question is operational: if the heater is switched on at a
known outdoor temperature, how long until the room reaches a comfortable
$20\,\degC$, and how much energy does that warm-up cost? The answer must
be good to about ten minutes in time and about ten percent in energy,
the precision a tenant would care about and a controls engineer would
need to size a pre-heat schedule.

The data are a temperature log. A thermometer logs the indoor air
temperature every two minutes for the first three hours after the heater
switches on, with the outdoor temperature logged in parallel and the
heater's electrical power read from a plug-in meter. The thermometer's
accuracy, from a Volume I Chapter 3 calibration, is $\pm 0.3\,\degC$;
the power meter reads to $\pm 2\%$; the clock is exact on this
timescale. A representative log starts at the outdoor temperature of
$6\,\degC$, rises slowly for the first twenty minutes, then climbs
through the comfort band and levels toward a steady value near
$21\,\degC$ as the thermostat begins to cycle.

\subsection{First candidate: one thermal lump}

The natural first abstraction treats the room as a single thermal
capacitance. The air and contents hold a quantity of heat measured by an
effective heat capacity $C$; heat leaks to the outdoors through an
effective thermal resistance $R$; the heater supplies power $P$. The
first law gives a single first-order equation,
\[
C \dv{T}{t} = P - \frac{T - T_{\text{out}}}{R},
\]
whose steady state is $T_\infty = T_{\text{out}} + PR$ and whose
approach is exponential with time constant $\tau = RC$:
\[
T(t) = T_\infty - (T_\infty - T_0)\, e^{-t/\tau}.
\]
Two parameters, $R$ and $C$, and a clean prediction. The abstraction
discipline of section 18.2 endorses the model as the simplest candidate:
one compartment, one resistance, the structure of the single-tank model
of section 18.5 applied to heat. The dominant scale is $\tau = RC$,
which the reader estimates before fitting at a few hours from the visible
levelling time of the log.

\subsection{Identification of the first model}

The model is nonlinear in $\tau$ but linear once $T_\infty$ is fixed.
With $T_\infty$ read from the late-time plateau at $21\,\degC$, the
transformed variable $\ln(T_\infty - T)$ is linear in $t$ with slope
$-1/\tau$. Linear regression on the log-excess returns
$\hat\tau \approx 52\,\text{min}$ with a regression standard error that
puts a 95\% interval of roughly $\pm 4\,\text{min}$ on $\tau$. From the
steady-state relation and the measured power $P = 1.5\,\text{kW}$, the
resistance is $\hat R = (T_\infty - T_{\text{out}})/P = 15/1500 =
0.010\,\si{\kelvin\per\watt}$ and the capacity is $\hat C = \hat\tau/\hat
R \approx 3120\,\text{s} / 0.010 \approx 3.1\times 10^{5}\,\si{\joule\per\kelvin}$,
which for a small room's air plus light furnishings is the right order
of magnitude. The identification has returned physical parameters, not
just curve-fit constants, and their magnitudes pass the
order-of-magnitude check.

\subsection{Validation, and the residual that breaks the model}

The validation step is where the first model fails honestly. The fit
residuals, plotted against time, are systematically
negative in the first half hour (the model predicts the room warms
faster than it does) and systematically positive afterward. The pattern
is the signature of a missing dynamic, and its sign points to the
mechanism: the model assumes the heat goes instantly into a single
capacitance, but the room's mass (walls, floor, furniture) must warm
before the air settles, and the wall mass lags the air. The single lump
cannot produce the S-shaped, delayed rise the data show; it can only
produce a prompt exponential.
Figure~\ref{fig:vol02:ch18:thermal-mass-revision}(a) shows the leading
fit and the structured residual.

The discipline of section 18.5 is explicit about the move here. The
residual structure is the next model's specification, not noise to be
covered by a wider interval. The room has at least two thermal
masses, the fast air and the slow walls, coupled by a resistance between
them. The revision is not free invention; it is the term the residuals
demanded.

% Figure: the model-revision arc for the building thermal-mass case study.
% (a) one-capacitance model: systematic residual lag. (b) two-capacitance
% model: residuals collapse to noise. Shows revision driven by structure.
\begin{figure}[htbp]
\centering
\begin{tikzpicture}[font=\small]
% ===== Panel (a): single-lump model =====
\begin{scope}[xshift=0cm]
\draw[->, thick, draw=engblack!70] (0,0) -- (5.4,0)
  node[right, font=\scriptsize] {$t$ (h)};
\draw[->, thick, draw=engblack!70] (0,0) -- (0,4.0)
  node[above, font=\scriptsize] {$T_{\text{in}}$ (\si{\degreeCelsius})};
\node[font=\small, anchor=south west] at (0.1,3.8) {\textbf{(a) one lump}};
% measured indoor temperature (data): rises with a delay after heater on
\foreach \px/\py in {0.3/0.5, 0.9/0.55, 1.5/0.75, 2.1/1.25, 2.7/2.0, 3.3/2.7, 3.9/3.1, 4.5/3.25} {
  \filldraw[engred!85] (\px,\py) circle (0.055);
}
% single-lump fit (no delay): rises too early
\draw[very thick, draw=engblue!80, smooth, samples=60, domain=0.2:4.8]
  plot (\x, { 3.3*(1 - exp(-0.7*\x)) });
\node[anchor=west, font=\scriptsize, color=engblack!70] at (1.2, 3.5)
  {fit leads data};
\end{scope}
% ===== Panel (b): two-lump model =====
\begin{scope}[xshift=7.0cm]
\draw[->, thick, draw=engblack!70] (0,0) -- (5.4,0)
  node[right, font=\scriptsize] {$t$ (h)};
\draw[->, thick, draw=engblack!70] (0,0) -- (0,4.0)
  node[above, font=\scriptsize] {$T_{\text{in}}$ (\si{\degreeCelsius})};
\node[font=\small, anchor=south west] at (0.1,3.8) {\textbf{(b) two lumps}};
\foreach \px/\py in {0.3/0.5, 0.9/0.55, 1.5/0.75, 2.1/1.25, 2.7/2.0, 3.3/2.7, 3.9/3.1, 4.5/3.25} {
  \filldraw[engred!85] (\px,\py) circle (0.055);
}
% two-lump fit (S-shaped, captures the lag): T = A(1 - (1+bt)e^{-bt}) like
\draw[very thick, draw=engblue!80, smooth, samples=80, domain=0.2:4.8]
  plot (\x, { 3.3*(1 - (1 + 0.95*\x)*exp(-0.95*\x)) });
\node[anchor=west, font=\scriptsize, color=engblack!70] at (1.5, 3.5)
  {fit tracks lag};
\end{scope}
\end{tikzpicture}
\caption[Model revision driven by residual structure: building
thermal mass]{The indoor-temperature response of a room after the heater
switches on, modelled two ways. (a) A single-lump model (one thermal
capacitance, one resistance) predicts an immediate exponential rise; the
data lag behind because the wall mass must warm before the air does, and
the residuals carry a systematic early-time deficit. The residual
structure is the model's own statement that a feature is missing. (b) A
two-lump model (air capacitance and wall capacitance in series) produces
the S-shaped response the data show, and the residuals collapse to
measurement noise. The revision was not free invention; it was the term
the panel-(a) residuals demanded.}
\label{fig:vol02:ch18:thermal-mass-revision}
\end{figure}


\subsection{Second candidate: two thermal lumps}

The two-capacitance model gives the air a capacity $C_a$ and the wall
mass a capacity $C_w$, couples them through an internal resistance
$R_{aw}$, and lets the air lose heat to the outdoors through $R_{ao}$.
The heater drives the air node. The equations are a coupled pair,
\[
C_a \dv{T_a}{t} = P - \frac{T_a - T_w}{R_{aw}} - \frac{T_a - T_{\text{out}}}{R_{ao}},
\qquad
C_w \dv{T_w}{t} = \frac{T_a - T_w}{R_{aw}},
\]
which is a linear system $\dot{\mathbf{T}} = \mathbf{A}\mathbf{T} +
\mathbf{b}$ of exactly the kind Volume II Chapter 7 solves and Chapter
10 diagonalises. The two eigenvalues of $\mathbf{A}$ are two time
constants, a fast one (the air responding to the heater) and a slow one
(the wall mass charging through $R_{aw}$). The general solution is a sum
of two exponentials,
\[
T_a(t) = T_\infty + c_1 e^{-t/\tau_1} + c_2 e^{-t/\tau_2},
\]
and the sum of a fast and a slow exponential is exactly the S-shaped
rise the single lump could not produce. The model now has four
parameters ($C_a, C_w, R_{aw}, R_{ao}$), one more than the question
strictly needs, which the abstraction discipline flags as a cost to be
justified by the residual evidence, and the evidence is the structured
residual of the first model.

\subsection{Identification of the second model}

Four parameters fit by nonlinear least squares on the numerically
integrated $T_a(t)$, with the integration done by the Runge-Kutta
method of Chapter 16. The fit returns the two time constants directly:
$\hat\tau_1 \approx 12\,\text{min}$ (air) and $\hat\tau_2 \approx
95\,\text{min}$ (walls), with the four physical parameters recovered
from the eigenstructure. The identifiability check matters here, because
four parameters fit to a single warm-up curve can be weakly
constrained. The Jacobian condition number at the fit is around $40$,
below the $10^3$ weak-identifiability threshold but well above the
single-lump model's near-unity value: the two time constants are
separated enough to be distinguished by data that span both, which this
log does. A reader who logged only the first twenty minutes would see
only the fast constant and would find $\tau_2$ practically
unidentifiable, the same lesson the logistic model taught about $K$.

\subsection{Validation of the second model, and the prediction}

The two-lump model's residuals collapse to the thermometer noise floor,
with no surviving time structure
(figure~\ref{fig:vol02:ch18:thermal-mass-revision}(b)). A held-out
second warm-up, logged on a different day with a different outdoor
temperature, is predicted to within the thermometer accuracy once
$T_{\text{out}}$ is updated; the model has generalised to data the
identification did not use. The cycle has closed.

The engineering answer follows. The time to reach $20\,\degC$ from a
$6\,\degC$ start is read from the fitted $T_a(t)$ as
$t_{20} \approx 41\,\text{min}$, against the single-lump model's
$t_{20} \approx 28\,\text{min}$, a thirteen-minute error that exceeds
the question's ten-minute tolerance and is exactly the error the
unrevised model would have shipped. The energy cost is the power times
the time the heater is actually on, which over the warm-up is close to
$P \cdot t_{20} = 1.5\,\text{kW} \times 0.68\,\text{h} \approx
1.0\,\text{kWh}$, plus the standing loss thereafter. The prediction
interval, propagating the parameter intervals through the integrated
solution and inflating for the residual structural error (negligible
here, the inflation factor is near one because the two-lump residuals
are white), is $t_{20} = 41 \pm 5\,\text{min}$ and energy $1.0 \pm
0.1\,\text{kWh}$. Both meet the question's precision.

\subsection{The look-back and the lesson}

Does the result use all the data? Yes; the fit used the full three-hour
log and the held-out day validated it. Does it explain what the first
model could not? Yes; the wall lag is the slow eigenvalue, and the slow
eigenvalue is the term the first model's residuals demanded. Does it
admit an independent check? Yes; the recovered $C_a$ and $C_w$ can be
estimated independently from the room's air volume and wall mass and
specific heats, and the estimates agree to within a factor that the
crude geometry justifies. Does the model survive hostile review within
its domain (a single room, electric resistive heating, outdoor
temperature roughly constant over the warm-up, no solar gain, no
occupancy heat)? Yes, and the domain statement names the regimes (solar
gain on a sunny wall, a second occupant, a cold front during the warm-up)
where the model would need a third input.

The lesson is the one the chapter exists to install. The first model was
not wrong because the modeller was careless; it was the correct simplest
first candidate, and it failed in the specific, readable way that told
the modeller exactly what to add. The discipline is not to get the model
right on the first try. The discipline is to read the residual, to let it
specify the revision, and to validate the revision against data the
revision was not fit to. The two-lump model is not the last word either;
a reader heating from a cold start in winter sun would find its
residuals carrying a daily structure that demands a solar-gain term. The
cycle never closes permanently. It closes well enough for the question,
which is the only standard a model is ever held to.

\section{Worked example, in full: staffing a service under random demand}\pathtag{standard}

The two thermal examples ran the cycle on deterministic models, where the
randomness lived only in the measurement noise and the model itself was a
differential equation with a definite solution. A large class of
engineering questions has randomness in the system, not merely in the
instrument: customers arrive at a clinic at unpredictable instants, jobs
hit a server in bursts, parts reach an inspection station at a rate that
fluctuates around its mean. The quantity the engineer must predict, the
time a customer waits, is itself a random variable, and the model that
predicts it is a stochastic model. We work a staffing problem at full
length because it exercises a modelling archetype the chapter has not yet
touched, the discrete-event queue, and because it closes the loop in a way
the deterministic examples could not: the analytic model is validated
against a Monte-Carlo simulation of the same random process, two
independent computations of the same quantity that must agree or expose an
error. The example reaches across the volume: the probability of Chapter
14 supplies the arrival and service distributions, the linear algebra of
Chapter 10 solves for the steady state of the underlying Markov chain, the
optimisation of Chapter 11 finds the staffing level that minimises total
cost, and the numerical methods of Chapter 16 drive both the
transcendental Erlang formula and the simulation.

\subsection{The question and the data}

A walk-in clinic staffs a number of identical service desks. Patients
arrive without appointments; each desk serves one patient at a time; a
patient who finds every desk busy waits in a single shared line. The
engineering question is operational and has a cost attached: how many
desks $c$ should the clinic staff so that the expected wait stays below a
service target of five minutes, and what staffing level minimises the sum
of labour cost and the cost of patients' waiting time? The answer must be
an integer (a desk is staffed or it is not) and must be defensible to a
clinic manager who pays the wages and answers for the complaints.

The data are an arrival log and a service log. Over a representative busy
period, the clinic records the instant each patient arrives and the
duration of each service. A morning's log gives $480$ arrivals over four
hours, a mean arrival rate of $\lambda = 2$ patients per minute, and a
mean service duration of $\bar s = 1.5$ minutes per patient, so a single
desk serves at rate $\mu = 1/\bar s \approx 0.667$ patients per minute.
The logs also record the full distribution of inter-arrival gaps and of
service durations, because the distributions, not just their means, are
the modelling content.

\subsection{First candidate: the M/M/c queue}

The abstraction discipline of section 18.2 names the question first
(expected wait as a function of staffing) and then the dominant scales.
The dimensionless group that governs every queue is the
\engterm{utilisation} $\rho = \lambda/(c\mu)$, the fraction of the total
service capacity the arrivals consume. A system with $\rho \ge 1$ has
arrivals faster than the desks can clear them and its queue grows without
bound; a system with $\rho < 1$ reaches a statistical steady state. The
group is the queue's Reynolds number, the single ratio whose value
decides which regime the system occupies, and the engineer estimates it
before any equation: with $\lambda = 2$ and $\mu = 0.667$, a single desk
gives $\rho = 3$, hopelessly overloaded, and stability requires at least
$c = 4$ desks, where $\rho = 2/(4 \times 0.667) = 0.75$.

The simplest model that answers the question treats the arrivals as a
Poisson process (independent arrivals at constant average rate) and the
service durations as exponentially distributed (memoryless service), with
$c$ identical desks. This is the M/M/c queue, the two M's naming the
Markov (memoryless) arrival and service processes. The memorylessness is
the load-bearing assumption, and it is exactly what makes the model
tractable: the number of patients in the system is a continuous-time
Markov chain, the birth-death chain of
figure~\ref{fig:vol02:ch18:queue-state}(a), whose steady-state occupancy
probabilities $p_n$ satisfy the cut-balance relations $\lambda p_{n-1} =
\min(n, c)\,\mu\, p_n$. Solving the chain, a linear-algebra computation of
exactly the kind Chapter 10 performs on any transition operator, gives the
Erlang C formula for the probability that an arriving patient must wait,
\[
P_{\text{wait}}
  = \frac{\dfrac{(c\rho)^c}{c!}\dfrac{1}{1-\rho}}
         {\displaystyle\sum_{n=0}^{c-1}\frac{(c\rho)^n}{n!}
          + \frac{(c\rho)^c}{c!}\frac{1}{1-\rho}},
\]
and the mean wait in queue follows from it as
\[
W_q = \frac{P_{\text{wait}}}{c\mu - \lambda}.
\]

% Figure: the birth-death state diagram of an M/M/c queue and the
% nonlinear blow-up of mean wait as utilisation approaches one.
\begin{figure}[htbp]
\centering
\begin{tikzpicture}[font=\small,
  state/.style={circle, draw=engblack!75, fill=engblue!12, minimum size=0.8cm,
    font=\small},
  rate/.style={->, thick, draw=engblack!70, font=\scriptsize}]
% ===== Top: birth-death chain for a single server (c=1) =====
\begin{scope}[yshift=0cm]
\foreach \n/\lab in {0/0, 1/1, 2/2, 3/3} {
  \node[state] (q\n) at ({2.0*\n},0) {\lab};
}
\node (qdots) at (8.0,0) {$\cdots$};
% arrivals (lambda) rightward, top arcs
\foreach \a/\b in {q0/q1, q1/q2, q2/q3} {
  \draw[rate] (\a) to[bend left=35] node[above] {$\lambda$} (\b);
}
\draw[rate] (q3) to[bend left=35] node[above] {$\lambda$} (qdots);
% departures (mu) leftward, bottom arcs
\foreach \a/\b in {q1/q0, q2/q1, q3/q2} {
  \draw[rate] (\a) to[bend left=35] node[below] {$\mu$} (\b);
}
\draw[rate] (qdots) to[bend left=35] node[below] {$\mu$} (q3);
\node[font=\small, anchor=south west] at (-1.0, 0.7)
  {\textbf{(a) birth-death chain, single server}};
\node[font=\scriptsize, anchor=north, text width=8.5cm, align=left]
  at (3.5,-0.9) {State $n$ is the number in system. Balance of probability
  flux across each cut gives $\lambda p_{n-1} = \mu p_n$, hence
  $p_n = (1-\rho)\rho^n$ with $\rho = \lambda/\mu$.};
\end{scope}
% ===== Bottom: mean wait vs utilisation =====
\begin{scope}[yshift=-5.6cm]
\draw[->, thick, draw=engblack!70] (0,0) -- (8.5,0)
  node[right, font=\scriptsize] {utilisation $\rho$};
\draw[->, thick, draw=engblack!70] (0,0) -- (0,3.8)
  node[above, font=\scriptsize] {mean wait $W_q$ (scaled)};
\foreach \xt/\xv in {0/0, 2/0.25, 4/0.5, 6/0.75, 7.6/0.95} {
  \draw[draw=engblack!70] (\xt,0) -- (\xt,-0.07);
  \node[font=\scriptsize, anchor=north] at (\xt,-0.09) {\xv};
}
% W_q proportional to rho/(1-rho); scaled so curve fits the box
\draw[very thick, draw=engred!85, smooth, samples=120, domain=0:0.93]
  plot ({8.0*\x}, { min(3.6, 0.18*\x/(1-\x)) });
\draw[dashed, draw=engblack!45] (8.0,0) -- (8.0,3.8);
\node[font=\scriptsize, anchor=south east, text width=3.3cm, align=right,
  draw=engblack!20, fill=engblue!6, inner sep=3pt] at (7.7,2.6)
  {$W_q \propto \dfrac{\rho}{1-\rho}$: doubling staff to cut $\rho$ from
  $0.9$ to $0.45$ cuts the wait by roughly an order of magnitude.};
\node[font=\small, anchor=south west] at (0.2, 3.5)
  {\textbf{(b) the congestion nonlinearity}};
\end{scope}
\end{tikzpicture}
\caption[The M/M/c queue: state chain and the congestion
nonlinearity]{(a) The birth-death Markov chain for a single-server queue.
Arrivals raise the state at rate $\lambda$, service lowers it at rate
$\mu$, and the steady-state occupancy probabilities $p_n$ follow from the
cut balance, the same linear-algebra steady-state computation Volume II
Chapter 10 performs on any Markov transition operator. (b) The mean queue
wait $W_q$ as a function of utilisation $\rho = \lambda/(c\mu)$. The wait
is mild while $\rho$ is below about $0.7$ and diverges as $\rho \to 1$.
The nonlinearity is the entire reason a staffing model is worth building:
a system run at $95\%$ utilisation to save labour cost incurs waits an
order of magnitude longer than the same system at $70\%$, and a planner
who reasons about the mean load alone, without the queue model, will size
the service one server short of the point where the wait explodes.}
\label{fig:vol02:ch18:queue-state}
\end{figure}


Two structural facts are available before any data are fit, and both come
from the dimensionless group rather than from numerical work. First,
\engterm{Little's law}, $L = \lambda W$, relates the mean number in system
$L$ to the mean time in system $W$ through the arrival rate alone; it holds
for any stable queue regardless of the arrival or service distribution,
and it gives the engineer a conservation-law check on any predicted pair
$(L, W)$. Second, the congestion nonlinearity of
figure~\ref{fig:vol02:ch18:queue-state}(b): $W_q$ grows as
$\rho/(1-\rho)$, mild while $\rho$ is below about $0.7$ and divergent as
$\rho \to 1$. The nonlinearity is why the model is worth building. A
planner who reasons about the mean load alone would note that four desks
have spare capacity ($\rho = 0.75$) and stop; the queue model shows that
$\rho = 0.75$ already sits on the shoulder of the blow-up, and the wait is
sensitive to one desk more or less in a way no mean-load argument reveals.

\subsection{Identification of the model}

The M/M/c model has no free parameters to fit once $\lambda$, $\mu$, and
$c$ are known: $\lambda$ and $\mu$ are estimated directly from the logs as
sample means, and $c$ is the decision variable. The identification work
is therefore the identification of the rate parameters and, more
importantly, the validation of the two distributional assumptions the
model rests on. The mean arrival rate is $\hat\lambda = 2.0$ per minute
with a Poisson sampling interval of roughly $\pm 0.09$ from the
$480$-arrival count ($\hat\lambda \pm \sqrt{\hat\lambda/T}$). The mean
service rate is $\hat\mu = 0.667$ per minute with an interval set by the
service-duration sample.

The distributional check is the part a careless modeller skips. For the
arrival process to be Poisson, the inter-arrival gaps must be exponential,
which the engineer checks by plotting the empirical gap distribution
against the exponential on a log scale (an exponential is a straight line
on a log-survival plot) or by a Kolmogorov-Smirnov test of the kind
Chapter 14 supplies. For the service to be exponential, the same check
applies to the service durations. The clinic's arrival gaps pass the
check well, as walk-in arrivals from a large independent population
genuinely are close to Poisson. The service durations do not: their
histogram is far less dispersed than an exponential, with a coefficient of
variation near $0.5$ rather than the exponential's $1.0$, because a
consultation has a typical length and rarely runs to many multiples of it.
The residual that the deterministic examples found in a time-series plot,
the queue model finds in a distribution-shape mismatch, and it is the same
kind of evidence: the data refuse the model's assumption, and the refusal
specifies the revision.

\subsection{Validation against a Monte-Carlo simulation}

Before revising, the engineer validates the M/M/c model as it stands,
because a model that cannot reproduce its own assumptions' consequences
has a coding error, not a modelling error, and the two must be separated.
The validation tool for a stochastic model is a \engterm{discrete-event
simulation}: a program that draws inter-arrival gaps and service durations
from the assumed distributions, advances an event clock from one arrival
or departure to the next, and records the wait each simulated patient
experiences. Averaging the waits over a long run, and over many replicated
runs to quantify the sampling spread, gives a Monte-Carlo estimate of
$W_q$ that depends on no closed-form formula. The simulation and the
Erlang C formula are two independent computations of the same quantity:
if they agree, both are trustworthy in their shared domain; if they
disagree, one has an error to find.

Figure~\ref{fig:vol02:ch18:queue-validation}(a) shows the agreement. The
analytic $W_q(c)$ curve falls within the simulation's standard-error bars
at every staffing level from $c = 3$ to $c = 7$, across $40$ replicated
runs of $10{,}000$ patients each. The validation is the cross-check the
deterministic examples got from held-out data: here the held-out evidence
is a second, structurally independent route to the same number. The
agreement also calibrates the simulation itself, so that when the engineer
later runs the simulation with the non-exponential service the formula
cannot handle, the simulation is already trusted.

% Figure: analytic M/M/c prediction against a discrete-event Monte-Carlo
% simulation, with the staffing-cost optimisation curve alongside.
\begin{figure}[htbp]
\centering
\begin{tikzpicture}[font=\small]
% ===== Left: analytic vs simulation, mean wait vs number of servers =====
\begin{scope}[xshift=0cm]
\draw[->, thick, draw=engblack!70] (0,0) -- (6.0,0)
  node[right, font=\scriptsize] {servers $c$};
\draw[->, thick, draw=engblack!70] (0,0) -- (0,4.0)
  node[above, font=\scriptsize] {mean wait $W_q$ (min)};
\foreach \xt/\xv in {1/3, 2/4, 3/5, 4/6, 5/7} {
  \draw[draw=engblack!70] (\xt,0) -- (\xt,-0.07);
  \node[font=\scriptsize, anchor=north] at (\xt,-0.09) {\xv};
}
\foreach \yt/\yv in {0/0, 1/10, 2/20, 3/30} {
  \draw[draw=engblack!70] (0,\yt) -- (-0.07,\yt);
  \node[font=\scriptsize, anchor=east] at (-0.09,\yt) {\yv};
}
% Analytic Erlang-C wait (decreasing, convex): tabulated points as a curve
\draw[very thick, draw=engblue!80] plot[smooth]
  coordinates {(1,3.6) (2,1.45) (3,0.62) (4,0.27) (5,0.11)};
% Simulation points with error bars (close to analytic)
\foreach \x/\y/\e in {1/3.45/0.30, 2/1.52/0.18, 3/0.58/0.10, 4/0.30/0.06, 5/0.10/0.04} {
  \draw[draw=engred!85, thick] (\x,{\y-\e}) -- (\x,{\y+\e});
  \fill[engred!85] (\x,\y) circle (1.6pt);
}
\node[font=\small, anchor=south west] at (0.1, 3.7)
  {\textbf{(a) prediction vs simulation}};
\draw[very thick, draw=engblue!80] (3.6,3.2) -- (4.0,3.2);
\node[anchor=west, font=\scriptsize] at (4.05,3.2) {Erlang C};
\fill[engred!85] (3.8,2.85) circle (1.6pt);
\node[anchor=west, font=\scriptsize] at (4.05,2.85) {sim.\ $\pm$SE};
\end{scope}
% ===== Right: total cost vs number of servers (the optimisation) =====
\begin{scope}[xshift=8.2cm]
\draw[->, thick, draw=engblack!70] (0,0) -- (6.0,0)
  node[right, font=\scriptsize] {servers $c$};
\draw[->, thick, draw=engblack!70] (0,0) -- (0,4.0)
  node[above, font=\scriptsize] {cost per hour (scaled)};
\foreach \xt/\xv in {1/3, 2/4, 3/5, 4/6, 5/7} {
  \draw[draw=engblack!70] (\xt,0) -- (\xt,-0.07);
  \node[font=\scriptsize, anchor=north] at (\xt,-0.09) {\xv};
}
% staffing cost: linear rising
\draw[thick, draw=enggray, dashed] plot coordinates
  {(1,0.8) (2,1.2) (3,1.6) (4,2.0) (5,2.4)};
\node[anchor=west, font=\scriptsize, text=enggray] at (4.6,2.5) {staff};
% waiting cost: convex falling (proportional to W_q * lambda)
\draw[thick, draw=engamber, dash dot] plot[smooth] coordinates
  {(1,3.6) (2,1.7) (3,0.85) (4,0.45) (5,0.25)};
\node[anchor=west, font=\scriptsize, text=engamber] at (4.7,0.35) {wait};
% total cost: U-shape with minimum near c=4
\draw[very thick, draw=engblue!85] plot[smooth] coordinates
  {(1,4.4) (2,2.9) (3,2.45) (4,2.45) (5,2.65)};
\node[anchor=west, font=\scriptsize, text=engblue!85] at (1.0,4.3) {total};
\draw[dashed, draw=engred!70] (4,0) -- (4,2.45);
\fill[engred!85] (4,2.45) circle (2.0pt);
\node[font=\scriptsize, anchor=south, text=engred!85] at (4,2.5) {$c^\star$};
\node[font=\small, anchor=south west] at (0.1, 3.7)
  {\textbf{(b) the staffing optimum}};
\end{scope}
\end{tikzpicture}
\caption[Queue model: validation against simulation and the staffing
optimum]{(a) The Erlang-C analytic prediction of mean wait against the
number of servers (solid line) compared with a discrete-event Monte-Carlo
simulation of the same arrival and service process (points, with one
standard error from $40$ replicated runs). The analytic curve falls
within the simulation's sampling interval at every staffing level, which
is the validation: two independent computations of the same quantity, one
closed-form and one stochastic, agree to within the simulation's noise.
(b) The optimisation the model exists to support. Staffing cost rises
linearly with $c$; expected waiting cost (lost time priced per customer
per minute) falls convexly as congestion eases. Their sum is U-shaped,
and its minimum $c^\star$ is the staffing level the model recommends. The
optimum sits at the staffing level where the marginal labour cost of one
more server equals the marginal waiting cost it removes, the discrete
analogue of the first-order condition Volume II Chapter 11 derives for
smooth objectives.}
\label{fig:vol02:ch18:queue-validation}
\end{figure}


\subsection{Revision: the service distribution the data demanded}

The distribution check refused the exponential service. The revision is
the M/G/c queue, where the G (general) service distribution carries the
measured coefficient of variation $c_s \approx 0.5$ rather than the
exponential's $1.0$. No simple closed form gives the exact wait for M/G/c,
but the Allen-Cunneen approximation corrects the Erlang C wait by the
variability factor $(c_a^2 + c_s^2)/2$, where $c_a$ is the arrival-gap
coefficient of variation (near $1$ for the Poisson arrivals) and $c_s$ the
service coefficient of variation. With $c_s = 0.5$ the factor is
$(1 + 0.25)/2 = 0.625$, so the M/M/c wait overstates the true wait by
about a third: the exponential model, by assuming service is more variable
than it is, predicts more congestion than the clinic actually suffers. The
discrete-event simulation, rerun with service durations drawn from the
measured distribution rather than an exponential, confirms the correction
directly, and the two revised routes again agree.

The revision changes the engineering answer. At $c = 4$ desks, the M/M/c
model predicts $W_q \approx 1.4$ minutes; the M/G/c revision predicts
$W_q \approx 0.9$ minutes. Both clear the five-minute target, so for the
target question the simpler model would not have misled. The revision
matters for the cost optimisation, where a third less waiting time shifts
the balance between labour cost and waiting cost and can move the optimal
staffing by a full desk.

\subsection{The optimisation, and the prediction}

The staffing decision is an optimisation, the discrete analogue of the
first-order conditions Chapter 11 derives for smooth objectives. The total
hourly cost is the sum of a labour cost, $c \cdot w$ for wage rate $w$,
and a waiting cost, $\lambda \cdot W_q(c) \cdot \kappa$, where $\kappa$
prices a patient-minute of waiting. Labour cost rises linearly in $c$;
waiting cost falls convexly as the congestion nonlinearity eases. Their
sum is U-shaped, with an interior minimum at the staffing level $c^\star$
where the marginal labour cost of one more desk equals the marginal
waiting cost it removes
(figure~\ref{fig:vol02:ch18:queue-validation}(b)). Because $c$ is an
integer, the optimisation is a one-dimensional search over a handful of
values rather than a calculus problem, and the engineer evaluates the
total cost at $c = 4, 5, 6$ and reads off the minimum. With a wage of
\$30 per hour and a waiting cost of \$1 per patient-minute, the M/G/c
model places $c^\star = 4$ at the morning peak, with $c = 5$ a close and
more conservative second that buys headroom against demand surges.

The prediction the engineer reports carries an interval, and the interval
has two sources the chapter has named throughout. The parameter
uncertainty propagates the sampling intervals on $\hat\lambda$ and
$\hat\mu$ through the Erlang formula. The structural uncertainty is larger
and is the honest part of the report: the model assumes a stationary
arrival rate, but a real clinic has a morning peak and a midday lull, so a
single $\rho$ understates the worst-hour congestion. The discipline of
section 18.5 applies. The engineer reports the staffing recommendation
peak-hour by peak-hour rather than as a single daily number, states that
the stationary model is validated only within a peak-load window, and
notes the regime, a demand surge or a desk going offline, in which the
recommendation would need re-running.

\subsection{The look-back and the lesson}

Does the result use all the data? Yes; the rate parameters came from the
full log and the distribution shapes from the full samples. Does it
explain what the first model could not? Yes; the service-distribution
mismatch was the residual, and the M/G/c variability factor is the term
the residual demanded. Does it admit an independent check? Yes, twice: the
discrete-event simulation validated the analytic formula, and Little's law
gives a conservation check on the predicted $(L, W)$ pair. Does the model
survive hostile review within its domain (a single peak-load window,
stationary arrivals, a shared queue, identical desks, no patient
abandonment)? Yes, and the domain statement names the regimes (time-varying
demand, patients who leave when the line is long, priority triage) where a
richer model would be needed.

The lesson generalises the deterministic examples to the stochastic case.
The randomness in a queue lives in the system, so the model is a
distribution rather than a number, and the validation tool is a second
independent computation, the Monte-Carlo simulation, rather than a
held-out data set. The discipline is unchanged: name the question, choose
the simplest model, check its assumptions against the data, validate it
against an independent route, read the residual (here the
distribution-shape mismatch), let it specify the revision, and report the
prediction with an interval that exposes the structural uncertainty the
stationary assumption hides.

\begin{mastery}
The discrete-event simulation and the Erlang formula are an instance of a
distinction the chapter has used implicitly and is worth naming:
\engterm{verification} versus \engterm{validation}. Verification asks
whether the model is solved correctly, that the code computes what the
equations say; validation asks whether the equations describe the world.
The simulation-versus-formula agreement of
figure~\ref{fig:vol02:ch18:queue-validation}(a) is a verification: both
compute the M/M/c wait, so their agreement says the implementations are
correct, not that the M/M/c model fits the clinic. The
distribution-shape check against the service log is the validation: it
compares the model's assumption to the world and finds it wanting. A
modeller who confuses the two reports a verified model as a validated one,
the most common form of misplaced confidence in computational work. The
rule is mechanical: agreement between two computations of the same model
verifies; agreement between the model and held-out data from the system
validates; only the second licenses a prediction about the world. The
aerospace and nuclear communities, where a wrong simulation costs lives,
codified the separation into formal verification-and-validation standards,
and the engineering reader takes from them the habit of labelling every
cross-check as one or the other before trusting it.
\end{mastery}

\section{Worked example: a biological system, a financial system, a structural system}\pathtag{standard}

The seven moves of the previous example apply to systems whose
mathematics is more interesting than Newton's cooling. We sketch
three at engineering depth, with enough algebra to show the
moves and not so much that the section becomes a survey.

\subsection{A biological system: logistic growth}

A bacterial culture in a flask grows by cell division. The cells
are exposed to a finite resource (glucose, dissolved oxygen,
space). The candidate model is the logistic equation
\[
\dv{N}{t} = r N \left(1 - \frac{N}{K}\right),
\]
where $N(t)$ is the cell count at time $t$, $r$ is the
intrinsic growth rate, and $K$ is the carrying capacity. The
analytical solution is
\[
N(t) = \frac{K}{1 + \left(K/N_{0} - 1\right) e^{-rt}}.
\]

The reader measures $N(t_{i})$ at hourly intervals for a 24-hour
culture. Identification fits $r$ and $K$ by nonlinear least
squares. Validation holds out the last six hours of data and
checks the fit's prediction; the residuals are inspected for
trend.

The structural error is consequential. The model assumes
unstructured well-mixed exponential growth that saturates at
$K$. Real cultures show a lag phase at the start (cells
adjusting to the medium), a stationary phase at the end
(cells maintaining count without dividing), and a death phase
beyond. The logistic model captures only the middle. A reader
who applies the logistic model to data extending into the
death phase will fit a value of $K$ that is biased downward
and a value of $r$ that is biased downward as well; the fit
will look good, the predictions outside the data range will
be wrong. The discipline is to state the time window in which
the model applies and to refuse the fit on data outside that
window.

The engineering use is real. A bioreactor designer who needs
to know how long the culture takes to reach a productive
density, or how much substrate to supply, uses the logistic
model with stated $r$ and $K$ intervals to size the reactor.
The model's structural error sets the safety margin: a
reactor sized at the model's nominal time, without margin
for the lag phase, will be ready later than the engineer
expected. The engineer adds the lag phase as a fixed offset
or moves to a structured model with a lag term.

\subsection{A financial system: a discounted cash flow}

A small business produces an annual cash flow $C_{i}$ for
$i = 1, \ldots, n$ years and is sold at year $n$ for residual
value $R$. The candidate model is the discounted cash flow:
the present value is
\[
PV = \sum_{i=1}^{n} \frac{C_{i}}{(1 + r)^{i}}
   + \frac{R}{(1 + r)^{n}},
\]
where $r$ is the discount rate, taken to reflect the
opportunity cost of capital and the riskiness of the cash
flows. The model is everywhere in finance and infrastructure
analysis; it is a working object on which acquisition,
investment, and infrastructure decisions are made.

Identification of the parameters is partly empirical and
partly judgement. The cash flows can be projected from
operating data; the discount rate is chosen by reference to
comparable instruments, the firm's cost of capital, or
regulatory rules. The residual value is judgemental.

The structural error dominates the report. The cash flows
are projections, not measurements; their interval is
typically wide and asymmetric, with a long upside and a
bounded downside (or the reverse). The discount rate
embodies risk; the choice of $r$ is itself a model. The
residual value at year $n$ depends on conditions $n$ years
from now, which the model cannot predict. A reader who
reports a single $PV$ value without intervals has produced a
number that will mislead any user who treats it as a
calibrated prediction.

The sensitivity apparatus of section 18.7 makes the dominance concrete.
Take a flat projected cash flow $C = \$100\text{k}$ per year for
$n = 10$ years with no residual value, and a discount rate $r = 8\%$. The
present value is the annuity sum $PV = C \sum_{i=1}^{10}(1+r)^{-i} =
100\text{k} \times 6.71 = \$671\text{k}$. The relative sensitivity to
the cash flow is $\tilde S_C = +1$ (the model is linear in $C$). The
relative sensitivity to the discount rate, computed from
$\partial PV/\partial r$, is only $\tilde S_r \approx -0.36$ at these
values: a one-percent relative change in $r$ (from $8.00\%$ to
$8.08\%$) moves the present value by about a third of a percent. By
sensitivity coefficient alone the cash flow is the stronger lever; what
makes the discount rate the one the report must expose is its far
larger uncertainty, whole percentage points against the few percent to
which a cash-flow projection anchored to operating data is known. A
modeller who pins $C$ to three significant figures while leaving $r$ as
an unexamined $8\%$ has controlled the input already well known and left
unquantified the one that is not. The
triangulation discipline applies directly: compute the present value at
$r = 6\%$, $8\%$, and $10\%$, report the resulting band of
$\$736\text{k}$ to $\$613\text{k}$, and let the decision turn on the
band.

The engineering discipline of this volume applies as it
does anywhere else. State the assumption set. Compute
intervals. Triangulate against alternative discount rates
and alternative growth rates. Report the range. The model is
only useful when its uncertainty is exposed, because the
decision the model supports (buy or do not buy, build or do
not build) turns on the interval rather than on the point.
A financial model that hides its interval is the same kind
of failure as an engineering model that hides its
structural error: the failure mode that section 18.8 names.

\subsection{A structural system: a beam under transverse load}

A simply supported beam of length $L$, cross-section moment of
inertia $I$, and Young's modulus $E$ carries a transverse load
$w$ per unit length. The Euler-Bernoulli beam equation gives
the deflection $y(x)$ from
\[
EI \dv[2]{y}{x} = M(x),
\]
where $M(x)$ is the bending moment. For uniform load on a
simply supported beam, the maximum deflection at midspan is
\[
y_{\max} = \frac{5}{384} \frac{w L^{4}}{E I}.
\]
The model is a working tool of structural engineering. It
has identifiable parameters: $L$, $I$, $E$, $w$, all measurable
or specifiable. It has a clean prediction: $y_{\max}$. It has
a verifiable check: load the beam, measure the deflection,
compare.

The model's structural error has known shape. The
Euler-Bernoulli theory assumes plane sections remain plane
(no shear deformation), small deflections (no geometric
nonlinearity), and linear elastic material (no yielding).
Each assumption fails at a known regime: shear matters for
short beams (span-to-depth ratio below about ten), geometric
nonlinearity matters at large deflection (relative to span,
above a few percent), yielding matters at high stress (above
the material's proportional limit). A beam designer who works
within the assumptions has a model whose predictions are
within a few percent of measured deflections; a designer who
works outside has a model whose predictions are wrong by an
amount the designer can estimate from the assumption that
fails.

The engineering discipline is to state the regime. A
data sheet reporting $y_{\max} = 5wL^{4}/(384 EI)$ for a
column of slender simply supported beams works within the
assumptions of the formula; the same formula applied to a
short stocky beam produces a deflection prediction that
under-estimates the true deflection by a factor that is
itself estimable from the shear-deformation correction.
The model's domain of validity is part of the model's
report.

The three sketched systems are different in their
quantities, in their data sources, and in their engineering
consequences. They are the same in their modelling
discipline. Identify the question, name the candidate
model, identify its parameters, validate against held-out
data, state the structural error, propagate the
sensitivity, and report the interval. The discipline is the
same; the apparatus changes.

\begin{mastery}
The three systems differ in one respect the chapter has so far kept in
the background: the source of their dominant uncertainty. The beam's
uncertainty is parametric and small; its structural form
(Euler-Bernoulli) is trustworthy within the slender-beam domain, and
the prediction interval is set by the material and dimensional
tolerances, which the tornado diagram of section 18.7 ranks. The
bacterial culture's uncertainty is structural and moderate; the
logistic form is an approximation that the lag and death phases violate,
and the honest interval is set by how far the data extend into the
regimes the logistic form omits. The discounted cash flow's uncertainty
is structural and dominant; the parametric uncertainty on a single
year's projected flow is dwarfed by the uncertainty in the discount
rate, the growth trajectory, and the residual value, none of which is a
measurement. The three thus span the spectrum from a model whose report
reduces to its parameter interval to a model whose parameter
interval is almost irrelevant beside its structural uncertainty. The
modeller's first act of judgement, before any fitting, is to place the
problem on this spectrum, because the placement decides where the
modelling effort should go: into tighter measurement for the beam, into
extending the data window for the culture, into scenario triangulation
for the cash flow. A modeller who spends a beam-style measurement effort
on a cash-flow problem has optimised the negligible term.
\end{mastery}

\section{Worked example, in full: contagion on a contact network}\pathtag{standard}

The SIR model of section~18.6 and the queue of section~18.9 both reduced a
population to an average. SIR assumed every individual mixes with every
other at the same rate; the queue assumed every patient draws from the same
arrival process. A large and growing class of engineering questions cannot
make that assumption, because the structure of who contacts whom, which
component depends on which, which node routes to which, is the thing that
governs the answer. The contact network of a disease, the dependency graph
of a power grid, the citation graph of a literature, the follower graph of
a platform: in each, a quantity propagates along edges whose pattern is the
model's content. We work a contagion-on-a-network problem at full length
because it exercises a modelling archetype the chapter has not yet touched,
the agent-based or network model, in which the system is a collection of
interacting individuals rather than a single aggregate equation. The
example closes the loop in a way the aggregate models could not: the
network model is validated against the mean-field SIR it generalises, and
the disagreement between them is itself the structural information the
heterogeneity carries. The example reaches across the volume: the
probability of Chapter 14 supplies the degree distribution and the
infection draws, the linear algebra of Chapter 10 supplies the adjacency
matrix and its spectral radius, the graph language of Chapter 17 supplies
the network itself, and the numerical methods of Chapter 16 drive the
individual-level simulation.

\subsection{The question and the data}

A workplace of $N$ people wants to bound how far an infection will spread
from a single index case before it burns out, and whether closing the
canteen, the one place everyone passes through, would change the answer.
The engineering question has two parts with a decision attached: what
fraction of the workforce does the model predict will be infected, and does
removing the highest-degree contact point (the hub) move that fraction
enough to justify the disruption. The answer must be good to about ten
percent in the infected fraction, the precision a facilities manager would
act on.

The data are a contact log. Over a representative week, a proximity sensor
records, for each pair of employees, whether they were in sustained close
contact. The log yields a contact network: a node per employee, an edge
per recorded contact pair. The network is not a number; it is the data, and
its structure, the distribution of how many contacts each person has, is
the modelling content. The log shows a right-skewed degree distribution: most
employees contact five to ten others, a handful (reception staff, the
canteen server) contact dozens. The mean degree is $\bar{d} \approx 9$,
but the variance is large, and the chapter's whole lesson about this
example turns on that variance.

\subsection{First candidate: the mean-field SIR (already in hand)}

The abstraction discipline names the question first (infected fraction and
its sensitivity to the hub) and the dominant scale next. The dimensionless
group that governs every contagion is the basic reproduction number $R_0$,
the expected number of secondary cases one infectious individual produces
in a fully susceptible population. The simplest model that answers the
question is the one the chapter already built: the homogeneous-mixing SIR
of section~18.6, with $R_0 = \beta/\gamma$ estimated from the contact rate
and the recovery time. With a per-contact transmission probability and the
mean degree $\bar{d} = 9$, the mean-field estimate is $R_0 \approx 1.8$,
comfortably above the threshold, and the final-size relation of
section~18.6's mastery box predicts an infected fraction near $0.73$.

The model is the correct first candidate and it answers the first half of
the question. It cannot answer the second half. The mean-field SIR has no
representation of the hub; it knows only the average degree, and removing
one high-degree node changes the average by less than one part in $N$. A
planner who stopped here would conclude, wrongly, that closing the canteen
is irrelevant. The residual that the deterministic examples found in a
time-series plot and the queue found in a distribution shape, this example
finds in a question the aggregate model structurally cannot address. The
refusal specifies the revision: the model must carry the network structure
the mean field threw away.

\subsection{Second candidate: the network model}

The revision keeps every individual as an agent and the contact log as the
coupling. Each node $i$ is in one of the three states $S$, $I$, $R$. An
infectious node infects each susceptible neighbour independently at rate
$\beta$ per edge; each infectious node recovers at rate $\gamma$. The model
is no longer a small system of ordinary differential equations; it is a
stochastic process on the graph, $N$ coupled state variables whose coupling
is the adjacency matrix $\mathbf{A}$ of Chapter 10, with $A_{ij} = 1$ when
$i$ and $j$ are in contact. The mean field is the special case in which
every $A_{ij}$ equals $\bar{d}/N$, the graph with all structure averaged
away.

The structural insight is again dimensional, and it is the payoff of
carrying the network. On a heterogeneous network the threshold is governed
not by the mean degree but by the ratio of the second moment of the degree
distribution to the first,
\[
R_0^{\text{net}} \;\approx\; \frac{\beta}{\gamma}\,
  \frac{\langle d^2\rangle - \langle d\rangle}{\langle d\rangle},
\]
a result that follows from tracking infection along edges rather than
across the population average. The hubs inflate $\langle d^2\rangle$ far
above $\langle d\rangle^2$, so a network with the same mean degree as a
homogeneous one can have a much larger effective reproduction number, and a
network whose hubs are removed can drop below threshold while its mean
degree barely moves. Albert, Jeong, and Barab{\'a}si established the
companion result for network robustness
\cite{paper:v6c08-albert-barabasi-2000}: scale-free networks tolerate
random node loss but collapse when their hubs are removed, the same
hub-dominance the contagion threshold expresses. The equivalent spectral
statement, that the epidemic threshold is set by the largest eigenvalue of
$\mathbf{A}$, makes the linear algebra of Chapter 10 the governing
apparatus: the contagion grows when $\beta/\gamma$ exceeds the reciprocal
of the adjacency matrix's spectral radius. Removing the hub lowers that
spectral radius disproportionately, which is exactly the lever the
mean-field model could not see.

% Figure: a contact network with heterogeneous degree, and the contagion
% curve on the network compared with the homogeneous-mixing mean-field SIR.
\begin{figure}[htbp]
\centering
\begin{tikzpicture}[font=\small]
% ===== Left: the contact network =====
\begin{scope}[xshift=0cm]
\node[font=\small, anchor=south west] at (-0.3, 4.0)
  {\textbf{(a) a contact network}};
% a small heterogeneous graph: one hub plus periphery
\tikzstyle{node0}=[circle, draw=engblack!70, fill=engblue!12,
  minimum size=4mm, inner sep=0pt]
\tikzstyle{hub}=[circle, draw=engblack!80, fill=engred!30,
  minimum size=7mm, inner sep=0pt, font=\scriptsize\bfseries]
\node[hub]   (h) at (2.0,2.0) {hub};
\node[node0] (a) at (0.2,3.4) {};
\node[node0] (b) at (0.0,2.0) {};
\node[node0] (c) at (0.4,0.7) {};
\node[node0] (d) at (1.8,0.2) {};
\node[node0] (e) at (3.4,0.6) {};
\node[node0] (f) at (3.9,2.1) {};
\node[node0] (g) at (3.3,3.5) {};
\node[node0] (i) at (1.9,3.8) {};
\node[node0] (j) at (4.4,3.4) {};
\foreach \n in {a,b,c,d,e,f,g,i} { \draw[draw=engblack!45] (h) -- (\n); }
\draw[draw=engblack!45] (a) -- (b);
\draw[draw=engblack!45] (b) -- (c);
\draw[draw=engblack!45] (f) -- (j);
\draw[draw=engblack!45] (g) -- (j);
\draw[draw=engblack!45] (f) -- (g);
\node[font=\scriptsize, anchor=north] at (2.0,-0.3)
  {degree-heterogeneous: a few high-degree hubs};
\end{scope}
% ===== Right: contagion curve, network vs mean-field =====
\begin{scope}[xshift=7.2cm]
\node[font=\small, anchor=south west] at (-0.3, 4.0)
  {\textbf{(b) infectious fraction over time}};
\draw[->, thick, draw=engblack!70] (0,0) -- (6.0,0)
  node[right, font=\scriptsize] {time $t$};
\draw[->, thick, draw=engblack!70] (0,0) -- (0,3.6)
  node[above, font=\scriptsize] {$I(t)/N$};
% mean-field SIR: later, taller peak
\draw[very thick, draw=engblue!80, smooth, samples=80, domain=0:5.6]
  plot (\x, { 3.0*0.46*exp(-((\x-3.0)*(\x-3.0))/0.9) });
% network: earlier, lower, fatter peak (hub seeds faster, heterogeneity
% blunts the peak)
\draw[very thick, draw=engred!85, dashed, smooth, samples=80, domain=0:5.6]
  plot (\x, { 3.0*0.30*exp(-((\x-2.2)*(\x-2.2))/1.7) });
\draw[very thick, draw=engblue!80] (3.4,3.3) -- (3.8,3.3);
\node[anchor=west, font=\scriptsize] at (3.85,3.3) {mean-field};
\draw[very thick, draw=engred!85, dashed] (3.4,2.9) -- (3.8,2.9);
\node[anchor=west, font=\scriptsize] at (3.85,2.9) {network sim.};
\end{scope}
\end{tikzpicture}
\caption[A contact network and the contagion it carries]{(a) A contact
network in which most individuals have low degree and a few hubs have
high degree, the degree heterogeneity that homogeneous mixing erases. The
infection threshold on such a network is governed not by the mean degree
but by the ratio of the second moment of the degree distribution to the
first, which the hubs inflate. (b) The infectious fraction over time,
computed on the network by individual-level stochastic simulation
(dashed) and by the homogeneous-mixing mean-field SIR of
section~\ref{fig:vol02:ch18:compartment-sir} (solid). The hubs seed the
epidemic earlier, so the network curve rises sooner; the degree
heterogeneity spreads the infection across a wider range of effective
contact rates, so the network peak is lower and broader than the
mean-field model predicts. The gap between the two curves is the
structural error the homogeneous-mixing assumption introduces, and its
sign and shape are the network model's specification.}
\label{fig:vol02:ch18:network-contagion}
\end{figure}


\subsection{Identification and the validation against the mean field}

The network model has the same two rate parameters as the SIR, $\beta$ and
$\gamma$, estimated the same way from the contact and recovery data; the
network structure itself is observed, not fitted. Before trusting the
model's answer to the hub question, the engineer validates it against the
aggregate model in the regime where the two must agree. Run the
individual-level simulation on a homogeneous network (every node the same
degree $\bar{d}$, edges placed at random) and the simulated trajectory must
reproduce the mean-field SIR curve to within the simulation's sampling
noise. This is the same cross-check the queue used: two independent
computations of the same quantity, one aggregate equation and one
stochastic simulation, that must agree or expose an error. They agree, and
the agreement licenses the simulation as a trustworthy instrument before it
is run on the structure the formula cannot handle.

Then run it on the real, heterogeneous network.
Figure~\ref{fig:vol02:ch18:network-contagion}(b) shows the divergence: on
the heterogeneous network the epidemic starts earlier (a hub, once
infected, seeds many neighbours at once) and peaks lower and broader (the
heterogeneity spreads the effective contact rate across a wide range, so
the synchronised wave the mean field predicts does not form). The final
infected fraction on the network is near $0.66$, against the mean-field
$0.73$: the homogeneous model overstates the spread, because it treats the
low-degree majority as if they mixed at the hub's rate. The gap is the
structural error the homogeneous-mixing assumption introduced, with the
sign and shape the network model now explains.

\subsection{The intervention, and the prediction}

The hub question is the one the network model exists to answer. Re-run the
simulation with the highest-degree node removed (the canteen contact point
closed). The final infected fraction drops from $0.66$ to $0.41$, a
reduction of a quarter of the workforce, because removing the hub cuts the
spectral radius of $\mathbf{A}$ and pushes the effective reproduction
number down toward threshold. The mean-field model predicted no effect; the
network model predicts a large one, and the difference is the decision. The
facilities manager now has a quantified trade-off: a quarter of the
workforce spared against the disruption of closing the canteen.

The prediction carries an interval from two sources the chapter has named
throughout. The parameter uncertainty propagates the sampling intervals on
$\hat\beta$ and $\hat\gamma$, amplified here because the network model is
more sensitive to $\beta$ near threshold than the mean field is. The
structural uncertainty is larger and is the honest part of the report: the
contact log is one week of one configuration, the proximity sensor misses
brief contacts that may still transmit, and the model assumes the network
is static while real contacts rewire daily. The engineer reports the
infected fraction as $0.66 \pm 0.08$ with the hub present and $0.41 \pm
0.10$ with it removed, states that both rest on the static-network
assumption, and names the regime (a fast-rewiring contact pattern, a
second hub the log undercounted) in which the recommendation would need
re-running on a dynamic-network model.

\subsection{The look-back and the lesson}

Does the result use all the data? Yes; the rate parameters came from the
contact and recovery logs and the structure from the full network. Does it
explain what the first model could not? Yes; the hub sensitivity is carried
by the spectral radius of the adjacency matrix, the quantity the mean field
averages away. Does it admit an independent check? Yes, twice: the
homogeneous-network simulation reproduced the mean-field SIR, and the
network threshold formula $R_0^{\text{net}}$ predicts the direction of the
hub effect that the full simulation then confirms. Does the model survive
hostile review within its domain (a closed workplace, a static one-week
contact network, a single index case, independent per-edge transmission)?
Yes, and the domain statement names the regimes (dynamic rewiring, repeated
introductions, correlated transmission within teams) where a richer model
would be needed.

The lesson generalises the aggregate examples to the structured case. When
the structure of the interactions governs the answer, the aggregate model
is not a coarser version of the right model; it is a different model that
answers a different question, and substituting it silently is the failure
the chapter's abstraction discipline exists to prevent. The discipline is
unchanged: name the question, choose the simplest model, find the question
the simple model cannot answer, let it specify the revision, validate the
revision against the simple model in their shared regime, and report the
prediction with an interval that exposes the structural uncertainty the
static-network assumption hides.

\begin{mastery}
The choice between an agent-based simulation and an aggregate equation is a
choice on a spectrum of \engterm{model resolution}, and the spectrum has a
principled middle. The fully aggregate mean-field model tracks one number
per compartment; the fully individual agent-based model tracks one state
per node. Between them sits the \engterm{degree-based} or
\engterm{pair-approximation} family, which tracks one compartment per
degree class (all degree-$d$ nodes lumped together) or one variable per
edge type, recovering most of the heterogeneity at a fraction of the
individual model's cost. The modelling judgement is to climb the resolution
ladder only as far as the question requires: the hub question forced at
least degree-resolution, but it did not require naming individuals, so a
degree-based model would have answered it more cheaply than the full
simulation. The general rule, which recurs in fluid dynamics (direct
simulation versus turbulence closure), in chemistry (every molecule versus
reaction-rate equations), and in traffic (every vehicle versus flow
density), is that resolution is a cost paid for a specific structural
question, and a modeller who simulates individuals when a degree class
would do has made the overfitting error in its computational form, buying
fidelity the decision cannot use. The agent-based model also raises a
verification burden the aggregate model does not: a stochastic simulation
must be run many times and its sampling spread quantified before its mean
is trusted, the discrete-event discipline of section~18.9 applied to a
process on a graph.
\end{mastery}

\section{Worked example, in full: closing the loop with feedback}\pathtag{standard}

Every model the chapter has built so far describes a system that does what
its dynamics dictate and is then observed. A controlled system is
different: an engineer measures its output, compares the measurement to a
goal, and feeds a correction back into the input, so that the system's
behaviour is set by the loop rather than by the plant alone. The
thermostat-controlled room, the cruise control holding a speed, the
autopilot holding a heading, the chemical reactor holding a temperature:
each is a feedback loop, and the model that predicts its behaviour must
model the loop, not just the plant. We work a feedback example at full
length because it exercises the control-loop archetype the chapter has
named in passing (Meadows's balancing loops, section~18.2) but never run
through the cycle, and because it reaches back to the room thermal model of
section~18.7 and turns it from a system to be predicted into a plant to be
controlled. The example draws on the ODEs and eigenvalues of Chapter 7, the
linear-system apparatus of Chapter 10, and the stability language the
volume has built toward.

\subsection{The question and the plant}

The room of section~18.7 warmed under a heater that was simply switched on.
A real thermostat does not switch the heater on and leave it; it modulates
the heating power to hold the room at a setpoint despite a fluctuating
outdoor temperature. The engineering question is one of control design: how
should the controller set the heating power as a function of the measured
temperature error so that the room reaches the setpoint quickly, holds it
against disturbances, and does not overshoot into discomfort or waste. The
answer is a controller gain, and the precision required is enough to keep
the steady-state error below half a degree and the overshoot below a degree.

The plant is the two-lump room model already identified: a fast air node
and a slow wall node, with the heating power $P$ driving the air. In the
open-loop model of section~18.7 the power was a fixed input; in the closed
loop it becomes the controller's output, a function of the temperature the
sensor reports.

\subsection{First candidate: proportional control}

The abstraction discipline names the question (hold the setpoint) and the
simplest controller that could answer it: proportional control, which sets
the heating power proportional to the error between the setpoint and the
measured temperature,
\[
P(t) = K_p\,\bigl(T_{\text{set}} - T_a(t)\bigr),
\]
with a single gain $K_p$ to choose. Substituting this into the air-node
equation of the two-lump plant closes the loop and turns the forced linear
system into an autonomous one whose dynamics the gain sets. The closed-loop
system is still the coupled pair of section~18.7, but the constant forcing
$P$ is replaced by the error-proportional term, which moves the system's
eigenvalues: the same room, under a controller, has different time
constants and a different steady state.

% Figure: the closed-loop block diagram of a thermostat-controlled room,
% and the closed-loop step response for three controller gains.
\begin{figure}[htbp]
\centering
\begin{tikzpicture}[font=\small,
  block/.style={draw=engblack!75, fill=engblue!10, minimum width=1.6cm,
    minimum height=1.0cm, align=center, font=\scriptsize},
  sum/.style={draw=engblack!75, circle, inner sep=0pt, minimum size=5mm},
  line/.style={->, thick, draw=engblack!70}]
% ===== Top: block diagram =====
\begin{scope}[yshift=0cm]
\node[font=\small, anchor=south west] at (-1.0, 1.1)
  {\textbf{(a) the closed loop}};
\node (ref) at (-1.0,0) {$T_{\text{set}}$};
\node[sum] (sum) at (0.2,0) {};
\node[font=\scriptsize] at (0.2,0) {$-$};
\node[block] (ctrl) at (2.0,0) {controller\\$K(s)$};
\node[block] (plant) at (4.6,0) {room\\(2-lump)};
\node (out) at (6.8,0) {$T(t)$};
\draw[line] (ref) -- (sum);
\draw[line] (sum) -- (ctrl);
\draw[line] (ctrl) -- node[above, font=\scriptsize] {$P$ (heat)} (plant);
\draw[line] (plant) -- (out);
% feedback path
\coordinate (fb) at (6.0,0);
\draw[thick, draw=engblack!70] (6.0,0) |- (6.0,-1.2) -| (0.2,-1.2);
\draw[line] (0.2,-1.2) -- (sum);
\node[block, minimum width=1.2cm, minimum height=0.8cm] (sens)
  at (3.0,-1.2) {sensor};
\node[font=\scriptsize, anchor=south] at (3.0,-1.5)
  {measured temperature feeds back};
\end{scope}
% ===== Bottom: step response for three gains =====
\begin{scope}[yshift=-5.0cm]
\node[font=\small, anchor=south west] at (-0.3, 3.7)
  {\textbf{(b) closed-loop step response}};
\draw[->, thick, draw=engblack!70] (0,0) -- (7.0,0)
  node[right, font=\scriptsize] {time $t$};
\draw[->, thick, draw=engblack!70] (0,0) -- (0,3.4)
  node[above, font=\scriptsize] {$T(t)$};
\draw[dashed, draw=engblack!45] (0,2.4) -- (7.0,2.4)
  node[right, font=\scriptsize] {$T_{\text{set}}$};
% low gain: slow, no overshoot
\draw[very thick, draw=enggray, smooth, samples=90, domain=0:6.8]
  plot (\x, { 2.4*(1 - exp(-0.55*\x)) });
% moderate gain: faster, slight overshoot
\draw[very thick, draw=engblue!80, smooth, samples=120, domain=0:6.8]
  plot (\x, { 2.4*(1 - exp(-0.9*\x)*cos(2.0*\x r)) });
% high gain: fast, large overshoot, ringing
\draw[very thick, draw=engred!85, dashed, smooth, samples=140, domain=0:6.8]
  plot (\x, { 2.4*(1 - exp(-0.5*\x)*cos(3.4*\x r)) });
\draw[very thick, draw=enggray] (4.6,3.1) -- (5.0,3.1);
\node[anchor=west, font=\scriptsize] at (5.05,3.1) {low gain};
\draw[very thick, draw=engblue!80] (4.6,2.75) -- (5.0,2.75);
\node[anchor=west, font=\scriptsize] at (5.05,2.75) {tuned};
\draw[very thick, draw=engred!85, dashed] (4.6,2.4) -- (5.0,2.4);
\node[anchor=west, font=\scriptsize] at (5.05,2.4) {high gain};
\end{scope}
\end{tikzpicture}
\caption[A thermostat as a closed loop]{(a) The thermostat-controlled room
drawn as a feedback loop. The controller compares the measured temperature
to the setpoint $T_{\text{set}}$, computes a heating power, and drives the
room, whose two-lump dynamics are the plant identified earlier in the
chapter. The sensor closes the loop. The open-loop room model becomes a
closed-loop system whose dynamics the controller, not the room alone, sets.
(b) The closed-loop step response for three controller gains. Low gain
approaches the setpoint slowly with no overshoot; a tuned gain reaches it
quickly with a small, decaying overshoot; high gain overshoots far and
rings, the signature of a loop pushed toward instability. The same plant
produces all three responses; the controller chooses among them. The
trade-off between speed and overshoot is the closed-loop analogue of the
abstraction trade-off the chapter opened with.}
\label{fig:vol02:ch18:feedback-loop}
\end{figure}


Two structural facts are available before any gain is fitted, and both come
from the eigenvalue analysis of Chapter 7 applied to the closed-loop
matrix. First, proportional control alone leaves a \engterm{steady-state
offset}: at equilibrium the heater must supply the standing heat loss, but
it can only do so if the error is nonzero, so the room settles slightly
below the setpoint, by an amount that shrinks as $K_p$ rises but never
reaches zero. Second, raising $K_p$ to reduce that offset moves the
closed-loop eigenvalues toward the imaginary axis: the response speeds up,
then begins to overshoot, then rings, and at high enough gain the slow wall
lag turns the loop oscillatory. Figure~\ref{fig:vol02:ch18:feedback-loop}(b)
shows the three regimes. The trade-off between offset and overshoot is the
proportional controller's signature, and it is the residual that specifies
the revision.

\subsection{Identification, validation, and the residual that breaks it}

The gain is tuned by running the closed loop, measuring the response to a
setpoint step, and adjusting $K_p$ until the response is as fast as it can
be without unacceptable overshoot. The validation is a held-out
disturbance: drop the outdoor temperature and check whether the controller
holds the room. Here proportional control fails its own check honestly. The
steady-state offset, predicted by the eigenvalue analysis, appears in the
data: the room holds about a degree below setpoint when the outdoor
temperature falls, and the offset grows with the disturbance. The residual
is not noise; it is the structural signature of a controller that has no
memory of accumulated error. The discipline of section~18.5 applies: the
offset is the next controller's specification.

\subsection{Second candidate: proportional-integral control}

The revision adds an integral term that accumulates the error over time,
\[
P(t) = K_p\,e(t) + K_i \int_0^t e(\tau)\,\dd\tau,
\qquad e(t) = T_{\text{set}} - T_a(t),
\]
so that any persistent error, however small, integrates into a growing
correction that drives the steady-state offset to zero. The integral term
is the controller's memory, the accumulation the proportional term lacked,
and it is exactly the stock-versus-flow distinction of section~18.2: the
proportional term responds to the flow of error, the integral term to its
accumulated stock. Adding it raises the closed-loop system's order by one,
introducing a third eigenvalue, and the gains $K_p$ and $K_i$ must now be
chosen together so that all three eigenvalues sit in the stable half-plane
with adequate damping. {\AA}str{\"o}m and Murray's text
\cite{text:astrom-murray2021} is the standard engineering treatment of the
loop-shaping this requires, and Strogatz's nonlinear-dynamics text
\cite{text:v2ch07-strogatz-nonlinear} supplies the eigenvalue-and-phase-plane
language the analysis rests on.

\subsection{Validation of the revision, and the prediction}

The proportional-integral controller's response to the same held-out
disturbance shows no surviving offset: the room returns to setpoint after
the outdoor temperature drops, with the integral term supplying the extra
standing power the proportional term could not. The cycle has closed. The
engineering answer is a gain pair that meets both specifications:
steady-state error driven to zero by the integral action, overshoot held
below a degree by the proportional gain and the integral gain chosen
together for adequate damping. The prediction the engineer reports is the
closed-loop settling time and overshoot, with an interval propagating the
uncertainty in the identified plant parameters ($C_a$, $C_w$, the
resistances) through the closed-loop eigenvalues, and a structural caveat:
the linear plant model is validated only for small excursions around the
setpoint, and a controller tuned on the linear model can misbehave if the
heater saturates (the integral term winds up against a power ceiling it
cannot exceed), the regime the next revision would have to address.

\subsection{The look-back and the lesson}

Does the result use all the data? Yes; the plant came from the section~18.7
identification and the gains from the closed-loop step and disturbance
tests. Does it explain what the first controller could not? Yes; the
steady-state offset was the residual, and the integral term is the
accumulation the residual demanded. Does it admit an independent check?
Yes; the closed-loop eigenvalues predicted from the identified plant and
the chosen gains match the settling time and overshoot measured in the
step response. Does the model survive hostile review within its domain
(small excursions, no actuator saturation, the plant time-invariant over
the test)? Yes, and the domain statement names the regimes (large setpoint
changes, a saturating heater, a plant that drifts as the building ages)
where a richer controller (anti-windup, gain scheduling, re-identification)
would be needed.

The lesson is the chapter's discipline applied to a system the engineer
acts on rather than merely observes. A feedback loop is a model whose
output the engineer chooses, and the modelling cycle runs the same way:
name the question, choose the simplest controller, read the residual (here
the steady-state offset), let it specify the revision (the integral term),
validate against a held-out disturbance, and report the prediction with the
domain of validity the linear model's small-signal assumption sets. The
control archetype recurs throughout the later volumes, wherever an engineer
closes a loop around a plant, and the modelling discipline travels with it
unchanged.

\section{Worked examples: four sharpened tools}\pathtag{standard}

The full-cycle examples ran the whole loop. Four shorter worked examples
sharpen individual tools the chapter has named, each on a question the
running examples raised but did not pursue: how correlated parameters
defeat identification, how a nonlinear model's uncertainty outruns the
linear formula, how two competing models are discriminated by evidence
rather than by fit, and how a scaling law is read off the dimensions before
any data are taken.

\subsection{Identifiability when parameters trade off}

The pendulum's $L$-and-$g$ ambiguity was a clean structural
non-identifiability: the model depended only on the ratio. The more common
and more insidious case is two parameters that are each identifiable in
principle but trade off against each other so strongly that the data
constrain only their combination. Consider the two-lump room model fit to a
warm-up curve that, by misfortune, was logged only over the first half
hour. The fast air constant $\tau_1$ and the coupling resistance $R_{aw}$
both shape the early rise, and over a short window an increase in one can
be compensated almost exactly by a decrease in the other: the fit's
residuals are nearly unchanged along a ridge in the $(\tau_1, R_{aw})$
plane. The point estimate sits somewhere on the ridge, the marginal
interval on each parameter is wide, and a naive report of the two marginal
intervals hides the fact that the parameters are not independently known.

The diagnostic is the parameter covariance matrix, the inverse of the
Fisher information the fit returns. Its off-diagonal entries, normalised to
correlations, expose the trade-off: a correlation near $\pm 1$ between two
parameters is the signature of the ridge. For the short-window room fit the
$(\tau_1, R_{aw})$ correlation is about $-0.97$, a near-perfect trade-off,
and the honest report is not two marginal intervals but the joint
confidence ellipse, long and thin along the ridge. The remedy is the one
identifiability always demands: data that break the trade-off. Logging the
warm-up long enough to see the slow wall constant separates $\tau_1$ from
$R_{aw}$, because the slow phase depends on them differently, and the
correlation drops toward zero as the ellipse rounds out. The lesson, which
the profile-likelihood mastery box anticipated, is that a model can pass
the structural-identifiability test and still be practically
unidentifiable along a parameter combination, and only the joint
covariance, not the marginal intervals, reveals it.

\subsection{Uncertainty propagation: analytic against Monte-Carlo}

Section~18.4 propagated parameter intervals through a prediction by the
first-order formula $\Var(y) \approx \sum_i S_i^2 \Var(\theta_i)$, exact
when the model is linear over the parameter intervals. The cooling-cup
inversion $t_{60} = k^{-1}\ln(\cdots)$ stretches that assumption, because
$t_{60}$ is nonlinear in $k$ and the interval on $k$ is not negligible. The
worked comparison runs both methods on the same numbers. The analytic
delta method gives $\sigma_{t} = |dt/dk|\,\sigma_k$, a symmetric interval
centred at $t_{60}$ evaluated at $\hat k$. The Monte-Carlo method draws ten
thousand values of $k$ from its fitted Gaussian, computes $t_{60}$ for
each, and reports the empirical distribution.

The two agree on the central value but disagree on the shape, and the
disagreement is instructive. Because $t_{60} \propto 1/k$ is convex, the
Monte-Carlo distribution is right-skewed: the draws with small $k$ map to
disproportionately large $t_{60}$, stretching the upper tail beyond what
the symmetric analytic interval covers.
Figure~\ref{fig:vol02:ch18:montecarlo-propagation} shows the effect on a
representative nonlinear model. The Monte-Carlo mean is shifted above the
linearisation point, and the upper $97.5$th percentile sits noticeably
higher than $\hat t_{60} + 1.96\,\sigma_t$. For the cooling cup, where the
relative interval on $k$ is six percent, the shift is small and the
analytic interval is adequate; for a model fit to sparser data, where the
relative interval reaches twenty or thirty percent, the skew matters and
the Monte-Carlo interval is the honest report. The rule the comparison
installs is to use the cheap analytic propagation when the input intervals
are small relative to the model's curvature scale, and to switch to
Monte-Carlo when they are not, with the changeover diagnosed by checking
whether the analytic and Monte-Carlo intervals agree on a single test case
before either is trusted in production.

% Figure: analytic first-order error propagation (a Gaussian) overlaid on
% the Monte-Carlo output histogram for a nonlinear model, showing the skew
% the linear approximation misses.
\begin{figure}[htbp]
\centering
\begin{tikzpicture}[font=\small]
\draw[->, thick, draw=engblack!70] (0,0) -- (9.0,0)
  node[right, font=\scriptsize] {output $y$};
\draw[->, thick, draw=engblack!70] (0,0) -- (0,4.0)
  node[above, font=\scriptsize] {density};
% Monte-Carlo histogram: right-skewed bars
\foreach \x/\h in {0.6/0.25, 1.2/0.7, 1.8/1.5, 2.4/2.6, 3.0/3.3,
  3.6/3.0, 4.2/2.3, 4.8/1.7, 5.4/1.2, 6.0/0.85, 6.6/0.55, 7.2/0.35,
  7.8/0.2, 8.4/0.1} {
  \fill[engamber!35, draw=engamber!70] (\x-0.3,0) rectangle (\x+0.3,\h);
}
% analytic Gaussian (symmetric, centred near the linearisation point):
\draw[very thick, draw=engblue!85, smooth, samples=120, domain=0.2:6.6]
  plot (\x, { 3.35*exp(-((\x-3.2)*(\x-3.2))/(2*1.0*1.0)) });
% mean markers
\draw[dashed, draw=engblue!80] (3.2,0) -- (3.2,3.5);
\node[font=\scriptsize, text=engblue!80, anchor=south, rotate=90]
  at (3.05,1.6) {analytic mean};
\draw[dashed, draw=engred!75] (3.85,0) -- (3.85,2.9);
\node[font=\scriptsize, text=engred!75, anchor=south, rotate=90]
  at (3.7,1.4) {MC mean};
% legend
\fill[engamber!35, draw=engamber!70] (6.6,3.5) rectangle (6.95,3.75);
\node[anchor=west, font=\scriptsize] at (7.0,3.62) {Monte-Carlo};
\draw[very thick, draw=engblue!85] (6.6,3.2) -- (6.95,3.2);
\node[anchor=west, font=\scriptsize] at (7.0,3.2) {first-order Gaussian};
\end{tikzpicture}
\caption[Analytic propagation versus Monte-Carlo]{The output distribution
of a nonlinear model whose inputs carry wide intervals. The first-order
error-propagation formula returns a symmetric Gaussian (solid) centred at
the model evaluated at the input means. The Monte-Carlo computation, which
runs the full nonlinear model on thousands of input draws, returns a
right-skewed histogram (bars) whose mean (red dashed) is shifted from the
linearisation point (blue dashed) and whose upper tail is heavier than the
Gaussian allows. The shift and the skew are exactly what the linear
approximation cannot represent: the curvature of the model maps a symmetric
input spread into an asymmetric output spread. When the intervals are
small and the model is nearly linear over them, the two agree and the cheap
analytic formula suffices; when the intervals are wide or the model
strongly curved, the Monte-Carlo distribution is the honest report and the
Gaussian understates the tail risk.}
\label{fig:vol02:ch18:montecarlo-propagation}
\end{figure}


\subsection{Model discrimination by evidence, not by fit}

Section~18.5's model selection penalised complexity to choose among nested
candidates. The complementary problem is discriminating between two
non-nested models that fit comparably well, where neither is a special case
of the other. The cooling cup admits a Newton's-law model (exponential
decay toward ambient) and a competing lumped-radiative model (a different
nonlinear decay law); both fit the thirty-point notebook with similar
residual sums. Choosing by training residual alone would be arbitrary,
because the two are close, and choosing by parameter count is moot because
both have the same number of parameters.

The discriminating quantity is the evidence the data give each model,
formalised as the likelihood ratio or, in the Bayesian framing, the Bayes
factor. The likelihood-ratio test compares the maximised likelihoods of the
two models, with the larger likelihood favouring the model that makes the
observed data more probable. For the cooling cup at modest temperature
excess the ratio is near one, the data cannot tell the two models apart,
and the honest report is that both are admissible and the choice should be
made on other grounds (Newton's law is simpler and its parameter is
physically interpretable). For a hotter notebook, where the radiative term
of the competing model bites, the likelihood ratio swings decisively toward
the radiative model, and the data now discriminate. The lesson is that
model discrimination is a property of the data's reach, not of the
modeller's preference: two models that the available data cannot separate
should both be carried, the triangulation discipline of section~18.5, until
data in a regime that separates them are acquired. The information criteria
of section~18.5 are the practical implementation, and Hastie, Tibshirani,
and Friedman \cite{text:v2ch18-hastie-statistical-learning-2009} connect
the likelihood ratio, the information criteria, and cross-validation as
three estimates of the same held-out predictive performance.

\subsection{Dimensional analysis reads a scaling law off the dimensions}

The Buckingham-Pi machinery of section~18.3 fixed the drag law and the
cantilever frequency before any dynamics were written. A third worked case
shows the move on a heat-transfer question the reader can pose about the
cooling cup itself, and it earns its place by deriving a result that the
chapter's compartment model assumed. How long does a solid body of
characteristic size $\ell$ take to cool through its interior, as opposed to
losing heat from its surface? The relevant variables are the size $\ell$
($\mathsf{L}$), the thermal diffusivity $\alpha$
($\mathsf{L}^2\mathsf{T}^{-1}$), and the time $t$ ($\mathsf{T}$). Three
variables, two primary dimensions, so the Buckingham count gives one
dimensionless group, the \engterm{Fourier number}
\[
\mathrm{Fo} = \frac{\alpha t}{\ell^2}.
\]
Dimensional homogeneity forces the interior temperature to depend on time
only through $\mathrm{Fo}$, so the internal equilibration time scales as
$t \sim \ell^2/\alpha$: the diffusion time grows with the square of the
size. This is the structural fact the lumped-capacitance model of
section~18.7 quietly assumed when it treated the room's air as a single
well-mixed node. The assumption is valid only when the interior
equilibrates fast compared with the surface loss, a condition the
\engterm{Biot number} (the ratio of internal to surface resistance) makes
precise: a small Biot number licenses the single-lump abstraction, and a
large one demands the spatial resolution the lump throws away. The
dimensional analysis thus both supplies the scaling law and names the
dimensionless threshold at which the chapter's own earlier abstraction
breaks, a closing demonstration that the cheapest structural information in
any model comes, as Rule Three insisted, from the dimensions before the
dynamics.

\begin{mastery}
The Fourier-and-Biot analysis points to the most general form of the
abstraction the chapter has used repeatedly without naming: \engterm{model
order reduction}. A spatially distributed system obeys a partial
differential equation with infinitely many degrees of freedom; the
lumped-capacitance model replaces it with one. The replacement is justified
when one mode (here the spatially uniform mode) dominates the response in
the regime of interest, and the systematic version of the move projects the
full equation onto the few modes that carry the energy. \engterm{Proper
orthogonal decomposition} is the data-driven form: collect snapshots of the
full field, compute the dominant modes by the singular value decomposition
of Chapter 10, and keep the handful whose singular values dominate. The
reduced model evolves only the coefficients of those modes, turning a
field equation into a small system of ordinary differential equations that
runs orders of magnitude faster while reproducing the dominant dynamics.
The discipline of this chapter applies to the reduced model unchanged: it
has a domain of validity (the regime spanned by the snapshots it was built
from), it has a structural error (the energy in the discarded modes), and
it fails the regime-extension test if queried outside the snapshot regime,
the same failure mode the surrogate model of section~18.4's mastery box
exhibits. Reduced-order modelling, surrogate modelling, and the
single-lump abstraction are three points on one continuum: each trades
fidelity for speed by keeping the modes the question needs and discarding
the rest, and each is honest only when the discarded modes are documented
as the model's structural error.
\end{mastery}

\begin{mastery}
The parameter-identification problem this chapter treats by least squares
and maximum likelihood has a Bayesian counterpart that is worth the
engineer's acquaintance because it handles the cases the point-estimate
methods strain on: wide priors, strong parameter correlations, and the
genuine quantification of what the data leave unknown. \engterm{Bayesian
calibration} treats the parameters as random variables with a prior
distribution encoding what was known before the data, updates the prior by
the likelihood to a posterior distribution encoding what is known after,
and reports the posterior rather than a point and an interval. The
posterior is the honest object: its spread is the parameter uncertainty,
its correlations are the trade-offs the covariance matrix approximates, and
its shape carries the skew and the bounds that a Gaussian interval cannot.
For the room model, a Bayesian calibration would return a joint posterior
over $(C_a, C_w, R_{aw}, R_{ao})$ whose ridge structure exposes the
short-window non-identifiability directly, without the linearisation the
condition number assumes. The computation is the inverse problem of
inferring causes (the parameters) from effects (the data), and it is
carried out by sampling the posterior with Markov-chain Monte-Carlo, the
apparatus Gelman and coauthors' text \cite{text:ch08-gelman-bda3-2013}
develops in full. The Bayesian framing also makes the structural-error
question sharper: the Kennedy-O'Hagan formulation adds an explicit
discrepancy term between the model and the world, calibrated alongside the
parameters, so that the predictive interval carries the model's own
inadequacy rather than absorbing it silently into the parameter estimates.
The engineering reader takes from this the working rule that when the
parameters are correlated, the priors are informative, or the structural
error must be quantified rather than bounded, the Bayesian posterior is the
report the decision deserves, at the cost of the sampling machinery the
point estimate avoids.
\end{mastery}

\begin{mastery}
The full-cycle examples each ended with a domain-of-validity statement and
a warning that the model would need re-running if the system drifted. The
contemporary name for a model kept continuously synchronised with its
physical counterpart is a \engterm{digital twin}: a model that ingests the
live sensor stream from the system it represents, re-identifies its
parameters as the system ages, and runs ahead of the system to predict the
consequences of a contemplated action before the action is taken. The
two-lump room model becomes a digital twin when it is fed the building's
live temperature log, re-fits $C_a$ and the resistances as the insulation
degrades or the furniture changes, and predicts the pre-heat schedule for
tomorrow morning. The digital twin is not a new kind of model; it is this
chapter's modelling cycle run continuously rather than once, with
identification, validation, and revision automated and repeated as new data
arrive. Its characteristic failure is the chapter's own simulation-became-
the-system mode (section~18.8) at industrial scale: a twin trusted as a
substitute for measurement, its predictions acted on without the held-out
validation that would catch its drift, becomes a confident liar exactly
when the system has changed enough to matter. The discipline that keeps a
digital twin honest is the discipline that keeps any model honest, the
held-out validation against data the re-identification did not use,
performed on a schedule, with rising validation residuals as the trigger to
question the twin rather than the system. The engineering reader takes the
digital twin as a deployment pattern for the chapter's cycle, not as an
escape from its discipline.
\end{mastery}

\section{Worked examples: the moves on unfamiliar systems}\pathtag{standard}

The full-cycle examples ran the loop on cooling, heating, queues, networks,
and feedback. The moves are the same on systems whose mathematics the
reader has not yet met in a modelling context. We work a set of shorter
vignettes on chemistry, biochemistry, fluid friction, and online
estimation, each pursuing one move to its useful conclusion, and then two
applied vignettes at full length on problems an engineer is actually paid
to answer.

\subsection{A second Buckingham collapse: pipe friction}

Section~18.3 collapsed the drag on a sphere onto a single curve. The same
machinery collapses the pressure drop in a pipe, and the second case earns
its place because it is the one an engineer sizing a water main or a
ventilation duct actually uses. The pressure drop $\Delta p$
($\mathsf{M}\mathsf{L}^{-1}\mathsf{T}^{-2}$) along a pipe can depend on the
pipe length $\ell$ ($\mathsf{L}$), the diameter $d$ ($\mathsf{L}$), the mean
flow speed $U$ ($\mathsf{L}\mathsf{T}^{-1}$), the fluid density $\rho$
($\mathsf{M}\mathsf{L}^{-3}$), the viscosity $\mu$
($\mathsf{M}\mathsf{L}^{-1}\mathsf{T}^{-1}$), and the wall roughness
$\varepsilon$ ($\mathsf{L}$). Seven variables, three primary dimensions, so
the count gives four dimensionless groups. Choosing $\rho$, $U$, $d$ as the
repeating set, the groups are the pressure coefficient
$\Delta p/(\rho U^2)$, the Reynolds number $\mathrm{Re} = \rho U d/\mu$, the
length ratio $\ell/d$, and the relative roughness $\varepsilon/d$. The
pressure coefficient is known empirically to be linear in $\ell/d$ (a pipe
twice as long drops twice the pressure at fixed everything else), which
peels that group off and leaves the \engterm{Darcy friction factor}
\[
f = \frac{\Delta p}{(\ell/d)\,\tfrac{1}{2}\rho U^2}
  = \phi\!\left(\mathrm{Re},\ \frac{\varepsilon}{d}\right),
\]
a function of two groups rather than of six raw variables. Dimensional
analysis has reduced a six-variable experiment to a two-variable chart.

Figure~\ref{fig:vol02ch18:pipe-friction-collapse} is that chart, the Moody
collapse \cite{paper:moody-1944}. In the laminar regime the dynamics supply
the constant that dimensional analysis could not: the Hagen-Poiseuille
solution gives $f = 64/\mathrm{Re}$ exactly, a straight line of slope $-1$
independent of roughness. In the turbulent regime the friction factor
flattens toward a plateau set by $\varepsilon/d$, the roughness the laminar
branch never felt. The modelling lesson repeats the sphere-drag lesson with
one addition: a single scalar function of two groups, measured once,
predicts the pressure drop for every pipe, fluid, and flow rate an engineer
will meet, and the regime boundary between the two branches, near
$\mathrm{Re} \approx 2300$, is itself a dimensionless threshold the chart
makes visible.

% Figure: friction factor versus Reynolds number for pipe flow, a second
% Buckingham-Pi collapse distinct from the sphere-drag case. The laminar
% branch f = 64/Re and a turbulent plateau, with two roughness curves.
\begin{figure}[htbp]
\centering
\begin{tikzpicture}[font=\small]
% axes
\draw[->, thick, draw=engblack!70] (0,0) -- (8.4,0)
  node[right, font=\scriptsize] {$\log_{10}\mathrm{Re}$};
\draw[->, thick, draw=engblack!70] (0,0) -- (0,5.2)
  node[above, font=\scriptsize] {$\log_{10} f$ (friction factor)};
% x ticks: Re = 10^2 .. 10^7
\foreach \xt/\xv in {0.6/2, 2.0/3, 3.4/4, 4.8/5, 6.2/6, 7.6/7} {
  \draw[draw=engblack!70] (\xt,0) -- (\xt,-0.07);
  \node[font=\scriptsize, anchor=north] at (\xt,-0.09) {\xv};
}
% y ticks: log f from -2 to -0.5, drawn on a shifted scale
\foreach \yt/\yv in {0.6/-2.0, 2.1/-1.5, 3.6/-1.0, 5.0/-0.5} {
  \draw[draw=engblack!70] (0,\yt) -- (-0.07,\yt);
  \node[font=\scriptsize, anchor=east] at (-0.09,\yt) {\yv};
}
% laminar branch: f = 64/Re, so log f = log 64 - log Re, slope -1.
% map: at Re=10^2 (x=0.6) f=0.64 -> log f = -0.19; at transition Re~2300.
% Draw a decreasing line from upper-left to the transition.
\draw[very thick, draw=engblue!80] (0.6, 4.9) -- (2.35, 3.15);
\node[font=\scriptsize, color=engblue!80, anchor=south west, rotate=-38] at (0.9,4.4)
  {laminar $f=64/\mathrm{Re}$};
% transition region (dashed vertical band)
\draw[dashed, draw=enggray] (2.35,0) -- (2.35,4.6);
\node[font=\scriptsize, color=enggray, anchor=south] at (2.35,4.6) {transition};
% turbulent branch, smooth pipe: gently decreasing then flattening
\draw[very thick, draw=engred!80]
  (2.6,2.9) .. controls (4.0,2.3) and (5.6,1.85) .. (7.8,1.55);
\node[font=\scriptsize, color=engred!80, anchor=south] at (6.3,1.75) {smooth pipe};
% turbulent branch, rough pipe: flattens to a plateau set by roughness
\draw[very thick, draw=engorange!85]
  (2.6,3.15) .. controls (4.2,2.85) and (5.4,2.72) .. (7.8,2.68);
\node[font=\scriptsize, color=engorange!85, anchor=south] at (6.3,2.85)
  {rough pipe ($\varepsilon/d$ fixed)};
% data points collapsed onto the curves (three pipe diameters, one symbol each)
\foreach \px/\py in {0.9/4.55, 1.4/4.05, 1.9/3.55, 3.2/2.55, 4.6/2.1, 6.0/1.8, 7.2/1.62} {
  \filldraw[engblack!75] (\px, \py) circle (0.055);
}
\end{tikzpicture}
\caption[Pipe-friction factor versus Reynolds number as a Buckingham-Pi
collapse]{Friction factor $f$ against Reynolds number for flow in a
circular pipe, a second data collapse of the kind the drag-on-a-sphere
case produced. Pressure drop measured on pipes of different diameter,
length, and fluid, plotted as raw pressure drop against flow rate, forms a
scatter of separate families; plotted as the dimensionless $f$ against
$\mathrm{Re}$ the families collapse onto two universal branches. The
laminar branch $f = 64/\mathrm{Re}$ is a straight line of slope $-1$ that
dimensional analysis and the Hagen-Poiseuille solution together fix
exactly; the turbulent branch flattens toward a plateau set by the second
dimensionless group, the relative roughness $\varepsilon/d$. A modeller
who measured one pipe reads the whole family off the collapsed curve, the
same reduction the sphere-drag case bought.}
\label{fig:vol02ch18:pipe-friction-collapse}
\end{figure}


\begin{estimation}
\textbf{Question.} A garden hose $20\,\text{m}$ long and $12\,\text{mm}$ in
bore carries water at $0.2\,\text{L/s}$. Estimate the Reynolds number, name
the regime, and estimate the pressure drop, before consulting the chart.

\smallskip
\textit{Estimate, before any calculation.}

\smallskip
The cross-section is $A = \pi d^2/4 \approx \pi (0.012)^2/4 \approx
1.1\times 10^{-4}\,\text{m}^2$, so the mean speed is $U = Q/A \approx
2\times 10^{-4}/1.1\times 10^{-4} \approx 1.8\,\text{m/s}$. With water's
$\rho \approx 10^3\,\si{\kilogram\per\cubic\meter}$ and $\mu \approx
10^{-3}\,\text{Pa}\cdot\text{s}$, the Reynolds number is $\mathrm{Re} =
\rho U d/\mu \approx 10^3 \cdot 1.8 \cdot 0.012/10^{-3} \approx 2.2\times
10^4$, comfortably turbulent. A smooth-pipe turbulent friction factor at
that Reynolds number is near $f \approx 0.026$ from the chart. The pressure
drop is $\Delta p = f (\ell/d)\tfrac{1}{2}\rho U^2 \approx 0.026 \cdot
(20/0.012) \cdot 0.5 \cdot 10^3 \cdot 1.8^2 \approx 0.026 \cdot 1667 \cdot
1620 \approx 7\times 10^4\,\text{Pa}$, roughly $0.7\,\text{atm}$.

\smallskip
\textbf{Postmortem.} The dimensionless groups did the work. The Reynolds
number placed the flow on the turbulent branch before any friction factor
was read, and the length ratio $\ell/d \approx 1700$ shows why a long thin
hose drops so much pressure: the drop scales with $\ell/d$ directly. An
engineer who wanted half the drop would reach first for a larger bore,
because the drop scales as roughly $d^{-5}$ at fixed flow rate (through both
the $\ell/d$ factor and the speed's dependence on area), the strongest
lever the dimensional form exposes.
\end{estimation}

\subsection{Linearise to identify: enzyme kinetics}

The chapter's identifications linearised where they could: the cooling cup
regressed a logarithm, the pendulum regressed $T^2$ on $L$. Enzyme kinetics
is the case where the linearising transform is famous, instructive, and a
trap the modelling discipline must name. The Michaelis-Menten rate law
\cite{paper:michaelis-menten1913},
\[
v = \frac{V_{\max}\,[S]}{K_m + [S]},
\]
relates a reaction rate $v$ to a substrate concentration $[S]$ through two
parameters: the maximum rate $V_{\max}$ and the half-saturation constant
$K_m$, the concentration at which the rate reaches half its ceiling.
Figure~\ref{fig:vol02ch18:michaelis-menten}(a) shows the saturating curve.
The parameters are nonlinear in the rate, but the double-reciprocal
transform
\[
\frac{1}{v} = \frac{K_m}{V_{\max}}\,\frac{1}{[S]} + \frac{1}{V_{\max}}
\]
is linear in $1/[S]$, so a straight-line fit of $1/v$ on $1/[S]$ reads
$1/V_{\max}$ off the intercept and $K_m/V_{\max}$ off the slope, the
Lineweaver-Burk plot of
figure~\ref{fig:vol02ch18:michaelis-menten}(b). The transform is the
related-problem move of P\'olya's discipline, and for a quick graphical
estimate it is exactly right.

The trap is the noise. The reciprocal $1/v$ magnifies the error on small
rates, which are the measurements at low substrate, so the transformed fit
weights the least reliable points most heavily and returns a biased
$\hat K_m$. The modelling discipline of section~18.4 names the defence: the
transform supplies starting values, and the honest estimate comes from a
nonlinear least-squares fit on the untransformed rate, with weights that
reflect the actual measurement variance. A reader who reports the
Lineweaver-Burk intercepts as final parameters has committed the noise-model
error the least-squares section warned against, in one of its most cited
historical forms. The vignette is small; the discipline is the one the
whole chapter installs: a linearising transform changes the noise model, and
a modeller who transforms to identify must fit on the scale the noise is
additive on.

% Figure: Michaelis-Menten saturation and its Lineweaver-Burk linearisation.
% (a) rate versus substrate, saturating; (b) double-reciprocal straight line.
\begin{figure}[htbp]
\centering
\begin{tikzpicture}[font=\small]
% ---- panel (a): saturation curve ----
\begin{scope}
\draw[->, thick, draw=engblack!70] (0,0) -- (5.2,0)
  node[right, font=\scriptsize] {$[S]$};
\draw[->, thick, draw=engblack!70] (0,0) -- (0,4.0)
  node[above, font=\scriptsize] {$v$};
% v = Vmax S/(Km+S), Vmax=3.4 (drawn), Km at half-max
\draw[very thick, draw=engblue!80] plot[domain=0:5,samples=60]
  (\x, {3.4*\x/(1.2+\x)});
% Vmax asymptote
\draw[dashed, draw=enggray] (0,3.4) -- (5,3.4);
\node[font=\scriptsize, color=enggray, anchor=east] at (5.0,3.55) {$V_{\max}$};
% half-max line and Km
\draw[dashed, draw=engred!70] (0,1.7) -- (1.2,1.7) -- (1.2,0);
\node[font=\scriptsize, color=engred!70, anchor=north] at (1.2,-0.05) {$K_m$};
\node[font=\scriptsize, anchor=north] at (2.6,-0.5) {(a) saturation};
% data points
\foreach \px in {0.3,0.7,1.2,1.9,2.8,4.0} {
  \pgfmathsetmacro\vy{3.4*\px/(1.2+\px) + 0.09*sin(\px*300)}
  \filldraw[engblack!75] (\px, {\vy}) circle (0.05);
}
\end{scope}
% ---- panel (b): Lineweaver-Burk ----
\begin{scope}[xshift=7.2cm]
\draw[->, thick, draw=engblack!70] (-1.4,0) -- (5.2,0)
  node[right, font=\scriptsize] {$1/[S]$};
\draw[->, thick, draw=engblack!70] (0,-0.2) -- (0,4.0)
  node[above, font=\scriptsize] {$1/v$};
% 1/v = (Km/Vmax)(1/S) + 1/Vmax; slope Km/Vmax, intercept 1/Vmax.
% choose intercept 0.6, slope 0.55; x-intercept = -1/Km
\draw[very thick, draw=engblue!80] (-1.1, 0.0) -- (5.0, 3.36);
% y-intercept marker
\filldraw[engred!80] (0,0.605) circle (0.05);
\node[font=\scriptsize, color=engred!80, anchor=west] at (0.15,0.75) {$1/V_{\max}$};
% x-intercept marker
\filldraw[engred!80] (-1.1,0) circle (0.05);
\node[font=\scriptsize, color=engred!80, anchor=south] at (-1.1,0.1) {$-1/K_m$};
% slope label
\node[font=\scriptsize, color=engblue!80, anchor=north west, rotate=28] at (2.5,1.7)
  {slope $=K_m/V_{\max}$};
% data points on the line
\foreach \px in {-0.9,-0.4,0.4,1.3,2.4,3.7} {
  \pgfmathsetmacro\vy{0.605+0.55*\px + 0.06*sin(\px*400)}
  \filldraw[engblack!75] (\px, {\vy}) circle (0.05);
}
\node[font=\scriptsize, anchor=north] at (2.0,-0.5) {(b) double reciprocal};
\end{scope}
\end{tikzpicture}
\caption[Michaelis-Menten kinetics and the Lineweaver-Burk
linearisation]{Enzyme reaction rate against substrate concentration. Panel
(a): the rate $v = V_{\max}[S]/(K_m+[S])$ rises toward the saturation
$V_{\max}$, with the half-saturation constant $K_m$ read where the rate
reaches half its ceiling. Panel (b): the double-reciprocal transform
$1/v = (K_m/V_{\max})(1/[S]) + 1/V_{\max}$ turns the saturating curve into
a straight line whose intercepts identify both parameters. The
transformation is a worked instance of the chapter's linearise-to-identify
move, with the working caution that the reciprocal reweights the noise:
low-substrate points, where $1/[S]$ is large, dominate the fit, so the
graphical read gives starting values while the honest interval comes from
a nonlinear fit on the untransformed rate.}
\label{fig:vol02ch18:michaelis-menten}
\end{figure}


\subsection{Two abstractions for one chemistry: the reactor}

The single-tank compartment model of section~18.5 treated a vessel as one
well-mixed stock. A reactor engineer has a choice of two abstractions for
the same reaction, and the choice changes the answer by a factor the
dimensionless group makes precise. The continuous stirred-tank reactor
(CSTR) keeps the single-compartment abstraction: fresh feed mixes instantly
into the vessel, and the whole vessel sits at the outlet concentration. The
plug-flow reactor (PFR) refuses that abstraction: the concentration falls
along the tube, and no single lumped state describes it. For a first-order
reaction with rate constant $k$ and mean residence time $\tau$, the
governing dimensionless group is the Damk\"ohler number $\mathrm{Da} =
k\tau$, the ratio of the residence time to the reaction timescale. The two
models give
\[
X_{\text{CSTR}} = \frac{\mathrm{Da}}{1 + \mathrm{Da}},
\qquad
X_{\text{PFR}} = 1 - e^{-\mathrm{Da}},
\]
for the conversion $X$, the fraction of feed reacted.
Figure~\ref{fig:vol02ch18:reactor-abstraction} plots both.

At the same residence time the tube out-converts the tank, and the gap
widens with $\mathrm{Da}$: to reach $90\%$ conversion the plug-flow reactor
needs $\mathrm{Da} = \ln 10 \approx 2.3$, while the tank needs
$\mathrm{Da} = 9$, four times the residence time and therefore four times
the volume. The reason is structural, not numerical: the tank dilutes fresh
high-concentration feed into already-converted product, so its reaction runs
everywhere at the low outlet concentration, while the tube lets the feed
react at its full concentration on entry. The modelling decision, which
abstraction to use, is not a matter of taste; it is the difference between a
reactor sized right and one four times too large. The engineer who picks the
single-tank abstraction by habit, because the compartment model was simplest
in section~18.5, has let a modelling convenience set a capital cost.

% Figure: two reactor abstractions for the same chemistry. (a) schematic of
% a well-mixed tank (CSTR) versus a plug-flow tube (PFR); (b) conversion
% versus dimensionless residence time (Damkohler number) for each.
\begin{figure}[htbp]
\centering
\begin{tikzpicture}[font=\small]
% ---- panel (a): schematics ----
\begin{scope}
% CSTR: a tank with a stirrer
\draw[thick, draw=engblack!75] (0,1.2) rectangle (1.8,3.0);
\draw[very thick, draw=engblue!80] (0.9,3.0) -- (0.9,2.3);
\draw[very thick, draw=engblue!80] (0.5,2.3) -- (1.3,2.3);
\node[font=\scriptsize, anchor=north] at (0.9,1.1) {CSTR (well mixed)};
% inflow / outflow arrows
\draw[->, thick, draw=enggray] (-0.7,2.7) -- (0.0,2.7);
\draw[->, thick, draw=enggray] (1.8,1.6) -- (2.5,1.6);
% concentration inside: uniform (one shade)
\filldraw[engblue!12, draw=none] (0.05,1.25) rectangle (1.75,2.95);
% PFR: a tube with a gradient
\draw[thick, draw=engblack!75] (3.6,2.2) rectangle (6.6,3.0);
\node[font=\scriptsize, anchor=north] at (5.1,2.0) {PFR (plug flow)};
% gradient blocks along the tube
\foreach \i/\sh in {0/40, 1/30, 2/22, 3/15, 4/9, 5/5} {
  \pgfmathsetmacro\xa{3.6+\i*0.5}
  \filldraw[engorange!\sh, draw=none] (\xa,2.22) rectangle (\xa+0.5,2.98);
}
\draw[->, thick, draw=enggray] (2.9,2.6) -- (3.6,2.6);
\draw[->, thick, draw=enggray] (6.6,2.6) -- (7.3,2.6);
\end{scope}
% ---- panel (b): conversion vs Damkohler ----
\begin{scope}[yshift=-4.6cm, xshift=1.2cm]
\draw[->, thick, draw=engblack!70] (0,0) -- (6.4,0)
  node[right, font=\scriptsize] {Damk\"ohler number $\mathrm{Da}$};
\draw[->, thick, draw=engblack!70] (0,0) -- (0,3.8)
  node[above, font=\scriptsize] {conversion $X$};
% full conversion line
\draw[dashed, draw=enggray] (0,3.4) -- (6.0,3.4);
\node[font=\scriptsize, color=enggray, anchor=east] at (6.0,3.6) {$X=1$};
% PFR (first order): X = 1 - exp(-Da)
\draw[very thick, draw=engorange!85] plot[domain=0:6,samples=60]
  (\x, {3.4*(1-exp(-\x))});
\node[font=\scriptsize, color=engorange!85, anchor=west] at (3.2,3.15) {PFR};
% CSTR (first order): X = Da/(1+Da)
\draw[very thick, draw=engblue!85] plot[domain=0:6,samples=60]
  (\x, {3.4*\x/(1+\x)});
\node[font=\scriptsize, color=engblue!85, anchor=north west] at (3.2,2.35) {CSTR};
\end{scope}
\end{tikzpicture}
\caption[Two reactor abstractions for one chemistry]{The same first-order
reaction modelled two ways. Panel (a): the continuous stirred-tank reactor
treats the vessel as one well-mixed compartment at a uniform concentration,
the single-tank abstraction of the compartment section; the plug-flow
reactor treats the tube as a concentration that falls along its length, a
distributed model with no single lumped state. Panel (b): conversion
against the dimensionless residence time, the Damk\"ohler number. For the
same residence time the plug-flow reactor reaches higher conversion, its
$X = 1 - e^{-\mathrm{Da}}$ climbing faster than the tank's
$X = \mathrm{Da}/(1+\mathrm{Da})$, because the tank dilutes fresh feed into
already-converted product while the tube does not. The choice between the
two abstractions is a modelling decision the dimensionless group makes
quantitative, not a detail of geometry.}
\label{fig:vol02ch18:reactor-abstraction}
\end{figure}


\subsection{The estimate that updates itself: recursive identification}

Every identification in the chapter processed a batch: all the data in hand,
one fit, one interval. A large class of engineering systems cannot wait for
the batch, because the data arrive one at a time and a decision is needed at
every step. A navigation filter estimates position as each satellite fix
arrives; a process controller estimates a drifting setpoint as each sample
lands; a battery gauge estimates charge as each current reading accumulates.
The modelling move is \engterm{recursive estimation}: keep a current best
estimate and its uncertainty, and update both as each measurement arrives,
without re-processing the history.

For a constant quantity measured with independent noise, the update is a
precision-weighted average. Let $\hat\theta_{n-1}$ be the estimate after
$n-1$ measurements, with variance $\sigma_{n-1}^2$, and let a new
measurement $y_n$ arrive with variance $\sigma_y^2$. The updated estimate
weights the two by their precisions (inverse variances),
\[
\hat\theta_n = \hat\theta_{n-1} +
  \frac{\sigma_{n-1}^2}{\sigma_{n-1}^2 + \sigma_y^2}\,
  \bigl(y_n - \hat\theta_{n-1}\bigr),
\qquad
\frac{1}{\sigma_n^2} = \frac{1}{\sigma_{n-1}^2} + \frac{1}{\sigma_y^2},
\]
the second relation showing the precision accumulating by simple addition.
The correction is the measurement residual $y_n - \hat\theta_{n-1}$ scaled
by a gain that shrinks as the estimate sharpens: an early measurement, when
the estimate is uncertain, moves it a lot; a late measurement, when the
estimate is already tight, barely moves it.
Figure~\ref{fig:vol02ch18:recursive-estimation} shows the estimate
converging while its uncertainty band narrows as $1/\sqrt{n}$, the same rate
the batch mean achieves and, for this linear-Gaussian case, exactly the
batch answer computed incrementally. The recursive form is the seed of the
Kalman filter the later volumes develop, where the quantity also drifts
between measurements and the update carries a model of that drift; the
modelling content already visible here is that an estimate is a value and an
uncertainty travelling together, and every measurement is worth exactly the
precision it adds.

% Figure: recursive (online) estimation. The running estimate of a
% constant converges toward the truth as measurements arrive, and the
% uncertainty band shrinks as 1/sqrt(n).
\begin{figure}[htbp]
\centering
\begin{tikzpicture}[font=\small]
\draw[->, thick, draw=engblack!70] (0,0) -- (9.2,0)
  node[right, font=\scriptsize] {measurement number $n$};
\draw[->, thick, draw=engblack!70] (0,0) -- (0,4.6)
  node[above, font=\scriptsize] {estimate};
% truth line
\draw[thick, draw=enggray, dashed] (0,2.6) -- (8.8,2.6);
\node[font=\scriptsize, color=enggray, anchor=west] at (8.0,2.8) {truth};
% uncertainty band: shrinks like 1/sqrt(n) around the running mean.
% Build upper and lower envelopes as piecewise lines.
% running mean wanders then settles at 2.6; band half-width ~ 1.6/sqrt(n)
\foreach \n/\m in {1/3.7, 2/3.2, 3/2.2, 4/2.9, 5/2.45, 6/2.75, 7/2.5, 8/2.62, 9/2.58, 10/2.6} {
  \pgfmathsetmacro\xx{0.85*\n}
  \pgfmathsetmacro\hw{1.7/sqrt(\n)}
}
% band as a filled region (approximate): sample the envelope
\filldraw[engblue!12, draw=none]
  (0.85,{3.55+1.7}) --
  (1.70,{3.10+1.7/sqrt(2)}) --
  (2.55,{2.70+1.7/sqrt(3)}) --
  (3.40,{2.75+1.7/sqrt(4)}) --
  (4.25,{2.60+1.7/sqrt(5)}) --
  (5.10,{2.68+1.7/sqrt(6)}) --
  (5.95,{2.58+1.7/sqrt(7)}) --
  (6.80,{2.62+1.7/sqrt(8)}) --
  (7.65,{2.60+1.7/sqrt(9)}) --
  (8.50,{2.60+1.7/sqrt(10)}) --
  (8.50,{2.60-1.7/sqrt(10)}) --
  (7.65,{2.60-1.7/sqrt(9)}) --
  (6.80,{2.62-1.7/sqrt(8)}) --
  (5.95,{2.58-1.7/sqrt(7)}) --
  (5.10,{2.68-1.7/sqrt(6)}) --
  (4.25,{2.60-1.7/sqrt(5)}) --
  (3.40,{2.75-1.7/sqrt(4)}) --
  (2.55,{2.70-1.7/sqrt(3)}) --
  (1.70,{3.10-1.7/sqrt(2)}) --
  (0.85,{3.55-1.7}) -- cycle;
% running estimate line
\draw[very thick, draw=engblue!85]
  (0.85,3.55) -- (1.70,3.10) -- (2.55,2.70) -- (3.40,2.75) --
  (4.25,2.60) -- (5.10,2.68) -- (5.95,2.58) -- (6.80,2.62) --
  (7.65,2.60) -- (8.50,2.60);
% individual noisy measurements
\foreach \xx/\yy in {0.85/4.3, 1.70/2.1, 2.55/1.6, 3.40/3.4, 4.25/2.0, 5.10/3.1, 5.95/2.1, 6.80/2.9, 7.65/2.4, 8.50/2.7} {
  \filldraw[engred!70] (\xx,\yy) circle (0.05);
}
\node[font=\scriptsize, color=engred!70, anchor=west] at (1.9,1.4)
  {individual measurements};
\node[font=\scriptsize, color=engblue!85, anchor=south west] at (3.6,3.0)
  {running estimate $\pm$ band};
\end{tikzpicture}
\caption[Recursive estimation of a constant]{A running estimate of a
constant quantity updated as each measurement arrives. The individual
measurements scatter widely; the recursive estimate, which blends the
prior estimate with each new datum in proportion to their relative
precisions, converges toward the truth. The shaded uncertainty band
narrows as $1/\sqrt{n}$, the same rate the batch mean achieves, but the
recursive form produces the current best estimate and its interval at
every step without re-processing the whole history. This is the update
rule at the heart of online identification and of the Kalman filter the
later volumes develop; the modelling content is that an estimate carries
its uncertainty forward and shrinks it by exactly the information each new
measurement adds.}
\label{fig:vol02ch18:recursive-estimation}
\end{figure}


\begin{mastery}
The recursive update above assumed the estimated quantity is constant. When
it drifts, the update needs a model of the drift, and the resulting recursion
is the \engterm{Kalman filter}: a predict step that advances the estimate and
inflates its uncertainty by the process model between measurements, and an
update step that is exactly the precision-weighted correction above. The
filter is the chapter's whole cycle run at measurement cadence. The predict
step is the model making a prediction with a stated interval; the update step
is the validation against a new datum; the gain is the honest weighting of
model against measurement by their relative precisions; and a filter whose
predicted residuals (the innovations) show structure is a filter whose
process model is wrong, the residual-reads-the-revision discipline of
section~18.5 applied online. The characteristic failure is a filter tuned to
trust its model too far: a process-noise setting too small makes the gain
collapse, the filter stops listening to measurements, and it tracks its own
prediction off into error while reporting a confident interval, the
simulation-became-the-system failure of section~18.8 in recursive form. The
engineering reader takes the Kalman filter as the deployment of this
chapter's discipline for systems that are measured continuously and must be
estimated in real time, and takes the innovation sequence as the held-out
residual that keeps such a filter honest.
\end{mastery}

\subsection{Applied vignette: sizing a rainwater-harvesting tank}

An engineer is asked to size the storage tank for a building that will
supply its toilet flushing and garden irrigation from harvested rainwater.
The question has a number attached and a decision behind it: how large a
tank buys enough reliability to justify its cost, where reliability is the
fraction of demand the harvested supply meets over a typical year. The
answer must be a tank volume, good enough that a larger tank would not
change the recommendation and a smaller one would fail the reliability
target. The vignette runs the chapter's cycle end to end on a problem whose
data are a weather record and whose model is a stock and a flow.

\engterm{The question and the data.} The building has a roof catchment area
$A$, a daily water demand $D$, and access to a multi-year daily-rainfall
record for the site. The engineering question is the tank volume $V$ that
achieves, say, $95\%$ supply reliability, and the cost trade-off behind it:
tank cost rises roughly linearly with volume, while the mains-water cost
avoided rises with reliability and saturates. The data are the rainfall
record (the model's forcing), the roof area and a runoff coefficient (the
catchment efficiency), and the demand schedule.

\engterm{The candidate model.} The tank is a single stock, and the model is
the mass balance the compartment section named. On each day $t$ the stored
volume updates as
\[
V_t = \min\!\bigl(V_{\max},\ \max(0,\ V_{t-1} + I_t - D_t)\bigr),
\]
where $I_t = \eta A r_t$ is the inflow from that day's rainfall $r_t$
through the runoff coefficient $\eta$, $D_t$ is the demand, and the
$\min$ and $\max$ enforce the two hard limits: the tank cannot hold more
than its capacity $V_{\max}$ (excess overflows and is lost) nor less than
empty (a shortfall is met from the mains). The two clamps are the model's
nonlinearity, and they are what a naive average-balance calculation misses:
a tank sized on the annual average of inflow against demand ignores that a
full tank in a wet month cannot save water for a dry one past its capacity,
and an empty tank in a dry spell cannot borrow against future rain.
Figure~\ref{fig:vol02ch18:reservoir-reliability}(a) shows a stock
trajectory through a wet start and a dry spell that ends in a deficit.

\engterm{Identification and the dominant scale.} The model has no free
parameters to fit in the statistical sense; its inputs are measured
($\eta$ from a roof test or a standard table, $A$ from the drawings, $D$
from the demand schedule) and its forcing is the rainfall record. The
identification work is the dimensionless framing that tells the engineer
what tank size to even consider. The governing group is the ratio of tank
capacity to the demand accumulated over a characteristic dry spell,
$V_{\max}/(D\,t_{\text{dry}})$: a tank much larger than one dry spell's
demand rides through droughts, a tank much smaller cannot. The engineer
estimates $t_{\text{dry}}$ from the rainfall record (the longest common
dry run) before running anything, and the estimate brackets the tank sizes
worth simulating.

\engterm{Running the model and reading the reliability.} The model is run
over the full multi-year record for a range of candidate volumes, and for
each the reliability is the fraction of demand met from storage rather than
from the mains. The result is
figure~\ref{fig:vol02ch18:reservoir-reliability}(b): reliability rises
steeply with the first cubic metres of storage and then bends over into
diminishing returns, because past the knee the added capacity only helps in
the rare dry spells longer than the ones already covered. The knee is the
sizing answer. The engineer reads the smallest volume that meets the $95\%$
target and reports it, not the volume that maximises reliability, because
the last few percent cost far more storage than they return.

% Figure: rainwater tank sizing. (a) stock trajectory over a dry spell;
% (b) supply reliability as a saturating function of tank volume.
\begin{figure}[htbp]
\centering
\begin{tikzpicture}[font=\small]
% ---- panel (a): tank level over time ----
\begin{scope}
\draw[->, thick, draw=engblack!70] (0,0) -- (5.4,0)
  node[right, font=\scriptsize] {day};
\draw[->, thick, draw=engblack!70] (0,0) -- (0,4.0)
  node[above, font=\scriptsize] {stored volume};
% capacity line
\draw[dashed, draw=enggray] (0,3.4) -- (5.0,3.4);
\node[font=\scriptsize, color=enggray, anchor=east] at (5.0,3.6) {capacity $V$};
% level: fills in early rain, drops through dry spell, small refills
\draw[very thick, draw=engblue!80]
  (0,1.2) -- (0.6,3.4) -- (1.2,3.4) -- (2.6,1.0) -- (3.0,1.9) -- (4.4,0.15) -- (5.0,0.9);
% empty (deficit) marker
\filldraw[engred!80] (4.4,0.15) circle (0.06);
\node[font=\scriptsize, color=engred!80, anchor=north] at (4.4,0.05) {deficit};
\node[font=\scriptsize, anchor=north] at (2.6,-0.5) {(a) stock trajectory};
\end{scope}
% ---- panel (b): reliability vs tank size ----
\begin{scope}[xshift=7.0cm]
\draw[->, thick, draw=engblack!70] (0,0) -- (5.4,0)
  node[right, font=\scriptsize] {tank volume $V$};
\draw[->, thick, draw=engblack!70] (0,0) -- (0,4.0)
  node[above, font=\scriptsize] {reliability};
% 100% line
\draw[dashed, draw=enggray] (0,3.6) -- (5.0,3.6);
\node[font=\scriptsize, color=enggray, anchor=east] at (5.0,3.78) {$100\%$};
% saturating reliability curve R = 1 - exp(-V/V0)
\draw[very thick, draw=engblue!80] plot[domain=0:5,samples=60]
  (\x, {3.6*(1-exp(-\x/1.3))});
% knee marker
\filldraw[engorange!85] (1.9,{3.6*(1-exp(-1.9/1.3))}) circle (0.06);
\node[font=\scriptsize, color=engorange!85, anchor=north west] at (2.0,2.55)
  {diminishing returns};
% data points
\foreach \px in {0.4,0.9,1.5,2.3,3.3,4.4} {
  \pgfmathsetmacro\ry{3.6*(1-exp(-\px/1.3)) + 0.06*sin(\px*300)}
  \filldraw[engblack!75] (\px, {\ry}) circle (0.05);
}
\node[font=\scriptsize, anchor=north] at (2.6,-0.5) {(b) reliability curve};
\end{scope}
\end{tikzpicture}
\caption[Rainwater tank sizing as a stock-and-flow model]{A rainwater
harvesting tank modelled as a single stock. Panel (a): the stored volume
integrates the net flow of rainfall inflow minus demand outflow; through a
dry spell the stock falls, and a deficit occurs when it reaches empty
while demand continues. Panel (b): the fraction of demand met, the supply
reliability, rises with tank volume and saturates. The knee of the curve
is the sizing decision, the volume past which added capacity buys little
extra reliability, a diminishing-return trade-off the model exposes and a
rule of thumb cannot. The reliability curve is computed by running the
mass balance over a multi-year rainfall record, the daily-resolution
counterpart of the compartment model's time constant.}
\label{fig:vol02ch18:reservoir-reliability}
\end{figure}


\engterm{Validation, structural error, and the honest interval.} The
validation holds out one year of the record, sizes the tank on the rest,
and checks that the held-out year's reliability matches the prediction; a
tank sized on a wet decade and validated only on wet years would overstate
its reliability, the regime-extension failure applied to a weather record.
The structural error is named honestly: the daily model ignores sub-daily
intensity (a downpour that overflows a nearly full tank delivers less than
its daily total suggests), the runoff coefficient drifts with roof debris
and season, and the demand is treated as fixed when real demand rises in the
dry season exactly when supply falls, a correlation that makes the naive
model optimistic. The engineer reports the tank volume with a reliability
interval that spans the driest and wettest years in the record, states that
the recommendation assumes the historical rainfall distribution holds, and
names the regime (a shifting climate, a change of use that raises demand) in
which the sizing would need re-running. The deliverable is a volume with a
reliability band and a domain of validity, the same shape of report the
cooling cup produced, on a system whose forcing is a weather record and
whose model is one clamped stock.

\subsection{Applied vignette: predicting the range of an electric vehicle}

A second applied vignette runs the cycle on a quantity a driver reads off a
dashboard and an engineer is responsible for: the remaining range of an
electric vehicle. The number is a prediction with a decision behind it,
whether to press on to the next charger or divert, and a wrong prediction
strands the driver. The vignette exercises identification, the noise-model
care of section~18.4, validation against a held-out drive, and the honest
interval, on a system whose model couples a battery to a vehicle.

\engterm{The question and the data.} The engineering question is the
distance the vehicle can travel from its current charge before the battery
reaches its cutoff, good enough that the driver can trust it to within a
margin smaller than the spacing of chargers. The data are a controlled
discharge of the battery (current, terminal voltage, and accumulated charge
logged from full to empty at a known load) and a set of drive cycles
(speed against time, with the energy drawn logged) that stand in for the
range of real driving.

\engterm{The candidate model, in two coupled parts.} The model has a
battery part and a vehicle part. The battery is a lumped cell: an
open-circuit voltage that depends on the state of charge through a measured
curve, in series with an internal resistance that drops voltage under load,
\[
V_{\text{term}} = V_{\text{oc}}(\text{SOC}) - I R_{\text{int}},
\qquad
\text{SOC}(t) = \text{SOC}_0 - \frac{1}{Q_{\text{cap}}}\int_0^t I\,\dd t',
\]
where the state of charge is the running integral of current (coulomb
counting) against the cell's capacity $Q_{\text{cap}}$.
Figure~\ref{fig:vol02ch18:battery-soc}(a) shows the open-circuit-voltage
curve, whose long flat plateau is the structural fact that governs the
modelling choice: voltage is a sharp gauge of charge only near full and
empty and a useless one in the middle, so the estimator must integrate
current rather than read voltage, exactly the recursive-estimation move of
the previous vignette. The vehicle part converts charge into distance
through an energy model: the power the battery must supply is the sum of
rolling resistance, aerodynamic drag (which scales as speed squared, the
dimensional fact that makes highway range shorter than city range), and the
kinetic energy of accelerations less what regenerative braking returns.

\engterm{Identification, with the noise model watched.} The battery
parameters are identified from the controlled discharge: the open-circuit
curve from a slow discharge where the internal-resistance drop is
negligible, the internal resistance from the voltage drop when the load
steps, and the capacity from the total charge delivered.
Figure~\ref{fig:vol02ch18:battery-soc}(b) shows the terminal-voltage fit
tracking the measured discharge. The noise-model care of section~18.4
applies: the current sensor's error is roughly constant in absolute terms,
so coulomb counting accumulates a growing charge error over a long drive, a
drift the model must carry rather than ignore. The vehicle parameters
(rolling coefficient, drag area, mass) are identified from the drive-cycle
energy logs by regressing measured energy on the model's speed-dependent
terms, a linear identification once the model form is fixed.

% Figure: battery model. (a) open-circuit voltage versus state of charge,
% with a flat plateau; (b) terminal voltage under load, showing IR drop
% and the model fit against measured discharge points.
\begin{figure}[htbp]
\centering
\begin{tikzpicture}[font=\small]
% ---- panel (a): OCV vs SOC ----
\begin{scope}
\draw[->, thick, draw=engblack!70] (0,0) -- (5.4,0)
  node[right, font=\scriptsize] {state of charge};
\draw[->, thick, draw=engblack!70] (0,0) -- (0,4.0)
  node[above, font=\scriptsize] {open-circuit voltage};
% flat plateau OCV curve: steep rise, long plateau, steep drop at the ends
\draw[very thick, draw=engblue!80]
  (0.2,0.5) .. controls (0.9,2.9) and (1.4,3.1) .. (2.6,3.2)
  .. controls (3.8,3.3) and (4.3,3.4) .. (4.9,3.9);
% plateau annotation
\draw[dashed, draw=enggray] (1.4,3.15) -- (4.0,3.15);
\node[font=\scriptsize, color=enggray, anchor=south] at (2.7,3.15) {flat plateau};
% data points
\foreach \px/\py in {0.4/1.6, 0.9/2.95, 1.8/3.15, 2.8/3.22, 3.8/3.32, 4.6/3.7} {
  \filldraw[engblack!75] (\px,\py) circle (0.05);
}
\node[font=\scriptsize, anchor=north] at (2.6,-0.5) {(a) voltage plateau};
\node[font=\scriptsize, color=engred!80, anchor=west, text width=3.0cm]
  at (0.7,1.0) {steep near the ends};
\end{scope}
% ---- panel (b): terminal voltage under load, fit vs data ----
\begin{scope}[xshift=7.0cm]
\draw[->, thick, draw=engblack!70] (0,0) -- (5.4,0)
  node[right, font=\scriptsize] {discharge time};
\draw[->, thick, draw=engblack!70] (0,0) -- (0,4.0)
  node[above, font=\scriptsize] {terminal voltage};
% cutoff line
\draw[dashed, draw=engred!70] (0,0.7) -- (5.0,0.7);
\node[font=\scriptsize, color=engred!70, anchor=east] at (5.0,0.85) {cutoff};
% model fit curve: gentle sag then a knee down to cutoff
\draw[very thick, draw=engblue!80]
  (0.1,3.5) .. controls (1.8,3.0) and (3.2,2.7) .. (4.0,2.2)
  .. controls (4.4,1.7) and (4.6,1.0) .. (4.75,0.7);
% measured discharge points scattered on the fit
\foreach \px/\py in {0.3/3.45, 1.2/3.15, 2.2/2.85, 3.2/2.68, 3.9/2.28, 4.5/1.15} {
  \pgfmathsetmacro\jj{0.05*sin(\px*400)}
  \filldraw[engblack!80] (\px,{\py+\jj}) circle (0.055);
}
% range prediction marker
\draw[dashed, draw=engorange!85] (4.75,0) -- (4.75,0.7);
\node[font=\scriptsize, color=engorange!85, anchor=south west] at (4.0,0.15)
  {predicted range};
\node[font=\scriptsize, anchor=north] at (2.6,-0.5) {(b) discharge fit};
\end{scope}
\end{tikzpicture}
\caption[A battery model for range estimation]{A lumped battery model for
predicting electric-vehicle range. Panel (a): open-circuit voltage against
state of charge, with a long flat plateau that makes voltage a poor direct
gauge of remaining charge in the middle of the range and a sharp one only
near the extremes, the structural fact that forces the estimator to
integrate current (coulomb counting) rather than read voltage. Panel (b):
terminal voltage under a driving load, the sum of the open-circuit term
and an internal-resistance drop; the fitted model tracks the measured
discharge points, and the time at which the terminal voltage reaches the
cutoff sets the predicted range. The model is identified from a
controlled discharge and validated against a held-out drive cycle, the
chapter's cycle applied to a quantity a driver reads off the dashboard.}
\label{fig:vol02ch18:battery-soc}
\end{figure}


\engterm{Validation and the prediction.} Validation is a held-out drive
cycle the identification did not use: predict its energy draw from the
identified parameters and check it against the logged energy. A model fit to
gentle city cycles and validated only on them would mispredict a highway
run, because the drag term the city cycle barely exercised dominates at
speed, the regime-extension failure again. The range prediction inverts the
model: integrate the vehicle's energy demand along the anticipated route
until the battery reaches cutoff, and the accumulated distance is the range.
The engineer reports it with an interval from two sources the chapter has
named throughout. The parameter uncertainty propagates the intervals on
capacity, resistance, and drag through the integrated model. The structural
uncertainty is larger and is the honest part: the range depends on the
future drive, which the model cannot know, and on temperature (a cold
battery has less usable capacity and higher resistance), which a
single-temperature identification omits. The engineer reports the range as a
band conditioned on an assumed driving style and temperature, states that
the plateau of the voltage curve makes a voltage-only fuel gauge unreliable
in the middle of the range, and names the regimes (hard driving, cold
weather, a degraded aged battery) in which the prediction would need its
parameters re-identified. The deliverable is the dashboard number with its
honest band, produced by the same seven moves the cooling cup used, on a
system that couples an electrochemical stock to a mechanical energy balance.

\begin{mastery}
The two applied vignettes share a structure worth naming: each couples two
sub-models, a stock and its forcing in the tank, a battery and a vehicle in
the car, and the coupling is where the modelling risk concentrates.
\engterm{Model coupling} raises a hazard the single-model examples avoided:
an error in one sub-model can be hidden by a compensating error in the
other, so that the coupled system validates while both parts are wrong. A
battery capacity identified too high and a drag coefficient identified too
high can cancel on the drive cycle they were jointly fit to and diverge on
any other, the multi-parameter trade-off of section~18.10 crossing a
sub-model boundary. The defence is to identify and validate each sub-model
against its own data before coupling them: the battery against the
controlled discharge, the vehicle against the energy logs, so that each
part earns its parameters independently and the coupled validation tests the
coupling rather than compensating for two separate errors. The general rule,
which recurs wherever an engineer assembles a system model from component
models (a building from its rooms, a grid from its lines, a supply chain
from its stages), is that a coupled model is only as trustworthy as the
weakest independently validated sub-model, and a coupling validated only end
to end has not ruled out the compensating-error failure that a component
model, hiding inside, can carry undetected. The discipline of this chapter
applies to each sub-model and then, separately, to the coupling.
\end{mastery}

\section{Worked examples: four more moves, and a part that must not break}\pathtag{standard}

The running examples exercised the cycle on cooling, heating, queues,
networks, feedback, chemistry, and estimation. Four more short vignettes
finish the tour of moves the chapter has named, each on a system whose
mathematics the earlier examples did not reach: a field that collapses onto
one master curve, a nonlinearity that turns a flowing road into a jammed one,
a design that reads its answer off a phase lag, and a life prediction that
sums the damage of a random load. A final applied vignette runs the full
cycle at length on a part an engineer is answerable for, a component that must
survive its service life without cracking.

\subsection{A field that collapses: similarity and the diffusion profile}

The Buckingham collapses of section~18.3 reduced a scalar (drag, pressure
drop) to a curve in one group. The same move reduces a whole spatial field to
one master profile when the field has no intrinsic length or time scale of its
own. A thick slab, initially uniform at $T_0$, has its surface dropped to
$T_s$ at $t=0$; the interior cools as heat diffuses inward. The temperature
$T(x,t)$ depends on depth $x$ ($\mathsf{L}$), time $t$ ($\mathsf{T}$), and the
thermal diffusivity $\alpha$ ($\mathsf{L}^2\mathsf{T}^{-1}$). Three variables,
two primary dimensions, so the count gives one dimensionless group, and the
group is the \engterm{similarity variable}
\[
\eta = \frac{x}{2\sqrt{\alpha t}}.
\]
The slab is thick enough that no far boundary enters, so no length scale
competes with $\eta$, and dimensional homogeneity forces the normalised
temperature to depend on $x$ and $t$ only through their combination $\eta$:
$(T - T_s)/(T_0 - T_s) = \mathrm{erf}(\eta)$, the diffusion equation's
self-similar solution. Barenblatt's monograph
\cite{text:barenblatt1996} is the standard treatment of when a problem admits
such a similarity reduction and when a hidden scale spoils it.

The modelling payoff is the same as the sphere-drag collapse, raised from a
scalar to a field. Figure~\ref{fig:vol02ch18:similarity-collapse} shows it: raw
profiles measured at several times spread rightward as three separate curves,
and plotted against $\eta$ they fall onto one. An experimenter who did not know
the similarity would map the field on a grid of depths and times; the reduction
turns the grid into a single profile, measured once. The move also delivers a
scaling law for free, the $x\sim\sqrt{\alpha t}$ penetration depth that
section~18.13's Fourier-number vignette asserted: the collapse is the reason the
depth grows as the square root of time, because $\eta$ held fixed is exactly that
statement. The discipline the collapse rests on is a refusal, not an assertion:
the reduction is honest only while no length scale competes with $\eta$, and a
slab thin enough for its far face to matter breaks the similarity and demands the
full two-variable solution the collapse discarded.

% Figure: similarity collapse of a transient-diffusion profile. Temperature
% profiles at several times, each a spreading curve when plotted against
% depth, collapsing onto one master curve when plotted against the
% similarity variable depth over root time.
\begin{figure}[htbp]
\centering
\begin{tikzpicture}[font=\small]
% ===== Left panel: raw profiles at several times =====
\begin{scope}[xshift=0cm]
\draw[->, thick, draw=engblack!70] (0,0) -- (5.2, 0)
  node[right, font=\scriptsize] {depth $x$ (mm)};
\draw[->, thick, draw=engblack!70] (0,0) -- (0, 4.3)
  node[above, font=\scriptsize] {$(T-T_s)/(T_0-T_s)$};
\node[font=\small, anchor=south west] at (0.2, 4.1) {\textbf{(a) raw profiles}};
% Three erf-like profiles for increasing time: rise from 0 at x=0 to 1
% at large x, with a spreading front. Approximate erf(x/(2 sqrt(a t)))
% by a smooth curve; larger t spreads the front rightward.
\foreach \s/\col in {0.9/engred!85, 1.6/engblue!85, 2.5/engamber!90} {
  \draw[very thick, draw=\col, smooth, samples=60, domain=0:4.9]
    plot (\x, { 3.6*(1 - exp(-(\x/\s)*(\x/\s))) });
}
\node[anchor=west, font=\scriptsize] at (2.7, 1.3) {\textcolor{engred!85}{$t_1$}};
\node[anchor=west, font=\scriptsize] at (3.3, 2.0) {\textcolor{engblue!85}{$t_2$}};
\node[anchor=west, font=\scriptsize] at (3.9, 2.7) {\textcolor{engamber!90}{$t_3$}};
\end{scope}
% ===== Right panel: collapsed onto the similarity variable =====
\begin{scope}[xshift=7.2cm]
\draw[->, thick, draw=engblack!70] (0,0) -- (5.2, 0)
  node[right, font=\scriptsize] {$\eta = x/(2\sqrt{\alpha t})$};
\draw[->, thick, draw=engblack!70] (0,0) -- (0, 4.3)
  node[above, font=\scriptsize] {$(T-T_s)/(T_0-T_s)$};
\node[font=\small, anchor=south west] at (0.2, 4.1) {\textbf{(b) collapsed}};
% Single master curve: the erf profile in eta.
\draw[very thick, draw=engblue!80, smooth, samples=70, domain=0:4.9]
  plot (\x, { 3.6*(1 - exp(-(\x/1.5)*(\x/1.5))) });
% Points from all three times now falling on the one curve.
\foreach \col in {engred!85, engamber!90} {
  \foreach \eta in {0.5,1.0,1.6,2.4,3.4} {
    \pgfmathsetmacro\yy{3.6*(1 - exp(-(\eta/1.5)*(\eta/1.5)))}
    \filldraw[\col] (\eta, {\yy}) circle (0.045);
  }
}
\node[anchor=west, font=\scriptsize, color=engblack!70] at (1.7, 1.4)
  {one curve: $\mathrm{erf}(\eta)$};
\end{scope}
\end{tikzpicture}
\caption[Similarity collapse of a transient-diffusion profile]{Temperature
in a half-space whose surface is suddenly cooled, sampled at three times.
(a) Plotted against depth, the profiles spread rightward as the front
diffuses inward, three separate curves. (b) Plotted against the similarity
variable $\eta = x/(2\sqrt{\alpha t})$, the three collapse onto one master
curve, the error function. The single dimensionless group $\eta$ that
Buckingham counting produces from depth, diffusivity, and time carries the
whole solution: the field depends on $x$ and $t$ only through their
combination $x/\sqrt{\alpha t}$, so one profile measured once predicts the
field at every depth and time. The collapse is the same reduction the sphere
drag showed, here applied to a spreading field rather than a scalar force.}
\label{fig:vol02ch18:similarity-collapse}
\end{figure}


\subsection{A nonlinearity that jams: the traffic fundamental diagram}

The queue of section~18.9 found a congestion nonlinearity in $\rho/(1-\rho)$,
the wait diverging as utilisation approached one. A road carries the same shape
of nonlinearity in a different quantity, and the vignette earns its place
because the nonlinearity is one every commuter meets and few can name. The
modelling question is the road's capacity: how many vehicles per hour can a lane
carry, and what happens when demand exceeds it. The abstraction discipline names
the state variables first: the traffic density $\rho$ (vehicles per kilometre)
and the mean speed $v$. The load-bearing assumption is a relation between them,
and the simplest defensible one, Greenshields's, is linear: drivers slow as the
road fills, from a free-flow speed $v_f$ on an empty road to a standstill at the
jam density $\rho_{\text{jam}}$,
\[
v(\rho) = v_f\left(1 - \frac{\rho}{\rho_{\text{jam}}}\right).
\]
The flow, the vehicles per hour crossing a point, is the product $q = \rho v$,
and substituting the speed relation gives the \engterm{fundamental diagram}
\[
q(\rho) = v_f\,\rho\left(1 - \frac{\rho}{\rho_{\text{jam}}}\right),
\]
an inverted parabola. Figure~\ref{fig:vol02ch18:traffic-fundamental} plots both
relations.

The structural insight is available before any data, from the shape of the
parabola alone. Flow rises with density through the free-flowing regime, peaks at
a capacity $q_{\max} = v_f\rho_{\text{jam}}/4$ where the density is exactly half
the jam value, and falls back toward zero as congestion sets in. The peak is the
nonlinearity that a mean-density argument misses entirely: a road run just past
its capacity carries \emph{fewer} vehicles per hour at higher density than one
run just below it, which is the mechanism of a jam. Two densities give the same
flow, one free-flowing and one congested, and the difference between them is
whether a small disturbance clears or amplifies. A traffic engineer sizing a lane
reads the capacity off the peak and holds the operating density below it, exactly
as the clinic held its utilisation off the shoulder of the queueing blow-up. The
identification fits $v_f$ and $\rho_{\text{jam}}$ from paired speed-density
counts, and the structural error is the linear speed law itself, which real
traffic violates near capacity where the flow is unstable and stop-and-go waves
form; the residual there is the specification for the kinematic-wave model that a
later volume develops.

% Figure: the traffic fundamental diagram. Speed falls linearly with
% density (left); flow, the product of speed and density, is the resulting
% inverted parabola with a capacity maximum (right).
\begin{figure}[htbp]
\centering
\begin{tikzpicture}[font=\small]
% ===== Left panel: speed-density =====
\begin{scope}[xshift=0cm]
\draw[->, thick, draw=engblack!70] (0,0) -- (5.0, 0)
  node[right, font=\scriptsize] {density $\rho$ (veh/km)};
\draw[->, thick, draw=engblack!70] (0,0) -- (0, 4.3)
  node[above, font=\scriptsize] {speed $v$ (km/h)};
\node[font=\small, anchor=south west] at (0.2, 4.1) {\textbf{(a) speed-density}};
% Linear Greenshields: v = v_f (1 - rho/rho_jam). v_f at rho=0, zero at jam.
\draw[very thick, draw=engblue!80] (0, 3.8) -- (4.4, 0);
\node[anchor=east, font=\scriptsize] at (-0.05, 3.8) {$v_f$};
\draw[draw=engblack!70] (4.4,0) -- (4.4,-0.07);
\node[font=\scriptsize, anchor=north] at (4.4,-0.09) {$\rho_{\text{jam}}$};
\end{scope}
% ===== Right panel: flow-density (fundamental diagram) =====
\begin{scope}[xshift=7.2cm]
\draw[->, thick, draw=engblack!70] (0,0) -- (5.0, 0)
  node[right, font=\scriptsize] {density $\rho$ (veh/km)};
\draw[->, thick, draw=engblack!70] (0,0) -- (0, 4.3)
  node[above, font=\scriptsize] {flow $q=\rho v$ (veh/h)};
\node[font=\small, anchor=south west] at (0.2, 4.1) {\textbf{(b) fundamental diagram}};
% q = v_f rho (1 - rho/rho_jam): inverted parabola, peak at rho_jam/2.
\draw[very thick, draw=engred!85, smooth, samples=60, domain=0:4.4]
  plot (\x, { 3.6*4*(\x/4.4)*(1 - \x/4.4) });
% Capacity marker at half jam density.
\pgfmathsetmacro\xcap{2.2}
\draw[dashed, draw=engblack!55] (\xcap,0) -- (\xcap, 3.6);
\draw[dashed, draw=engblack!55] (0, 3.6) -- (\xcap, 3.6);
\node[font=\scriptsize, anchor=south] at (\xcap, 3.62) {capacity $q_{\max}$};
\node[font=\scriptsize, anchor=north] at (\xcap,-0.09) {$\rho_{\text{jam}}/2$};
\node[anchor=west, font=\scriptsize, color=engblack!70] at (0.4, 1.0) {free flow};
\node[anchor=east, font=\scriptsize, color=engblack!70] at (4.3, 1.0) {congested};
\end{scope}
\end{tikzpicture}
\caption[The traffic fundamental diagram]{The Greenshields model of a
single-lane road. (a) Speed falls linearly from the free-flow value $v_f$ at
empty road to zero at the jam density $\rho_{\text{jam}}$, the load-bearing
abstraction the model makes. (b) The flow, the product $q=\rho v$, is the
resulting inverted parabola: it rises with density through the free-flowing
regime, peaks at the capacity $q_{\max}$ where $\rho=\rho_{\text{jam}}/2$,
and falls back to zero as congestion stops the traffic. The single curve
carries the road's whole operating envelope, and the capacity peak is the
nonlinearity that a mean-density argument misses: a road run just past the
peak carries less traffic at higher density, the congestion collapse the
model exists to predict.}
\label{fig:vol02ch18:traffic-fundamental}
\end{figure}


\subsection{A design read off a phase lag: vibration isolation}

The feedback example of section~18.11 tuned a controller by moving eigenvalues.
A large class of design problems needs no controller at all, only the right
choice of a passive parameter, and the move is to read the design off the
system's steady-state response to a periodic input rather than off its transient.
A sensitive instrument sits on a bench that vibrates at a known frequency
$\omega$ (a nearby pump, a building's structural resonance). The engineering
question is how to mount the instrument so it feels as little of the bench motion
as possible. The candidate model is a single mass $m$ on a spring $k$ with damping
$c$, driven through its base by the bench displacement; the quantity the design
turns on is the \engterm{transmissibility}, the ratio of the instrument's
motion amplitude to the bench's. Solving the forced oscillator of Chapter~7 for
the steady state gives the transmissibility as a function of the frequency ratio
$r = \omega/\omega_0$, where $\omega_0 = \sqrt{k/m}$ is the mount's natural
frequency, and the structural fact the design rests on is the shape of that
function: transmissibility exceeds one (the mount \emph{amplifies}) for
$r < \sqrt{2}$, equals one at $r = \sqrt{2}$ regardless of damping, and falls
below one only for $r > \sqrt{2}$.

The design rule falls out of the dimensionless threshold. To isolate, the mount
must be soft enough that the disturbance frequency sits well above the mount's
natural frequency, $r \gg \sqrt{2}$, which means choosing $k$ and $m$ so that
$\omega_0 \ll \omega$. A stiff mount, with $\omega_0$ near $\omega$, sits on the
resonance and amplifies the very vibration it was meant to block, the failure a
naive ``make it rigid'' instinct produces. The damping choice carries its own
trade-off the model exposes: more damping tames the resonance the instrument
must pass through during start-up, but above $r=\sqrt{2}$ more damping
\emph{worsens} the isolation, because it couples the high-frequency bench motion
back through the dashpot. The modelling lesson is that the design number, the
mount stiffness, is read off a dimensionless frequency ratio and a phase
relationship, not off a transient, and a designer who reasons only about static
stiffness has modelled the wrong regime.

\begin{estimation}
\textbf{Question.} An instrument of mass $8\,\text{kg}$ must be isolated from a
bench vibrating at $30\,\text{Hz}$. Estimate the mount stiffness that puts the
system at a frequency ratio of $r = 4$, and estimate the resulting
transmissibility, before any detailed design.

\smallskip
\textit{Estimate, before any calculation.}

\smallskip
A frequency ratio $r = 4$ at a drive frequency of $30\,\text{Hz}$ needs a mount
natural frequency of $f_0 = 30/4 = 7.5\,\text{Hz}$, so $\omega_0 = 2\pi f_0
\approx 47\,\text{rad/s}$. The stiffness follows from $\omega_0^2 = k/m$:
$k = m\omega_0^2 = 8 \times 47^2 \approx 1.8\times 10^4\,\si{\newton\per\meter}$,
a soft mount that deflects about $mg/k \approx 8\times 9.8/1.8\times 10^4 \approx
4\,\text{mm}$ under the instrument's weight, a sane static sag. For light damping
the transmissibility above resonance is roughly $1/(r^2 - 1) = 1/15 \approx
0.07$, so the instrument feels about seven percent of the bench motion.

\smallskip
\textbf{Postmortem.} The dimensionless ratio did the work: fixing $r=4$ set the
natural frequency, which set the stiffness, which the static-sag check then
confirmed as buildable. A designer tempted to stiffen the mount for
robustness would raise $\omega_0$ toward the drive frequency, drop $r$ toward
one, and land on the resonance where transmissibility exceeds one, amplifying
the disturbance. The isolation lever is softness, not stiffness, which the
dimensionless form exposes and the static instinct hides.
\end{estimation}

\subsection{Applied vignette: will the part survive its service life?}

An engineer is responsible for a component that flexes every time the machine
cycles: a bracket in a compressor, a suspension arm on a vehicle, a weld in a
crane boom. The question has a number and a liability behind it: will the part
survive its design life of, say, ten million cycles without a fatigue crack, and
what margin does the prediction carry. A wrong answer on the unsafe side cracks a
part in service; a wrong answer on the safe side ships a part heavier and more
expensive than the duty needs. The vignette runs the chapter's cycle on a model
whose data are a material's fatigue tests and whose forcing is a measured load
history, and it earns its length because the failure the model predicts is one
the earlier examples only named.

\engterm{The question and the data.} The engineering question is the predicted
fatigue life $N_f$ under the part's actual service loading, good enough that a
safety factor can be applied with confidence. The data are two streams. The
first is a set of constant-amplitude fatigue tests on the material: coupons cycled
at several stress amplitudes until they crack, giving pairs of stress amplitude
$S$ and cycles-to-failure $N$. The second is a measured service-load history: a
strain gauge on a prototype logs the stress the part actually sees over a
representative operating period, a jagged signal of large and small cycles mixed
together.

\engterm{The candidate model.} The model has two parts, and each is a small
model in its own right. The material part is the \engterm{S-N curve}, the
empirical relation between stress amplitude and life, which plots as a straight
line on log-log axes over the finite-life range,
\[
S^{m} N = C
\qquad\text{equivalently}\qquad
\log N = \log C - m \log S,
\]
the power law of section~18.4 in fatigue's clothing, with the exponent $m$ and
the constant $C$ the two parameters to identify. The damage-accumulation part is
\engterm{Miner's rule}: a load history made of $n_i$ cycles at stress amplitude
$S_i$, each with allowable life $N_i$ read off the S-N line, accumulates a
linear damage
\[
D = \sum_i \frac{n_i}{N_i},
\]
and the part is predicted to fail when $D$ reaches unity.
Figure~\ref{fig:vol02ch18:fatigue-life} shows both parts, the S-N line reading
off the allowable lives and the Miner sum climbing toward the failure line.
Stephens and coauthors' fatigue text \cite{text:stephens-fatigue} is the working
reference for both the S-N fitting and the linear-damage rule, and Dowling
\cite{text:dowling-mechanical-behavior} treats the mean-stress corrections the
bare S-N line omits.

% Figure: fatigue-life prediction. An S-N curve (stress amplitude versus
% cycles to failure, log-log straight line) on the left; a Miner's-rule
% damage accumulation, summing the damage each stress block contributes,
% on the right.
\begin{figure}[htbp]
\centering
\begin{tikzpicture}[font=\small]
% ===== Left panel: S-N curve =====
\begin{scope}[xshift=0cm]
\draw[->, thick, draw=engblack!70] (0,0) -- (5.2, 0)
  node[right, font=\scriptsize] {$\log_{10} N$ (cycles)};
\draw[->, thick, draw=engblack!70] (0,0) -- (0, 4.3)
  node[above, font=\scriptsize] {$\log_{10} S$ (stress amp.)};
\node[font=\small, anchor=south west] at (0.2, 4.1) {\textbf{(a) S-N curve}};
% Basquin straight line on log-log: log S = log A - (1/m) log N.
\draw[very thick, draw=engblue!80] (0.3, 3.8) -- (4.8, 0.5);
% Mark two operating stress blocks and read their lives off the line.
\draw[dashed, draw=engred!70] (0, 3.0) -- (1.35, 3.0) -- (1.35, 0);
\node[font=\scriptsize, anchor=east, color=engred!85] at (-0.05, 3.0) {$S_1$};
\node[font=\scriptsize, anchor=north, color=engred!85] at (1.35,-0.09) {$N_1$};
\draw[dashed, draw=engamber!80] (0, 1.7) -- (3.05, 1.7) -- (3.05, 0);
\node[font=\scriptsize, anchor=east, color=engamber!90] at (-0.05, 1.7) {$S_2$};
\node[font=\scriptsize, anchor=north, color=engamber!90] at (3.05,-0.09) {$N_2$};
\node[anchor=west, font=\scriptsize, color=engblack!70] at (2.3, 3.3)
  {slope $-1/m$};
\end{scope}
% ===== Right panel: Miner damage accumulation =====
\begin{scope}[xshift=7.2cm]
\draw[->, thick, draw=engblack!70] (0,0) -- (5.2, 0)
  node[right, font=\scriptsize] {service (blocks)};
\draw[->, thick, draw=engblack!70] (0,0) -- (0, 4.3)
  node[above, font=\scriptsize] {cumulative damage $D$};
\node[font=\small, anchor=south west] at (0.2, 4.1) {\textbf{(b) Miner sum}};
% Failure line at D = 1.
\draw[dashed, draw=engblack!55] (0, 3.4) -- (5.0, 3.4);
\node[font=\scriptsize, anchor=south east, color=engblack!70] at (5.0, 3.4) {$D=1$: failure};
% Staircase of damage increments, each block adding n_i/N_i.
\draw[very thick, draw=engred!85]
  (0,0) -- (0.7,0.5) -- (1.4,0.85) -- (2.1,1.5) -- (2.8,1.9)
  -- (3.5,2.6) -- (4.2,3.0) -- (4.7,3.4);
\filldraw[engred!85] (4.7,3.4) circle (0.06);
\node[font=\scriptsize, anchor=west, color=engred!85] at (3.0, 1.6)
  {$D=\sum_i n_i/N_i$};
\end{scope}
\end{tikzpicture}
\caption[Fatigue-life prediction by the S-N curve and Miner's rule]{Predicting
when a cyclically loaded part fails. (a) The S-N curve is a power law, a
straight line on log-log axes (the Basquin form $S^m N = \text{const}$), whose
slope $-1/m$ is read from constant-amplitude fatigue tests: a lower stress
amplitude buys a longer life. Two operating stress blocks $S_1$ and $S_2$ read
off their allowable lives $N_1$ and $N_2$. (b) Miner's linear rule sums the
fractional damage $n_i/N_i$ each block of $n_i$ cycles contributes; the part is
predicted to fail when the accumulated damage $D$ reaches unity. The rule is
the model, the S-N line is its identified parameter set, and the scatter in
real fatigue life around $D=1$ is the structural error the report must carry.}
\label{fig:vol02ch18:fatigue-life}
\end{figure}


\engterm{The abstraction the model makes, and the one it hides.} The service
history is a jagged signal, not a clean set of blocks, so a counting rule must
first turn it into a histogram of stress amplitudes before Miner's sum can be
formed. The standard rule pairs each load reversal with its matching one to
extract closed hysteresis loops, and the modelling decision to use it is the
decision that the damage depends on the range of each closed cycle and not on the
order in which the cycles arrive. That order-independence is exactly what Miner's
linear rule assumes, and it is the model's largest structural error: real fatigue
damage depends on sequence, because a large cycle can open a crack that
subsequent small cycles then grow, so a history's damage is not truly the sum of
its parts. The abstraction is defensible because the sequence effect is modest for
many spectra and the linear rule is simple and auditable, but the modeller names
it rather than hiding it.

\engterm{Identification and the noise model.} The S-N parameters are identified by
regressing $\log N$ on $\log S$ across the coupon tests, a linear fit on the log-log
scale exactly as the chapter's power laws demanded. The noise-model care of
section~18.4 applies with force here: fatigue life scatters by a factor of several
at a fixed stress amplitude, and the scatter is multiplicative (a lognormal life),
so the fit must be on $\log N$, where the noise is additive, not on $N$. A modeller
who regressed $N$ directly would let the long-lived outliers dominate and return a
biased, over-optimistic line. The identification returns the exponent $m$, the
constant $C$, and, because the scatter is the whole point, the standard deviation
of $\log N$ about the line, which is the parameter the safety margin will use.

\engterm{Validation, structural error, and the honest interval.} Validation holds
out a subset of the coupon tests, fits the S-N line on the rest, and checks that
the held-out lives fall within the predicted scatter band; a line fit only to
high-stress short-life coupons and extrapolated to the low-stress long-life regime
is the regime-extension failure the chapter has named repeatedly, and it is
especially dangerous here because the low-stress regime is where the part actually
operates and where tests are most expensive to run. The prediction the engineer
reports is a life with an interval, and the interval has the two sources the
chapter has tracked throughout. The parameter uncertainty propagates the intervals
on $m$ and $C$ and the coupon scatter through the Miner sum. The structural
uncertainty is larger and is the honest part: the linear damage rule ignores load
sequence, the S-N line ignores mean stress unless corrected, and the service
history is one representative sample of a loading that varies from duty to duty.
The engineer reports the predicted life as a band, states that the median
prediction assumes the historical duty cycle and no unusual overloads, applies a
safety factor sized to the scatter rather than to the median alone, and names the
regimes (an overload event that shifts the mean stress, a corrosive environment
that accelerates crack growth, a duty cycle heavier than the logged one) in which
the prediction would need its coupon data re-acquired. The deliverable is a life
with a scatter band and a domain of validity, the same shape of report the cooling
cup produced, on a system whose failure is a crack rather than a wrong
temperature.

\begin{mastery}
The fatigue model rests on a linear damage rule that is known to be wrong in a
specific, characterisable way, and the case is worth naming because it is the
purest illustration of the chapter's ``all models are wrong, but some are
useful'' epigraph. Miner's rule assumes damage accumulates independently of
sequence and sums to unity at failure; measured failures scatter around a Miner
sum anywhere from about $0.3$ to $3$, because a high-low load sequence damages
differently from a low-high one and the interaction the linear rule omits is real.
The engineering response is not to abandon the rule but to carry its error
explicitly: the observed spread of the failure-point Miner sum is a directly
measured structural-error distribution, and the safety factor is sized to cover
it. The more accurate alternatives, a nonlinear damage rule that tracks the crack
state, or a fracture-mechanics model that integrates the Paris crack-growth law
\cite{text:hertzberg-fatigue} cycle by cycle, buy fidelity at the cost of
parameters (the initial flaw size, the crack-growth exponents) that are themselves
hard to identify and scatter widely. The modeller's judgement is the chapter's
recurring one: climb the fidelity ladder only as far as the decision requires, and
where the simple linear rule is retained, report its measured scatter as the honest
interval rather than pretending the unit-damage threshold is sharp. A fatigue life
quoted as a single number, without the factor-of-several scatter band the material
data show, has committed the chapter's central error, reporting a model's point
output as if it were a calibrated prediction.
\end{mastery}

\section{Failure: models that fooled their builders}\pathtag{core}

A model that has fooled its builder is the most dangerous artefact
this volume produces. The builder, persuaded by the model's
internal consistency, the closeness of the fit to past data, and
the elegance of the equations, applies the model to a regime where
its assumptions break and is surprised by a result the model has
no apparatus to refuse.

We treat this section at the recognition level rather than the
case-narrative level. The named registry holds the cases where
the registry's discipline applies; this section names categories
of failure rather than specific incidents, and points to where
the named treatments live in later volumes.

\subsection{Calibration that overfit}

A model with many parameters, fit to a few data points, can
match the data perfectly. The fit is the worst kind of
validation: the model has not learned the system, it has
memorised the data. A reader who fits a tenth-degree polynomial
to twelve points has built a model with two degrees of freedom
left over for residuals; the residuals will be small by
construction, the prediction outside the data range will be
wild. The failure mode is overfitting; the discipline is the
train-validate split, restated above.

The engineering form of the failure occurs when a model is
calibrated to historical data of a system, used to predict the
system's behaviour in a regime the historical data did not
cover, and trusted in that regime because it agreed with the
historical record. The agreement was guaranteed by construction;
the trust was unearned. The next regime exposes the model's
extrapolation as a guess, and the engineer is surprised by the
disagreement.

The diagnostic is structural. A model whose parameters
outnumber the regimes the data span has not been validated for
extrapolation. The engineer who proposes to use the model
outside its calibration regime has the burden of additional
data, a structural argument, or both.

% Figure: overfit demonstration. Quadratic ground truth, twelve data
% points, two fits: a quadratic and a degree-10 polynomial that
% interpolates the points.
\begin{figure}[htbp]
\centering
\begin{tikzpicture}[font=\small]
% Axes
\draw[->, thick, draw=engblack!70] (-0.5, 0) -- (9.0, 0) node[right, font=\small] {$x$};
\draw[->, thick, draw=engblack!70] (0, -1.0) -- (0, 4.5) node[above, font=\small] {$y$};
% Ticks
\foreach \xt/\xv in {1/0.2, 3/0.4, 5/0.6, 7/0.8} {
  \draw[draw=engblack!70] (\xt, 0) -- (\xt, -0.08);
  \node[font=\scriptsize, anchor=north] at (\xt, -0.1) {\xv};
}
\foreach \yt/\yv in {1/1, 2/2, 3/3, 4/4} {
  \draw[draw=engblack!70] (0, \yt) -- (-0.08, \yt);
  \node[font=\scriptsize, anchor=east] at (-0.1, \yt) {\yv};
}

% Ground truth: y = (x/3)^2 + 0.5; quadratic on x in [0, 9]
\draw[thick, draw=engblue!80, dashed]
  plot[smooth, samples=80, domain=0:9] (\x, {(\x/3)^2 + 0.5});

% Quadratic fit (close to ground truth)
\draw[thick, draw=engblue!90]
  plot[smooth, samples=80, domain=0:9] (\x, {(\x/3)^2 + 0.55});

% Degree-10 polynomial: interpolates 12 points but wild between
\draw[thick, draw=engred!85]
  plot[smooth cycle=false, samples=200, domain=0:9, smooth]
    (\x, {(\x/3)^2 + 0.5 + 0.6 * sin(150 * \x) * exp(-((\x-4.5)^2)/9)});

% Twelve data points (small noise around ground truth)
\foreach \xt/\yval in {0.3/0.6, 1.0/0.65, 1.8/0.95, 2.5/1.3, 3.3/1.8,
  4.0/2.1, 4.8/2.7, 5.5/3.2, 6.3/3.95, 7.0/4.05, 7.7/4.55, 8.5/4.9} {
  \filldraw[engblack!90] (\xt, \yval) circle (0.07);
}

% Five held-out points; show degree-10 misses them badly
\foreach \xt/\yval in {0.7/0.65, 2.2/1.2, 5.0/2.55, 6.8/3.85, 7.4/4.25} {
  \draw[engblack!90, thick] (\xt - 0.1, \yval - 0.1) -- (\xt + 0.1, \yval + 0.1);
  \draw[engblack!90, thick] (\xt - 0.1, \yval + 0.1) -- (\xt + 0.1, \yval - 0.1);
}

% Legend
\node[anchor=west, font=\scriptsize] at (0.3, 4.3) {\textcolor{engblue!80}{(dashed)} ground truth};
\node[anchor=west, font=\scriptsize] at (0.3, 4.0) {\textcolor{engblue!90}{(blue solid)} quadratic fit (fit and predicts)};
\node[anchor=west, font=\scriptsize] at (0.3, 3.7) {\textcolor{engred!85}{(red solid)} degree-10 fit (fits, fails to predict)};
\node[anchor=west, font=\scriptsize] at (0.3, 3.4) {$\bullet$ training data, $\times$ held-out};
\end{tikzpicture}
\caption[Overfit: a polynomial that fits the data but fails on
prediction]{Twelve training points (filled circles) drawn from a
quadratic ground truth (dashed) with small Gaussian noise. The
quadratic fit (solid blue) recovers the ground truth almost
exactly. The degree-10 polynomial fit (solid red) interpolates the
training points so closely that its training residuals are
near-zero; on the held-out points (crosses) it produces excursions
several times the noise level. The mechanism is structural: a
ten-parameter model fit to twelve points has two residual degrees
of freedom and almost any wiggle between the points is permitted
by the fit. The defence is the train-validate split, restated in
section~18.3 and exercised in Exercise 18.16 (overfit, by
construction).}
\label{fig:vol02:ch18:overfit}
\end{figure}


Figure~\ref{fig:vol02:ch18:overfit} shows the failure mode at the
mechanism level. A quadratic ground truth produces twelve noisy
points; a quadratic fit recovers the ground truth, a degree-10
polynomial interpolates the points and explodes between them. The
training residuals of the degree-10 fit are smaller than the noise
(a sign that the fit has memorised noise as signal); the held-out
residuals are larger by an order of magnitude. The companion script
\texttt{overfit\_demo.py} runs the demonstration end-to-end; the
ratio it prints is the diagnostic the engineer can compute on any
candidate model. Hastie, Tibshirani, and Friedman
\cite{text:v2ch18-hastie-statistical-learning-2009} treat the
bias-variance decomposition behind the picture; the engineer takes
from it the working rule that complexity costs predictive
performance unless the data span the regime that the complexity is
needed to model.

% Figure: the bias-variance decomposition of test error as a function of
% model complexity, with the total error U-curve and its two components.
\begin{figure}[htbp]
\centering
\begin{tikzpicture}[font=\small]
\draw[->, thick, draw=engblack!70] (0,0) -- (8.0,0)
  node[right, font=\scriptsize] {model complexity (parameters $p$)};
\draw[->, thick, draw=engblack!70] (0,0) -- (0,4.4)
  node[above, font=\scriptsize] {expected squared error};
% bias^2: high at low complexity, falls
\draw[very thick, draw=engblue!80, smooth, samples=90, domain=0.3:7.6]
  plot (\x, { 3.6*exp(-0.65*\x) + 0.15 });
\node[font=\scriptsize, text=engblue!80, anchor=west] at (2.0,1.35)
  {bias$^2$};
% variance: low at low complexity, rises
\draw[very thick, draw=engred!85, smooth, samples=90, domain=0.3:7.6]
  plot (\x, { 0.12*exp(0.52*\x) + 0.1 });
\node[font=\scriptsize, text=engred!85, anchor=east] at (7.4,3.2)
  {variance};
% irreducible noise floor
\draw[thick, draw=enggray, dashed] (0,0.6) -- (7.6,0.6)
  node[right, font=\scriptsize, text=enggray] {noise floor};
% total = bias^2 + variance + noise: U-shape
\draw[very thick, draw=engblack!80, smooth, samples=120, domain=0.3:7.6]
  plot (\x, { 3.6*exp(-0.65*\x) + 0.15 + 0.12*exp(0.52*\x) + 0.1 + 0.6 });
\node[font=\scriptsize, text=engblack!80, anchor=south west] at (0.4,4.0)
  {total test error};
% optimum marker (rough minimum)
\draw[dashed, draw=engblack!50] (3.6,0) -- (3.6,1.55);
\fill[engblack!80] (3.6,1.55) circle (1.8pt);
\node[font=\scriptsize, anchor=south] at (3.6,1.6) {$p^\star$};
\node[font=\scriptsize, anchor=north, text=engblack!60] at (1.6,-0.35)
  {underfit};
\node[font=\scriptsize, anchor=north, text=engblack!60] at (6.0,-0.35)
  {overfit};
\end{tikzpicture}
\caption[The bias-variance decomposition]{The expected squared error of a
model's predictions on new data, decomposed into three parts as model
complexity grows. Bias squared (blue) is large for a too-simple model that
cannot represent the system and falls as complexity rises. Variance (red)
is small for a simple model and grows as added parameters let the fit chase
noise. The irreducible noise floor (grey dashed) is the part no model can
remove. Their sum (black) is U-shaped, with its minimum at the complexity
$p^\star$ that balances under-representation against over-flexibility. The
left branch is the underfit regime, where the model selection of
section~\ref{fig:vol02:ch18:model-selection} sits below the optimum; the
right branch is the overfit regime of figure~\ref{fig:vol02:ch18:overfit}.
The decomposition is why the simplest model that answers the question is
not merely a matter of taste: complexity beyond $p^\star$ raises the error
the model will incur on data it has not seen.}
\label{fig:vol02:ch18:bias-variance}
\end{figure}


Figure~\ref{fig:vol02:ch18:bias-variance} resolves the same picture into
its components. The error a model incurs on data it has not seen splits
into a bias term, large when the model is too simple to represent the
system, a variance term, large when the model is flexible enough to chase
the noise, and an irreducible noise floor no model removes. The
degree-10 polynomial of figure~\ref{fig:vol02:ch18:overfit} sits far out
on the variance branch: its bias is near zero (it can represent any
twelve-point pattern) but its variance is enormous (a different noise
draw produces a wildly different fit). The quadratic sits near the
minimum, where the two terms balance. The decomposition is why the
simplest model that answers the question is not a stylistic preference
but a quantitative optimum: complexity past $p^\star$ trades a small bias
reduction for a large variance increase, and the prediction error rises.

\subsection{Garbage in: the model that inherited its data's flaws}

The failure modes above assume good data and a wrong model. The
complementary failure is a defensible model fed data that quietly
violate the model's assumptions, and it is the more common failure in
practice because data quality is invisible in the fit's residuals when
the flaw is systematic. A least-squares fit reports the variance of the
residuals about the fitted curve; it cannot report that the
measurements were all biased by a miscalibrated instrument, that a
sensor saturated above a threshold the data never crossed in the
training window, or that a logging system dropped samples
preferentially during the events of interest. Each of these is a
violation the model's noise term does not see, and each produces a fit
that looks clean and predicts badly.

The mechanism is worth stating sharply because it resists the chapter's
other defences. Cross-validation does not catch a bias shared by all the
data, because every fold inherits the same bias; the held-out residuals
are as clean as the training residuals, and the model passes validation
while being wrong. The defence is not statistical but metrological: the
instrument is calibrated against a traceable standard (Volume I Chapter
3), the data are audited for missingness and saturation before fitting,
and the provenance of every datum is recorded. The modelling discipline
borrows from the measurement discipline here, and the borrowing is the
point. A model is only as trustworthy as the data's traceability, and a
modeller who accepts data without provenance has built a model whose
domain of validity is unknown regardless of how well it fits.

A concrete form: a sensor with a range ceiling that the training data
never reached. A flow meter calibrated and logged across a range that
happened to stay below its $10\,\text{L/s}$ ceiling produces clean data
and a clean fit. A model fit to that data predicts confidently at
$15\,\text{L/s}$, a regime the sensor cannot even report; the prediction
is a pure extrapolation dressed as a fit, and the engineer who reads the
training residuals sees no warning. The diagnostic is to plot the data's
coverage in every input dimension and to mark the prediction's location
relative to that coverage. A prediction outside the data's convex hull
is an extrapolation, whatever the residuals say, and the discipline is
to label it as one.

\begin{mastery}
A subtler instance arises when the data are generated by a process that
responds to the model. A model used to allocate inspection effort, fit to
historical defect data, directs inspectors to the locations the model
predicts; the locations the model rates low are inspected less, so fewer
defects are recorded there, so the next fit rates them lower still. The
data the model is retrained on are no longer an independent sample of
the world; they are a sample shaped by the model's own past predictions.
The fit's residuals shrink while the model's correspondence to the true
defect distribution degrades. This is the closed-loop version of garbage
in, and it is the mechanism behind a class of failures in
predictive policing, credit scoring, and recommendation systems, where
the model's outputs alter the distribution the model is later evaluated
against. Goodhart's observation, that a measure used as a target ceases
to be a good measure, is the operational summary. The defence is to hold
out a randomised audit sample that is collected independently of the
model's recommendations, so that at least one stream of data remains an
honest sample of the world; without it, no amount of cross-validation on
the model's own loop can detect the drift. Volume X returns to this
mechanism where the systems are large enough that the feedback dominates.
\end{mastery}

\subsection{Validated within the regime, applied outside}

A model can be honestly validated within a regime and then
applied outside the regime as if the validation extended.
Newton's law of cooling validated at modest temperature
excess does not validate the model at high temperature
excess where radiative cooling dominates. A linear stress-
strain relation validated below the yield stress does not
validate the model above yield. A logistic growth model
validated in the exponential phase does not validate the
model in the death phase.

The failure mode is the regime-extension fallacy. The
discipline is to state the validation regime explicitly, to
refuse the model outside the regime, and to acquire data in
the new regime before extending the model. The discipline is
mechanical; the failure persists because the engineer who has
seen the model work in one regime acquires confidence the
model has not earned for the next.

\subsection{The risk model that did not include its mechanism}

A risk model that quantifies the probability of an outcome
without including the mechanism by which the outcome arises
will report low probabilities for outcomes whose mechanism the
model does not see. The 2007-09 financial crisis produced
several discussions in which models that fit historical
default rates were used to price instruments whose default
correlations the historical data did not include; when the
correlated regime arrived, the models reported probabilities
that were orders of magnitude smaller than the realised
frequencies. The episode is a recurring case in the
engineering-modelling literature.

We name the category, not the case. The discipline that the
chapter has installed says: a model whose parameters are
identified from data of a regime cannot quantify the
probability of an outcome outside the regime. The model can
report a probability, but the probability is a working
number that has not been validated against the regime in
which it is being used. The reader who uses the number as if
it were calibrated has confused the model's within-regime
output with an across-regime prediction; the failure mode is
the same as the regime-extension fallacy, applied to a
probability rather than a magnitude.

\subsection{The simulation that became the system}

A computational model of a system, run at high resolution,
produces predictions of the system's behaviour. After many
iterations of refinement, the simulation matches the
historical data closely. The engineer, working within the
simulation, begins to reason about the system through the
simulation rather than against it. Decisions are made on the
basis of simulation runs; the simulation's outputs replace
direct measurement of the system. The model has become the
system, in the engineer's working reasoning.

The failure mode is described informally as ``believing the
simulation.'' Its mechanism is the gradual replacement of the
discipline of validating the model against the system with the
discipline of refining the model against the data the model
was previously fit to. The simulation's residuals get smaller;
its predictions on new regimes get worse, because the
refinement chases features that the data already showed and
ignores features that the data did not.

The defence is procedural and was named in section 18.3.
Validation against held-out data, repeated at regular intervals,
exposes the divergence between the simulation and the system.
A simulation whose validation residuals are growing over time
is a simulation that is being refined toward its training data
rather than toward the system, and the engineer's response is to
acquire new data and to question the refinements that have
narrowed the training residuals while widening the validation
ones.

\subsection{Two short-form cases at the recognition level}

The chapter avoids the named-cases registry in this volume (the
registry's discipline begins in Volume III, where the apparatus
applies to load-bearing engineering systems rather than to the
modelling apparatus itself). At the recognition level, two episodes
deserve naming because they are the textbook examples of two of the
four failure modes above. We name them once, point to where a
serious treatment of each can be found, and continue.

\engterm{Long-Term Capital Management, 1998}. A hedge fund founded
on a model of bond-spread mean reversion, calibrated against several
years of data from a regime in which sovereign-bond markets had
behaved smoothly, applied the model to instruments whose default
correlations the historical data did not include. When the 1998
Russian sovereign default arrived, the correlated regime arrived
with it; the model's reported probability of the realised loss was
many orders of magnitude below the realised frequency. The episode
is the textbook case of the regime-extension fallacy applied to a
probability. Roger Lowenstein's \emph{When Genius Failed} (2000)
and Donald MacKenzie's \emph{An Engine, Not a Camera} (2006) are the
standard non-technical and technical treatments respectively; the
engineering-modelling reader's takeaway is that a model's reported
probability is a within-regime output, not a calibrated prediction
across regimes.

\engterm{The Vasa, 1628}. The Swedish flagship \emph{Vasa} sank on
its maiden voyage in Stockholm harbour after sailing less than a
nautical mile, the result of a stability calculation that the
shipwrights had not performed and a stability test that the
admiral cut short under political pressure. The model existed
(the metacentric-height analysis was understood in principle); the
discipline of performing the test and acting on its result was
absent. The case is a textbook instance of the simulation that
became the system: refinements of the upper deck's armament were
pursued under royal pressure, refinements of the stability
calculation were not, and the model of an in-spec ship replaced the
discipline of measuring an in-spec ship. Richard Hocker's
\emph{Vasa: A Swedish Warship} (2011) is the modern primary-source
synthesis; the engineering takeaway is that the absence of a
modelling step is itself a modelling decision, and an undocumented
one is the most consequential variety.

We name these once, at the recognition level, and rely on the
chapter's apparatus for the structural argument. The reader who
wants the case-study depth that the named-cases registry provides
should consult the registry entries that Volumes III onward
maintain; the registry is the chapter's downstream apparatus, not
its own.

\subsection{The principle}

\begin{principle}[A model is a hypothesis until it is validated against
data the identification did not use]
A model fit to data is a description of the data. A model whose
predictions on data the identification did not use have been
compared, quantitatively, to those data is a validated model. The
two are different artefacts. A modelling report that conflates
them, by reporting fit residuals as validation residuals or by
reporting parameter intervals as predictive intervals, has not
yet entered the discipline. The engineering action is to set
aside data before identification, to compare predictions to those
data after identification, and to inflate the prediction interval
by a factor that reflects the structural uncertainty observed in
the validation.
\end{principle}

The principle is procedural rather than analytical. Its content
is not new; the structural defence the engineer practises is the
same defence Volume I named for measurements (record the
estimate before the calculation), the same defence Volume II
Chapter 14 named for hypothesis tests (state the test before
seeing the data), and the same defence the next ten volumes will
name for designs, prototypes, and regulations. The defence is
the discipline of producing a record before producing a result,
because human reasoning is too good at aligning records with
results retroactively.

\section{Historical sketch: how engineers learned to model}\pathtag{standard}

The modelling discipline this chapter codifies did not arrive whole. It
accreted over two centuries, as engineers and applied mathematicians
discovered, often through failure, which moves were load-bearing and
which were ornament. The sketch is partial and the dates carry their
own half-lives, but the arc is worth tracing because it shows the
discipline as a hard-won settlement rather than a set of arbitrary
rules.

\subsection{Dimensional reasoning becomes a method}

Fourier, in the 1822 \emph{Th\'eorie analytique de la chaleur}, was
explicit that physical equations must be dimensionally homogeneous, and
he used the requirement to check his heat-conduction results. The idea
that dimensions alone constrain the form of physical law sat as folk
practice through the nineteenth century, used by Rayleigh and others
without a general theorem. Buckingham's 1914 paper
\cite{paper:buckingham1914} supplied the theorem, naming the count of
dimensionless groups and the procedure to form them, and the method
became standard equipment in fluid mechanics and heat transfer within a
generation. Bridgman's 1922 lectures \cite{text:bridgman1922}
turned it into a coherent philosophy of measurement, and Barenblatt's
later work \cite{text:barenblatt1996} connected it to scaling,
self-similarity, and the intermediate asymptotics that govern how a
model behaves between its limiting regimes. The lesson the period
taught is that the cheapest structural information in any model comes
before the dynamics, from the dimensions of the variables alone.

\subsection{Rate and compartment models spread across domains}

The first-order rate law, $\dot{x} = -kx$, appeared independently in
radioactive decay (Rutherford and Soddy, 1902), in chemical kinetics
(the nineteenth-century law of mass action), and in Newton's much older
law of cooling. The recognition that these were the same mathematics
applied to different conserved quantities was itself a modelling
advance. Lotka \cite{hist:v6c08-lotka-1925} and Volterra
\cite{hist:v6c08-volterra-1926} extended the structure to coupled
nonlinear compartments for predator and prey in the 1920s, and Kermack
and McKendrick \cite{paper:v6c08-kermack-mckendrick-1927} did the same
for epidemics in 1927, deriving the threshold theorem that a
dimensionless reproduction number controls. By mid-century the
compartment model was the shared language of pharmacokinetics, reactor
engineering, ecology, and physiology, the same equations rewritten for
each field's conserved stock.

\subsection{Systems thinking and the honesty of uncertainty}

The control-theory and operations-research traditions of the
mid-twentieth century made feedback, delay, and stability the explicit
objects of modelling rather than incidental features, and Meadows's
\emph{Thinking in Systems} \cite{gen:meadows2008} later distilled the
stocks-flows-feedback discipline into a form usable without heavy
notation. In parallel, the statistical tradition made the honest
reporting of uncertainty a professional norm. Box, whose aphorism opens
this chapter, argued in his 1979 essay
\cite{paper:v6c08-box-1979} and his response-surface monograph
\cite{text:v2ch18-box-empirical-1987} that the modeller's job is to
find the model that is useful despite being wrong, and to say in what
respect it is wrong. Cobelli and DiStefano's 1980 review
\cite{paper:v2ch18-cobelli-distefano-1980} settled the language of
identifiability that the chapter uses. Saltelli's sensitivity-analysis
program \cite{text:v2ch18-saltelli-sensitivity-2008} and the
statistical-learning synthesis of Hastie, Tibshirani, and Friedman
\cite{text:v2ch18-hastie-statistical-learning-2009} brought the
validation, cross-validation, and bias-variance apparatus into the form
the working engineer now reaches for. The arc ends where the chapter
stands: a model is a hypothesis, its uncertainty is part of its report,
and its validation is against data it was not fit to.

\section{Modelling as a moral discipline}\pathtag{core}

We close Volume II where the volume began: with the claim that
mathematics is the engineer's working notation. The reader who
finishes this chapter has practised the apparatus of fifteen
hundred pages on a question of their own choice and has produced
a model that is reportable: data, fit, validation, structural
account, sensitivity, prediction interval, domain of validity,
omitted-variable list. The deliverable is the credential that
admits the reader to Volume III, where the language is taken for
granted.

The credential has professional content. An engineer who reports
a model with its interval and its domain of validity does not
make decisions on a number that the next engineer cannot trace.
An engineer who states the omissions of the model lets a
subsequent engineer revise the model when new data or new
questions arrive. An engineer who validates against held-out
data refuses to confuse fit with prediction. The three habits,
practised across projects, are the discipline that turns a model
from a private artefact of one engineer into a working object
that an organisation can use.

The discipline scales. A profession that models honestly produces
a record of which decisions were made on calibrated predictions,
which were made on uncalibrated ones, and which were made on
working numbers that were not yet ready to leave the desk. A
profession that models dishonestly loses the record; the
distinction between calibrated and uncalibrated decisions is what
gets erased first, and after the erasure the profession can no
longer say which of its predictions are trustworthy. The reader
who closes this volume with the discipline installed has done
the second-hardest part of becoming an engineer in the sense the
book uses the word; the hardest part, the discipline of
estimation, was installed in Volume I.

The remaining ten volumes will exercise the discipline. Volume III
applies it to forces; Volume IV to flows; Volume V to materials and
chemistry; Volume VI to biology and medicine; Volume VII to
information and computation; Volume VIII to machines and control;
Volume IX to systems and infrastructures; Volume X to failures and
their lessons; Volume XI to design and synthesis; Volume XII to the
civilization the engineered systems compose. Each volume will ask
the reader to model a quantity in the volume's domain. The seven
moves are the same: name the question, choose the abstraction,
identify the parameters, validate against held-out data, state the
structural error, propagate the sensitivity, look back. The reader
who has practised the seven moves on a quantity from the Volume I
notebook can practise them, with the same confidence, on any
quantity the rest of the work asks them to model.

\section{Project: from Volume I notebook to defended model}\pathtag{core}

The chapter's project is the reader's first act of independent
modelling. The deliverable is a 5-page modelling write-up that
takes one quantity from the Volume I capstone notebook and turns
it into a model the reader can defend.

\begin{project}[Modelling a quantity from the Volume I notebook]
\textbf{Track}: analysis only. \textbf{Hazard class}: none.
\textbf{Tools}: paper, pencil, calculator, the reader's own
Volume I measurement notebook. \textbf{Time}: ten to twenty
hours, distributed across one to two weeks. \textbf{Deliverable}:
a 5-page write-up plus a 1500-word reflection.

\smallskip
\textbf{Selection}. The reader chooses one quantity from their
Volume I capstone notebook. The quantity must satisfy three
criteria. It must have at least ten data points (so an
identification is possible). It must have an inferable
mechanism (so a candidate model exists). It must matter to a
question the reader can state in one sentence (so the model
has a purpose). Examples: the cooling-cup notebook of Chapter 4;
the small-circuit current measurements of Chapter 5; the
pendulum periods of Chapter 4; any of the Chapter 8 statistical
data sets the reader collected.

\smallskip
\textbf{The 5-page write-up}. The write-up has seven sections,
in this order, with the indicated approximate length.

\engterm{Section 1: question}, half a page. The engineering
question the model must answer, stated explicitly. The units of
the answer. The precision the answer requires. The decision the
answer supports.

\engterm{Section 2: data}, half a page. The data set used, with
its source, its conditions, its uncertainty (from the Volume I
calibration), and any data points excluded with reasons.

\engterm{Section 3: candidate model}, one page. The model's
equation, its parameters, its assumptions, and the
abstraction discipline applied: which features the model
keeps, which it throws away, which dimensionless groups
characterise its regime.

\engterm{Section 4: identification}, one page. The method
(linear regression, nonlinear least squares, maximum
likelihood, or other), the point estimates, the intervals, the
identifiability check.

\engterm{Section 5: validation}, one page. The held-out data
or cross-validation procedure, the validation residuals, the
diagnostic plots, the structural-error account from the
residual structure.

\engterm{Section 6: prediction}, half a page. The prediction
relevant to the engineering question, with its interval,
including both parameter and structural uncertainty.

\engterm{Section 7: domain and omissions}, half a page. The
regime in which the model is validated. The omitted variables
and the regime in which each becomes consequential. The
specific conditions under which the model would not be applied.

\smallskip
\textbf{The 1500-word reflection}. The reflection addresses five
questions. Which modelling decisions were the most consequential.
Which decisions, if reversed, would change the engineering
conclusion. Which features of the data forced which features of
the model. What the validation residuals revealed that the fit
residuals did not. What the difference is, for this model,
between a model the reader can defend and a model the reader
cannot.

\smallskip
\textbf{Defensibility}. The deliverable is judged against the
question: would the model survive a hostile review? A hostile
reviewer asks: what data are missing? What omitted variables
might dominate? What would change the conclusion? What is the
domain of validity, and how was it established? A model that
answers these questions in writing has earned its credential.

\smallskip
\textbf{Phases and milestones.} The ten-to-twenty hour budget is
divided into four phases, each with a milestone the reader can
verify before moving on. The phasing matters because the most
common project failure is to write the report before validating the
model, at which point the report's intervals are wrong and the
revision is more costly than the original work.

\engterm{Phase 1: question and data audit} (2--3 hours). The reader
selects the quantity, writes Section 1 of the write-up (question)
and Section 2 (data) in draft, and produces a one-paragraph audit
that states whether the data set passes the three selection
criteria. \textit{Milestone}: a written one-sentence engineering
question, a written data description, and a binary go/no-go on the
selection. If the data fail any criterion, the reader returns to
selection rather than forcing the model on a data set that does
not support one.

\engterm{Phase 2: abstraction and identification} (3--5 hours). The
reader drafts the candidate model (Section 3), names the dominant
scales and dimensionless groups, and performs the identification
(Section 4) using one of the chapter's companion scripts as a
template. \textit{Milestone}: a point estimate of each parameter
with a 95\% interval, an explicit identifiability check (the
Jacobian condition number from \texttt{logistic\_fit.py} or the
analogous check from \texttt{cooling\_fit.py}), and a one-paragraph
defence of why the chosen model is the simplest model that answers
the question.

\engterm{Phase 3: validation and structural error} (3--5 hours).
The reader holds out a fraction of the data (Section 5), refits on
the rest, computes the validation residuals, runs $k$-fold
cross-validation via \texttt{cross\_validation.py}, and produces
the residual-structure plot that figure 18.3 exemplifies.
\textit{Milestone}: a validation residual plot, an inflation factor
from cross-validation, and a written one-paragraph statement of
the structural error the residuals expose (or of why they expose
none).

\engterm{Phase 4: prediction and look-back} (2--3 hours). The reader
computes the prediction relevant to the engineering question with
its honest interval (Section 6, parameter plus structural
uncertainty propagated), states the domain of validity and the
omitted variables (Section 7), and writes the 1500-word reflection.
\textit{Milestone}: a single number with an interval, ready to be
handed to a downstream user; a written domain-of-validity statement;
and a reflection that answers the five questions above.

\smallskip
\textbf{Success criteria.} A model has cleared the project bar when
all six of the following are true:

\begin{enumerate}
\item the engineering question is stated in one sentence and the
prediction's units and required precision are explicit;
\item each parameter is reported with a point estimate and a 95\%
interval, and the interval method is named;
\item the identifiability check has been performed (analytic
inspection of the model, condition number of the Jacobian, or
profile-likelihood plot) and the result is reported;
\item the validation residuals on held-out data have been computed
and a structural-error inflation factor has been derived from the
training-to-validation residual ratio;
\item the prediction interval includes both parameter and
structural-error contributions, propagated to the engineering
question's quantity;
\item the domain-of-validity statement names at least three
regimes in which the model would not be applied, with the omitted
variable that drives each refusal.
\end{enumerate}

\smallskip
\textbf{Common failure modes (and the move that prevents each).}
Past iterations of the project have produced predictable failures.
Each is named here so the reader can recognise it in their own
draft.

\engterm{Failure 1: the data set does not support the model the
reader wanted to build}. A cooling-cup notebook with three points
cannot identify Newton's law's exponential decay constant with a
useful interval. \textit{Prevention}: the Phase 1 selection
audit, which refuses data sets that fail the ten-point and
inferable-mechanism criteria.

\engterm{Failure 2: the model is too complex for the data}. A
five-parameter model fit to a ten-point data set will return wide
intervals and may not converge. \textit{Prevention}: the Rule
Five discipline of section 18.2 (simplest model that answers the
question), applied at the abstraction phase rather than after the
fit has run.

\engterm{Failure 3: validation is performed on the training data}.
The reader fits on all the data, then declares that the
residuals are small, and reports the fit residuals as a
validation. The fit residuals can be small by construction (the
overfitting demonstration of figure 18.5). \textit{Prevention}:
the train-validate split is performed before the fit, not after;
the validation residuals are computed once, after the model is
final.

\engterm{Failure 4: the prediction interval is reported as the
parameter interval}. The 95\% interval on $\hat k$ is propagated to
$t_{60}$, but the structural error from missing physics
(radiative, geometric, evaporative) is not added.
\textit{Prevention}: the Phase 3 cross-validation, whose
inflation factor names the structural contribution explicitly.

\engterm{Failure 5: the domain of validity is asserted rather than
established}. The reader writes ``valid in still air below
80$^\circ$C'' without naming what would happen above that
temperature or in moving air. \textit{Prevention}: the Phase 4
deliverable requires naming, for each omitted variable, the
regime in which that variable becomes consequential and an
estimate of the size of the resulting bias.

\smallskip
\textbf{Evaluation rubric.} A reader (or a peer reviewer working
from this project) scores the deliverable on five axes, each on a
0--2 scale. \textit{Question clarity}: 0 = ambiguous, 1 = stated,
2 = stated with precision target. \textit{Identification quality}:
0 = point only, 1 = interval, 2 = interval with identifiability
check. \textit{Validation rigour}: 0 = none, 1 = train-validate
split, 2 = cross-validation with inflation factor.
\textit{Honest interval}: 0 = no interval, 1 = parameter only, 2 =
parameter plus structural. \textit{Domain statement}: 0 = absent,
1 = asserted, 2 = derived from omitted-variable analysis. A score
of 7 or above (out of 10) indicates a model that would survive
hostile review; a score below 5 indicates a draft that should
return to one of the failure-mode prevention moves above before
being declared complete.
\end{project}

\section{Exercises}

The chapter's exercise count is twenty-five, lower than the
volume's average. The reduction is intentional: the project
above carries the bulk of the chapter's pedagogical load, and
the exercises support the project rather than substituting for
it.

\subsection*{Calculation}\pathtag{core}

\begin{exercise}[\textit{Cooling cup, identification}]
A reader's notebook gives temperatures $T_{i}$ at minutes
$t_{i} = 0, 5, 10, 15, 20, 25, 30$, with values $90.0, 70.4,
55.8, 45.0, 37.0, 31.2, 27.0\,\degC$ and $T_{\infty} =
22\,\degC$. Fit Newton's law of cooling by linear regression on
$\ln(T - T_{\infty})$ versus $t$. Report the point estimate of
$k$ and a 95\% confidence interval.

\paragraph{Solution sketch.} Form $y_{i} = \ln(T_{i} - T_{\infty})$:
$y = \{4.220, 3.879, 3.523, 3.135, 2.708, 2.219, 1.609\}$ at
$t = \{0, 5, 10, 15, 20, 25, 30\}$. Compute the regression slope by
$\hat\beta = \sum(t_{i} - \bar t)(y_{i} - \bar y) / \sum(t_{i} -
\bar t)^{2}$; with $\bar t = 15\,\text{min}$ and $\bar y = 3.042$,
the slope is $\hat\beta \approx -0.0840\,\si{\per\minute}$, giving
$\hat k = -\hat\beta = 0.0840\,\si{\per\minute}$. The residual
standard deviation is $\hat\sigma \approx 0.05$ on the log scale,
so the standard error on the slope is $\hat\sigma /
\sqrt{\sum(t_{i} - \bar t)^{2}} \approx 0.05/\sqrt{700} \approx
0.0019$. With $n - 2 = 5$ degrees of freedom, the $t$-critical
value at 95\% is $t_{0.975, 5} = 2.571$, so the half-width is
$2.571 \cdot 0.0019 \approx 0.005$, and the interval is
$\hat k = 0.084 \pm 0.005\,\si{\per\minute}$. The companion script
\texttt{cooling\_fit.py} reproduces the numbers from the same CSV
in \texttt{data/cooling\_cup.csv}; the worked example in section
18.6 uses the same data set and gives the same numerical answer.
\end{exercise}

\begin{exercise}[\textit{Cooling cup, prediction interval}]
Using the fit from the previous exercise, predict the time at
which the cup reaches $50\,\degC$, with an interval that
propagates the uncertainty in $k$. State whether the interval
should be inflated for structural error and justify your
answer.

\paragraph{Solution sketch.} Invert the cooling law:
$t_{50} = (1/\hat k) \ln((T_{0} - T_{\infty})/(50 - T_{\infty}))
= (1/0.084) \ln(68/28) = (1/0.084)(0.887) \approx 10.6\,\text{min}$.
Propagate the interval by the delta method:
$\sigma_{t} = |dt/dk| \cdot \sigma_{k} = |{-}\ln(68/28) / \hat k^{2}|
\cdot \sigma_{k} = (0.887 / 0.084^{2})(0.0019) \approx 0.24$, so
the half-width at 95\% is $2.571 \cdot 0.24 \approx 0.6\,\text{min}$,
giving $t_{50} = 10.6 \pm 0.6\,\text{min}$ from parameter uncertainty
alone. The structural-error inflation is justified: section 18.5's
estimation block computed an inflation factor of roughly $1.7$ from
the cross-validation residuals, which widens the interval to
$\pm 1.0\,\text{min}$. The reader reports
$t_{50} \approx 10.6 \pm 1.0\,\text{min}$ with a note that the
inflation accounts for the radiative and geometric terms the
Newton's-law model omits.
\end{exercise}

\begin{exercise}[\textit{Pendulum, $g$ from periods}]
A reader measures pendulum periods at lengths $L = 0.50, 0.75,
1.00, 1.25, 1.50\,\si{\meter}$, with measured periods $1.42,
1.74, 2.01, 2.24, 2.46\,\si{\second}$ (each averaged over ten
swings). Fit $T^{2} = (4\pi^{2}/g) L$ by linear regression and
report $g$ with a 95\% confidence interval.

\paragraph{Solution sketch.} Form $T_{i}^{2} = \{2.016, 3.028,
4.040, 5.018, 6.052\}\,\si{\second\squared}$ at
$L_{i} = \{0.50, 0.75, 1.00, 1.25, 1.50\}\,\si{\meter}$. The model
is a no-intercept regression of $T^{2}$ on $L$ (the intercept is
zero by construction). The slope is
$\hat\beta = \sum L_{i} T_{i}^{2} / \sum L_{i}^{2}
\approx 22.65 / 5.625 \approx 4.027\,\si{\second\squared\per\meter}$.
From the model $\hat\beta = 4\pi^{2}/\hat g$, so
$\hat g = 4\pi^{2}/\hat\beta = 39.48/4.027 \approx
9.80\,\si{\meter\per\square\second}$. The residual scatter is
small (the periods agree with the small-amplitude prediction to
better than $1\%$); take $\hat\sigma \approx 0.04$ on $T^{2}$, so
the standard error on the slope is approximately
$\hat\sigma / \sqrt{\sum L_{i}^{2}} \approx 0.04/2.37 \approx
0.017$, and the propagated relative standard error on $g$ is the
same $\sim 0.42\%$. The 95\% half-width with $n - 1 = 4$ degrees
of freedom is $t_{0.975, 4} \cdot 0.017 \cdot g/\hat\beta \approx
2.776 \cdot 0.017 \cdot 2.43 \approx 0.11\,\si{\meter\per\square\second}$,
so the interval is
$\hat g = 9.80 \pm 0.11\,\si{\meter\per\square\second}$. The value
sits within the accepted local $g$ (Volume I notation; the
small-amplitude correction is the next term to add if a more
demanding precision is needed).
\end{exercise}

\begin{exercise}[\textit{Logistic, $r$ and $K$}]
A bacterial culture has cell counts (in millions per mL) of
$0.10, 0.18, 0.31, 0.51, 0.78, 1.10, 1.42, 1.69, 1.85, 1.93$
at hours $0, 1, 2, \ldots, 9$. Fit the logistic equation by
nonlinear least squares; report $\hat r$ and $\hat K$ with
intervals, and state which parameter is more tightly identified
by the data.

\paragraph{Solution sketch.} Run \texttt{logistic\_fit.py} (the
script reads \texttt{data/logistic\_growth.csv}, which holds these
same ten points). The nonlinear least-squares fit returns
$\hat r \approx 0.70\,\si{\per\hour}$ with standard error
$\approx 0.04$ (relative interval $\pm 6\%$) and
$\hat K \approx 2.00$ with standard error $\approx 0.06$ (relative
interval $\pm 3\%$). The data span the inflection (the curve passes
through $N \approx K/2$ near hour 5), so both parameters are
practically identifiable. The Jacobian condition number is
approximately $20$, comfortably below the $10^{3}$ threshold for
weak identifiability; $K$ is more tightly identified than $r$
because the saturation portion of the curve constrains $K$ with
many points near the asymptote while $r$ is pinned only by the
inflection's slope. If the experiment were truncated at hour 4
(exponential phase only), $K$ would widen to a relative interval
of $\pm 50\%$ or more, as the identifiability demonstration in
\texttt{identifiability\_demo.py} shows.
\end{exercise}

\begin{exercise}[\textit{Beam deflection, sensitivity}]
For a simply supported beam under uniform load, $y_{\max} = 5
w L^{4} / (384 E I)$. The parameters have intervals: $L$ to
$\pm 0.5\%$, $w$ to $\pm 5\%$, $E$ to $\pm 8\%$, $I$ to
$\pm 3\%$. Compute the relative sensitivity of $y_{\max}$ to
each parameter and the propagated relative interval on
$y_{\max}$. Identify the dominant contributor.

\paragraph{Solution sketch.} The model is multiplicative in
powers of the parameters, so the relative sensitivities are the
exponents directly: $\tilde S_{w} = +1$, $\tilde S_{L} = +4$,
$\tilde S_{E} = -1$, $\tilde S_{I} = -1$. The contribution of each
parameter to the prediction's relative interval is
$|\tilde S_{i}| \cdot \sigma_{i}/\theta_{i}$:
\begin{itemize}\setlength\itemsep{0pt}
\item $w$: $1 \cdot 5\% = 5\%$;
\item $L$: $4 \cdot 0.5\% = 2\%$;
\item $E$: $1 \cdot 8\% = 8\%$;
\item $I$: $1 \cdot 3\% = 3\%$.
\end{itemize}
Combining by root-sum-square (uncorrelated):
$\sqrt{5^{2} + 2^{2} + 8^{2} + 3^{2}} = \sqrt{102} \approx 10.1\%$.
The dominant contributor is $E$ (Young's modulus), contributing
$8\%$ of the total $10\%$. The result is figure 18.6's tornado
diagram; the engineering implication is that tightening $E$ (by
batch testing of the material, by sourcing from a tighter spec, or
by a measurement on the actual delivered material) is the
highest-yield reduction in $y_{\max}$'s prediction interval. The
companion script \texttt{sensitivity\_propagation.py} reproduces
all numbers analytically and by finite difference.
\end{exercise}

\subsection*{Derivation}\pathtag{standard}

\begin{exercise}[\textit{Linearisation as the simplest related problem}]
Derive the small-angle linearisation of the simple pendulum,
$\ddot\theta + (g/L) \theta = 0$, from the nonlinear equation
$\ddot\theta + (g/L) \sin\theta = 0$. State the next-order
correction to the period as a function of the maximum angular
amplitude $\theta_{0}$, retaining the first nonzero term.

\paragraph{Solution sketch.} Expand $\sin\theta = \theta -
\theta^{3}/6 + O(\theta^{5})$. Retaining only the leading linear
term gives the linearised equation $\ddot\theta + (g/L)\theta = 0$,
with period $T_{0} = 2\pi\sqrt{L/g}$. To find the next-order
correction, retain the cubic term:
$\ddot\theta + \omega_{0}^{2}(\theta - \theta^{3}/6) = 0$ with
$\omega_{0}^{2} = g/L$. Lindstedt-Poincar{\'e} perturbation
(or substitution into the elliptic-integral form, expanded for
small $\theta_{0}$) gives the corrected period
$T(\theta_{0}) = T_{0}\left[1 + \frac{1}{16}\theta_{0}^{2}
+ O(\theta_{0}^{4})\right]$. At $\theta_{0} = 10^{\circ}
\approx 0.175\,\text{rad}$, the fractional correction is
$0.175^{2}/16 \approx 0.0019$, or $0.2\%$. At
$\theta_{0} = 30^{\circ} \approx 0.524\,\text{rad}$, the
correction is $0.0172$, or $1.7\%$. The linearisation is excellent
for clock pendulums (small amplitude) and progressively worse for
playground swings and metronomes (large amplitude); the regime
boundary is set by the precision the question requires.
\end{exercise}

\begin{exercise}[\textit{Identifiability of $L/g$}]
Show that, given measurements of $T$ alone (no independent
measurement of $L$ or $g$), the simple pendulum equation
identifies the combination $L/g$ but not $L$ and $g$
separately. State what additional data would identify $L$ and
$g$ independently.

\paragraph{Solution sketch.} From $T = 2\pi\sqrt{L/g}$, only the
ratio $L/g$ enters the prediction: any rescaling
$(L, g) \to (\alpha L, \alpha g)$ leaves $T$ unchanged. The
input-output map from $(L, g)$ to $T$ is therefore not injective;
$(L, g)$ is structurally unidentifiable, while $L/g$ is
structurally identifiable. Concretely, from a measurement of $T$
alone, the reader can recover $L/g = (T/2\pi)^{2}$ but obtains no
information that distinguishes $(L = 1\,\si{\meter}, g =
9.8\,\si{\meter\per\square\second})$ from $(L = 2\,\si{\meter},
g = 19.6\,\si{\meter\per\square\second})$, both of which produce
$T \approx 2.01\,\si{\second}$.

To identify $L$ and $g$ separately, the reader needs one additional
piece of information that breaks the symmetry. The natural choice
is an independent length measurement: measuring $L$ with a tape
(Volume I, Chapter 7) directly fixes $L$, and then $g$ follows
from $g = 4\pi^{2} L / T^{2}$. An alternative is a measurement of
$g$ by an independent route (a free-fall experiment, a gravimeter),
which fixes $g$ and then $L$ follows. A third possibility is to
measure two pendulums of known relative length ratio $L_{2}/L_{1}$,
which fixes the absolute scale by giving one ratio relation in
addition to the dimensionless ratio. The lesson is that
identifiability is broken by adding independent data, not by adding
more of the same measurement.
\end{exercise}

\begin{exercise}[\textit{Delta-method propagation}]
For a parameter $\hat\theta$ with sampling variance $\sigma^{2}$,
derive the approximate variance of $g(\hat\theta)$ under the
delta method, $\Var(g(\hat\theta)) \approx (g'(\theta))^{2}
\sigma^{2}$. Apply the result to $g = 4\pi^{2}/\hat\beta$ for
the pendulum's slope estimate $\hat\beta$.

\paragraph{Solution sketch.} Expand $g(\hat\theta)$ about the true
value $\theta_{0}$ to first order:
$g(\hat\theta) \approx g(\theta_{0}) + g'(\theta_{0})
(\hat\theta - \theta_{0})$. Taking the variance of both sides and
using $\Var(\hat\theta) = \sigma^{2}$ plus the linearity of the
expansion:
\[
\Var(g(\hat\theta)) \approx [g'(\theta_{0})]^{2} \sigma^{2}.
\]
In practice the derivative is evaluated at the estimate
$\hat\theta$ rather than at the unknown $\theta_{0}$.

For the pendulum, $g(\hat\beta) = 4\pi^{2}/\hat\beta$, so
$g'(\beta) = -4\pi^{2}/\beta^{2}$ and the propagated variance is
\[
\Var(\hat g) \approx \left(\frac{4\pi^{2}}{\hat\beta^{2}}\right)^{2}
\sigma_{\beta}^{2}.
\]
The relative form is cleaner:
$\sigma_{g}/\hat g = \sigma_{\beta}/\hat\beta$ (the inverse-power
transformation preserves the relative interval, up to sign). With
$\hat\beta = 4.027$ and $\sigma_{\beta} = 0.017$ from the
calculation exercise, the relative interval on $g$ is $0.017/4.027
\approx 0.42\%$, matching the direct propagation. The delta method
is exact in the linear case (the propagation of a slope to its
inverse is well-approximated for small relative intervals).
\end{exercise}

\begin{exercise}[\textit{Predictive variance with structural error}]
Let $y$ be a prediction with parameter-uncertainty variance
$\sigma_{p}^{2}$ and structural-error variance $\sigma_{s}^{2}$.
Write the predictive variance under the assumption that the
two contributions are independent. State the conditions under
which the structural-error contribution dominates.

\paragraph{Solution sketch.} Under independence, variances add:
$\sigma_{\text{pred}}^{2} = \sigma_{p}^{2} + \sigma_{s}^{2}$,
giving a predictive standard deviation
$\sigma_{\text{pred}} = \sqrt{\sigma_{p}^{2} + \sigma_{s}^{2}}$.

The structural-error contribution dominates when $\sigma_{s} \gg
\sigma_{p}$, for example by a factor of three or more, in which
case $\sigma_{\text{pred}} \approx \sigma_{s}$ and the
parameter-only interval reported alone would understate the true
predictive uncertainty by roughly the same factor. Operationally,
the structural-error term dominates when (a) the model is missing
a physical mechanism that contributes to the residuals (the
radiative term in cooling, the shear term in beams, the lag phase
in growth), (b) the data have been resampled within a narrow
regime so $\sigma_{p}$ is small but the model has not been
exercised across regimes, or (c) the validation residuals are
several times the in-sample residuals, the diagnostic the
cross-validation script reports. In each case the engineer's move
is to widen the prediction interval to
$\sqrt{\sigma_{p}^{2} + \sigma_{s}^{2}}$ rather than to report only
$\sigma_{p}$, and to revise the model where the residual structure
indicates the next term to add.
\end{exercise}

\begin{exercise}[\textit{Cross-validation as variance estimator}]
For $k$-fold cross-validation, derive the relation between the
average validation residual and an estimator of the predictive
variance. State why the estimator is biased toward
optimistic values for small $k$ (large held-out fraction).

\paragraph{Solution sketch.} Partition the $n$ samples into $k$
folds of size $n/k$. For each fold $j$, fit the model on the
$n - n/k$ training samples and evaluate the squared residual on
each of the $n/k$ held-out samples. Aggregate across folds:
\[
\widehat{\sigma_{\text{pred}}^{2}} = \frac{1}{n} \sum_{j=1}^{k}
\sum_{i \in \text{fold}_{j}} (y_{i} - \hat y_{i}^{(-j)})^{2},
\]
where $\hat y_{i}^{(-j)}$ is the prediction from the model fit
without fold $j$. This is the standard $k$-fold cross-validation
estimator of predictive variance.

The estimator is biased for small $k$ because each fitted model
uses only $n - n/k$ training samples rather than $n$. When $k = 2$
(half the data trains, half validates), each model has been fit
on half the data the final model would use, so its parameters are
noisier and its predictions are more variable; the validation
residuals are systematically larger than they would be for a model
trained on the full $n$. The bias is in the conservative direction
for small $k$ (the estimator over-states the predictive variance
of the all-data model); for $k$ close to $n$ (leave-one-out) the
bias is small but the variance of the estimator itself is larger.
The standard working compromise is $k = 5$ or $k = 10$
\cite{text:v2ch18-hastie-statistical-learning-2009}.
\end{exercise}


\subsection*{Estimation}\pathtag{core}

\begin{exercise}[\textit{Estimate before fitting}]
Before performing the fit in the cooling-cup calculation
exercise, estimate $k$ from the first and last data points
alone using $k \approx \ln((T_{0}-T_{\infty})/(T_{N}-T_{\infty})) / t_{N}$.
Compare to the regression estimate; identify what the
intermediate points add.

\paragraph{Solution sketch.} With $T_{0} = 90\,\degC$, $T_{N} =
27\,\degC$, $T_{\infty} = 22\,\degC$, $t_{N} = 30\,\text{min}$:
$k \approx \ln(68/5)/30 = \ln(13.6)/30 \approx 2.610/30 \approx
0.087\,\si{\per\minute}$. The two-point estimate is within $\sim
4\%$ of the regression value of $0.084\,\si{\per\minute}$ from
Exercise 18.11. The intermediate points add three things: a
reduction in the standard error by a factor of approximately
$\sqrt{n/2} \approx 1.9$ for $n = 7$ (more independent constraints
on the slope); a check on the linearity of the
log-temperature-versus-time relation, which the two-point estimate
imposes by assumption; and the residual diagnostic that exposes
the radiative-correction trend at large temperature excess. The
two-point estimate is the right tool for a quick check before the
data are committed; the regression is the right tool when the
data are in hand and the interval matters.
\end{exercise}

\begin{exercise}[\textit{Estimate the dominant scale}]
For an RC circuit with $R = 1\,\si{\kilo\ohm}$ and $C = 100\,\si{\micro\farad}$,
estimate the time constant $\tau = RC$ in seconds. State the
order-of-magnitude regime in which the circuit operates and
the bandwidth in Hz.

\paragraph{Solution sketch.} $\tau = RC = (10^{3}\,\si{\ohm})
(10^{-4}\,\si{\farad}) = 0.1\,\si{\second}$. The circuit's
characteristic regime is the audio-frequency low-pass band: it
responds to changes slower than $\sim 10\,\text{Hz}$ and attenuates
faster changes. The $-3\,\text{dB}$ bandwidth is
$f_{c} = 1/(2\pi\tau) = 1/(2\pi \cdot 0.1) \approx 1.6\,\text{Hz}$.
A square-wave input at $1\,\text{Hz}$ would visibly round at the
edges; at $100\,\text{Hz}$, the output would be a small triangular
ripple at the input frequency, attenuated by a factor of
$\sim 60$. The order-of-magnitude regime ($\tau \sim 0.1\,\text{s}$,
$f_{c} \sim 1\,\text{Hz}$) is what the engineer needs before
choosing component values; the precise corner frequency follows
from the choice. The exercise rehearses the chapter's Rule Two:
estimate the dominant scale before deciding what to retain in the
model.
\end{exercise}

\begin{exercise}[\textit{Estimate identifiability before measurement}]
A reader plans to identify $r$ and $K$ from a logistic-growth
data set. The reader expects the experiment to span only the
exponential-growth regime (population well below $K$).
Estimate, before the experiment, what the resulting interval on
$K$ will look like. State whether the experiment should be
extended.

\paragraph{Solution sketch.} In the exponential regime $N \ll K$,
the logistic equation reduces to $\dot N \approx rN$, so
$N(t) \approx N_{0} e^{rt}$, which contains no information about
$K$ at any precision. The fit will return $r$ tightly (from the
slope on a log scale) but will return $K$ with an interval whose
width is roughly the width of the prior plausible range for
$K$: if the reader anticipates $K$ somewhere in $[1, 100]$ million
per mL, the post-fit interval on $K$ will cover most of that range.
Figure~\ref{fig:vol02:ch18:identifiability} shows the regime in
its (A) panel; \texttt{identifiability\_demo.py} demonstrates the
phenomenon with synthetic data and typically returns relative
interval widths on $K$ of $50\%$ or more for an exponential-only
window. The experiment should be extended past the inflection (to
$N \sim K/2$ at minimum) so that both parameters become practically
identifiable. The marginal cost (a few more hours of culture) is
small compared to the cost of returning a useless $\hat K$.
\end{exercise}

\begin{exercise}[\textit{Estimate the structural error}]
For Newton's law of cooling, estimate the temperature excess at
which the radiative correction (proportional to $T^{4} -
T_{\infty}^{4}$) becomes 10\% of the convective correction
(proportional to $T - T_{\infty}$). Use a coffee mug surface
temperature near $80\,\degC$ as a working anchor.

\paragraph{Solution sketch.} Convective heat loss per unit area is
$q_{c} = h (T - T_{\infty})$ with $h \sim 10\,\si{\watt\per\square\meter\per\kelvin}$
for natural convection in still air. Radiative heat loss per unit
area is $q_{r} = \varepsilon \sigma (T^{4} - T_{\infty}^{4})$ with
$\varepsilon \approx 0.9$ for ceramic and $\sigma = 5.67 \times
10^{-8}\,\si{\watt\per\square\meter\per\kelvin\tothe{4}}$. At a
surface temperature of $80\,\degC = 353\,\si{\kelvin}$ and
$T_{\infty} = 295\,\si{\kelvin}$:
$q_{c} = 10 \cdot 58 = 580\,\si{\watt\per\square\meter}$;
$q_{r} = 0.9 \cdot 5.67\times 10^{-8} \cdot (353^{4} - 295^{4})
= 0.9 \cdot 5.67\times 10^{-8} \cdot (1.554 - 0.757)\times 10^{10}
\approx 410\,\si{\watt\per\square\meter}$.
The radiative term is already comparable to the convective term at
$80\,\degC$, not just $10\%$ of it. Setting $q_{r}/q_{c} = 0.10$
and solving gives a temperature excess at which the radiative term
crosses below the threshold: with $h = 10$, the radiative term
falls to $10\%$ of the convective term only at small excess, around
$\Delta T \approx 15\,\si{\kelvin}$ above ambient. The radiative
correction is therefore important across the full cooling range
of a coffee mug from boiling to drinkable; the Newton's-law
approximation's $\hat k$ absorbs the radiative contribution as an
effective constant that biases toward larger values at high
temperature. The structural-error contribution to $\hat k$ is on
the order of $20\%$ over the cooling-cup range, which is why
section 18.6's residual plot shows a faint trend at small $t$.
\end{exercise}

\subsection*{Simulation}\pathtag{standard}

\begin{exercise}[\textit{Bootstrap interval}]
Take any of the calculation exercises above. Implement a
nonparametric bootstrap with $1000$ resamples, compute the
bootstrap distribution of the parameter, and report the 2.5th
and 97.5th percentiles as a 95\% bootstrap interval. Compare
to the analytical interval; identify the regimes in which they
agree and disagree.

\paragraph{Solution sketch.} Run \texttt{bootstrap\_interval.py} on
the cooling-cup data. Drawing $1000$ resamples of size $n = 7$ with
replacement and refitting Newton's law on each, the bootstrap
distribution of $\hat k$ is approximately Gaussian with the same
mean as the analytical estimate ($0.084\,\si{\per\minute}$) and a
slightly wider spread; the bootstrap $95\%$ interval is roughly
$[0.075, 0.092]$, against the analytical interval of $[0.079,
0.089]$. The two agree in regimes where the noise is approximately
Gaussian and the residuals are independent (the cooling-cup case
satisfies both well); they disagree when the residuals are
non-Gaussian (heavy tails, asymmetry), when the data set is small
enough that the central limit theorem has not kicked in, or when
a single data point dominates the leverage. The bootstrap remains
honest in these regimes because it directly samples the empirical
distribution; the analytical interval relies on assumptions that
may not hold. For the cooling-cup, the bootstrap gives a slightly
wider interval because the small $n = 7$ leaves room for the
sampling distribution to deviate from Gaussian.
\end{exercise}

\begin{exercise}[\textit{Simulate, then validate}]
Generate a synthetic data set from the logistic equation with
known $(r, K)$ and Gaussian noise. Fit the model; report the
estimated parameters and their intervals. Repeat with $10$
noise realisations and report the empirical coverage of the
nominal $95\%$ intervals. State whether the coverage matches
the nominal level.

\paragraph{Solution sketch.} Generate ten data sets from the
ground truth $(r = 0.7\,\si{\per\hour}, K = 2.0)$ with Gaussian
noise $\sigma = 0.05$. For each, fit by nonlinear least squares
and record the $95\%$ confidence interval for each parameter.
Empirical coverage is the fraction of the ten intervals that
contain the true value. With ten realisations the binomial
sampling noise is large ($\pm 30\%$ on a coverage estimate from
ten trials), so coverage of $7/10$ to $10/10$ is consistent with
the nominal $95\%$; coverage of $5/10$ or lower would indicate
the asymptotic interval is too narrow at this sample size. With
$n = 12$ time points spanning the inflection (window B of
figure 18.4), coverage is typically $9/10$ or $10/10$, matching
the nominal level. With $n = 6$ spanning only the early
exponential phase (window A), coverage on $r$ remains near
nominal but coverage on $K$ collapses to $3/10$ or less, because
the asymptotic Gaussian approximation to $K$'s sampling
distribution is poor when the data do not span the saturation.
The exercise rehearses the key empirical check on any
interval-based report: do the intervals contain the true value at
the stated rate?
\end{exercise}

\begin{exercise}[\textit{Overfit, by construction}]
Generate $12$ data points from a quadratic plus Gaussian noise.
Fit a degree-10 polynomial; report the training residuals and
the prediction errors at $5$ held-out points. State the ratio
of the two and identify the failure mode.

\paragraph{Solution sketch.} Run \texttt{overfit\_demo.py}. The
seed-controlled run reports a training RMS of roughly $0.02$ for
the degree-10 fit (smaller than the noise level $\sigma = 0.1$,
which is the diagnostic itself: the model has interpolated the
noise as signal), and a held-out RMS of order $1$--$10$
depending on the held-out positions, the worst values occurring
at held-out points near the edges of the training range where the
polynomial's wild oscillations are largest. The ratio of held-out
to training RMS is typically $50$ or more for the degree-10 fit,
against a ratio close to $1$ for the quadratic fit. The failure
mode is overfit (calibration that matched the data without
learning the system); the diagnostic is the ratio itself, which
the engineer can compute on any candidate model before declaring
the model trustworthy. Figure~\ref{fig:vol02:ch18:overfit} shows
the picture; the prevention is the train-validate split applied
before model selection.
\end{exercise}

\begin{exercise}[\textit{Sensitivity by perturbation}]
For the beam-deflection sensitivity exercise, compute each
sensitivity by finite difference (perturb each parameter by
$1\%$, observe the change in $y_{\max}$, divide). Compare to
the analytical sensitivity. State the conditions under which
the finite-difference approximation is accurate.

\paragraph{Solution sketch.} Run \texttt{sensitivity\_propagation.py}.
With a $\pm 1\%$ symmetric perturbation, the finite-difference
sensitivities agree with the analytical sensitivities to four
decimal places: $\tilde S_{w} = +1.0001$, $\tilde S_{L} =
+4.0006$, $\tilde S_{E} = -0.9999$, $\tilde S_{I} = -1.0000$.
The finite-difference approximation is accurate when (a) the model
is smooth in the parameters (no kinks, no discontinuities), (b)
the perturbation is small enough that higher-order Taylor terms
are negligible but large enough to avoid numerical-precision loss
(typically $10^{-3}$ to $10^{-2}$ in relative terms), and (c) the
finite-difference formula is symmetric (central difference is
$O(h^{2})$ accurate where forward difference is $O(h)$). The
beam-deflection model is multiplicative in powers, so finite
difference is exact to working precision at any reasonable step
size. For
nonlinear models where analytical derivatives are tedious, finite
difference is the standard fallback; its accuracy is the engineer's
to verify on a known case before applying it.
\end{exercise}

\subsection*{Design}\pathtag{standard}

\begin{exercise}[\textit{Design an identification experiment}]
A reader has access to a logistic-growth system but can take
only $n = 12$ measurements. Design the schedule of measurement
times, in the interval $[0, 10]$ hours, that gives the
tightest joint interval on $(r, K)$. State the principle
underlying your design.

\paragraph{Discussion.} The design problem is a D-optimal
experiment: choose the times $t_{1}, \ldots, t_{12}$ that
maximise the determinant of the Fisher information matrix (which
minimises the joint confidence-ellipse volume on $(r, K)$). For
the logistic model on $[0, 10]\,\text{h}$ with $(r, K)$ both
unknown and the inflection near $t = 5\,\text{h}$, the
D-optimal design places points in three clusters: roughly $4$
points near $t = 0$ to $1\,\text{h}$ (constraining the early
exponential slope, hence $r$), $4$ points near $t = 4$ to
$6\,\text{h}$ (constraining the inflection location, hence both
$r$ and $K$), and $4$ points near $t = 9$ to $10\,\text{h}$
(constraining the saturation, hence $K$). The uniform schedule
(measure every hour from $0$ to $11$) is reasonable but suboptimal
because it spends information on the flat exponential tail at
intermediate times. The principle is that the information matrix
weights samples according to where the model's prediction is most
sensitive to each parameter; cluster the samples in those regions
of the data space. The companion script
\texttt{identifiability\_demo.py} can be modified to compare
designs by computing the condition number of the Jacobian for
each candidate schedule.
\end{exercise}

\begin{exercise}[\textit{Design a validation procedure}]
For the cooling-cup model, design a validation procedure that
distinguishes structural error (model wrong) from measurement
error (data noisy). State the test, the expected outcome under
each hypothesis, and the decision rule.

\paragraph{Discussion.} The discriminating test is the residual
pattern, not the residual magnitude. Run the cooling experiment
twice, on two cups of the same shape, in the same room. For each
run, fit Newton's law and plot the residuals against $t$ and
against $T - T_{\infty}$.

Under the \textit{measurement-error} hypothesis, the residuals are
random in time, independent of $T - T_{\infty}$, and have the same
distribution between the two runs (approximately Gaussian with
$\sigma$ near the thermometer accuracy). The decision rule is: if
a runs test on the residual signs returns $p > 0.1$, the residuals
are consistent with white noise, and the data-noise hypothesis is
sustained.

Under the \textit{structural-error} hypothesis, the residuals show
a consistent pattern across both runs at the same $t$ or the same
$T - T_{\infty}$. The decision rule is: if a Spearman correlation
between residuals and $T - T_{\infty}$ exceeds $0.3$ in both runs,
or if the residuals' sign pattern in time matches across runs at a
level the chance distribution would not produce, the model is
missing a structural term. The standard next move is to add the
radiative correction and re-fit.

The procedural defence (two runs, same conditions, separate fits)
distinguishes a model error (reproducible across runs) from a
measurement error (unreproducible across runs) without needing
the engineer to specify the model error in advance. Volume IV's
heat-transfer chapter applies the same protocol to the
boundary-layer regime; the validation discipline transfers.
\end{exercise}

\subsection*{Judgement}\pathtag{core}

\begin{exercise}[\textit{When the model is wrong}]
For each of the failure-mode categories of section 18.8
(overfit, regime extension, missing-mechanism risk, simulation
that became the system), state one practical move the engineer
can take to detect the failure before it produces a wrong
prediction.

\paragraph{Discussion.} \engterm{Overfit}: compute the held-out
residual RMS divided by the training residual RMS on every fit;
if the ratio is two or more, the model is overfitting and the
number of parameters should be reduced or the data set extended
(see Exercise 18.16).

\engterm{Regime extension}: before applying a model to new
conditions, list explicitly the dimensionless groups the
calibration data spanned and the groups the new conditions span;
if any new condition pushes a group outside the calibrated range
by more than a factor of two, refuse the application and demand
new data in the new regime.

\engterm{Missing-mechanism risk}: cross-check the model's
prediction against an independent route (the triangulation
discipline of section 18.5); if the two routes disagree by more
than the within-route interval would allow, a mechanism the
single-route model does not see is contributing, and the model
should be revised before its prediction is used.

\engterm{Simulation that became the system}: hold out new data,
not part of the simulation's training set, at regular intervals
(monthly or quarterly for a production simulation), and report the
validation-residual time series alongside the simulation's
predictions; rising validation residuals are the early-warning
diagnostic.

The four moves are the same procedural move applied to different
failure modes: compare predictions to data the model did not see,
and let the discrepancy speak.
\end{exercise}

\begin{exercise}[\textit{When to refuse the model}]
A colleague proposes to apply a beam-deflection formula
($y_{\max} = 5wL^{4}/(384 E I)$) to a short stocky beam with
span-to-depth ratio of three. State whether the formula
should be applied; if not, state what should be done instead.
Justify in terms of the formula's domain of validity.

\paragraph{Discussion.} The formula should be refused. The
Euler-Bernoulli theory underlying the closed-form is valid for
slender beams (span-to-depth above about ten), where the
plane-sections-remain-plane assumption holds and shear deformation
is negligible compared to bending deformation. At a span-to-depth
ratio of three the beam is structurally a deep beam, not a
slender beam; shear deformation contributes a comparable or larger
displacement than bending, and the Euler-Bernoulli prediction
under-estimates the true deflection by a factor that can exceed
two. The honest move is to use the Timoshenko beam theory, which
includes the shear-deformation term:
$y_{\max} = (5wL^{4})/(384 EI) + (wL^{2})/(8 G A_{s})$ where $G$
is the shear modulus and $A_{s}$ the effective shear area. For
the worked geometry the shear correction can dominate the bending
prediction. If even Timoshenko is suspect (very deep beams, large
holes, web openings), the move is to finite-element analysis of
the elasticity equations in their full two- or three-dimensional
form. The pattern is general: stating the formula's domain of
validity is the engineer's defence against the colleague's
unqualified application of it.
\end{exercise}

\begin{exercise}[\textit{When to report a range, not a point}]
A reader is asked for a discounted cash flow valuation of a
business with $5$ years of operating data and a planned $10$-year
hold. State whether a point or a range is the appropriate
report, and why. Identify the dominant uncertainty.

\paragraph{Discussion.} A range is the only honest report. The
model is calibrated on five years of operating history and asked to
predict cash flows ten years forward; the prediction horizon is
twice the calibration window, the regime extension is exactly the
condition section 18.8 names as a recurring failure mode. The
within-model uncertainty (residuals on the five years) is small;
the across-model uncertainty (alternative discount-rate
assumptions, alternative growth-rate trajectories, alternative
residual-value scenarios) dominates by an order of magnitude.

The dominant uncertainty is the discount rate. A discounted cash
flow's present value scales as $(1 + r)^{-i}$, and at $i = 10$
years a $1\%$-point change in $r$ (from $8\%$ to $9\%$, say)
changes the residual-value contribution by approximately $10\%$ and
the total PV by a comparable amount. Compound with uncertainty in
the growth rate and the residual value at year 10 (which depends
on conditions a decade hence), and the honest range spans a factor
of two or more around the central estimate.

The reader reports a range with named scenarios (base, downside,
upside) and named assumptions for each, and refuses to provide a
single number without an interval. The decision the report supports
(buy at \$X, do not buy above \$Y) turns on the range, not on the
point.
\end{exercise}

\begin{exercise}[\textit{Defending the omissions}]
Take the model from the project deliverable. List the five
most consequential omissions. For each, state the regime in
which the omission becomes consequential and the data that
would identify whether the regime obtains. Rank the omissions
by the cost of being wrong.

\paragraph{Discussion.} The exercise asks the reader to perform
the omitted-variable analysis on their own project deliverable. A
worked example for the cooling-cup model: the five most consequential
omissions are (1) the radiative term ($T - T_{\infty}$ above
$\sim 50\,\degC$ excess; data identifying the regime: surface
temperature measurement above $80\,\degC$); (2) the evaporative term
(open surface in dry air; data: relative humidity, surface area);
(3) the geometry term (non-cylindrical cup; data: cup dimensions
and material thermal conductivity); (4) the moving-air convective
term (any non-still ambient; data: anemometer reading near the
cup); (5) the thermometer thermal-mass lag (fast cooling, small
thermal mass; data: thermometer time constant relative to the
cup's cooling timescale).

Rank by cost of being wrong: the radiative term is the largest in
the regime of interest (kettle cooling, very hot drinks), so it
tops the list; evaporation is next (significant for open mugs in
dry rooms); the moving-air term follows (a draught in a kitchen
doubles $h$); geometry is fourth (matters across cups of different
shape but is moderate for a fixed cup); the thermometer lag is
fifth (the smallest in normal conditions). The ranking is the
information the engineer carries forward when deciding which
omission to address first if the prediction interval has to
shrink.
\end{exercise}

\begin{exercise}[\textit{The closing judgement}]
Take any model the reader has built across Volumes I and II.
State, in two sentences, the engineering question the model
answers, and in one sentence, the regime in which the answer
is no longer trustworthy. Compare the two with the equivalent
statements written by a peer about the same model. Where the
two diverge, identify the modelling decision that produced
the divergence.

\paragraph{Discussion.} The exercise has no fixed answer; its
purpose is to expose the modelling decisions that the writer made
without naming. Pairs of readers who work on the same model
(cooling cup, pendulum, RC circuit, beam) reliably produce
divergent question statements and divergent domain-of-validity
statements, even when the two readers fit the same data and
report the same numbers.

Common sources of divergence: one reader frames the question as
``how long until the cup is drinkable'' and the other as ``what
is the thermal time constant of the cup'' (different output
quantities, different precision requirements); one reader sets the
domain as ``below $80\,\degC$ excess'' (radiative cutoff) and the
other as ``no draught'' (convective cutoff); one reader writes the
prediction interval as $\pm \sigma_{p}$ and the other as
$\pm \sqrt{\sigma_{p}^{2} + \sigma_{s}^{2}}$ (parameter-only vs
honest). Naming the divergent decision is the exercise; the
reader closes Volume II with the habit that the divergent decision
is the modelling decision worth recording. The reflection
discipline carries forward into Volume III, where every project
deliverable carries an explicit assumptions ledger.
\end{exercise}
